\documentclass[11pt, a4paper]{article}
\usepackage[a4paper, margin=2.5cm]{geometry}
\usepackage[T1]{fontenc}
\usepackage[utf8]{inputenc} % Gwarantuje prawidłowe czytanie polskich znaków
\usepackage{lmodern} % Standardowa, ulepszona czcionka Latin Modern
\usepackage[polish]{babel} % Czysta konfiguracja dla pdfLaTeX

\usepackage{amsmath}
\usepackage{amssymb}
\usepackage{booktabs}
\usepackage{enumitem}
\setlist[itemize]{label=-}
\usepackage{hyperref}

\title{\textbf{Opracowanie zagadnień na egzamin licencjacki z fizyki}}
\author{Fizyka I stopnia}
\date{}

\begin{document}

\maketitle
\tableofcontents
\newpage

\section{Zasady dynamiki Newtona. Pęd. Zasada zachowania pędu.}

Klasyczna mechanika opiera się na trzech postulatach sformułowanych przez Isaaca Newtona, które definiują związek między ruchem ciał a działającymi na nie siłami w inercjalnych układach odniesienia.

\textbf{I Zasada Dynamiki (Zasada bezwładności):} Jeżeli na ciało nie działa żadna siła zewnętrzna lub działające siły się równoważą (wypadkowa sił jest równa zeru, $\sum \vec{F} = 0$), to ciało to pozostaje w spoczynku lub porusza się ruchem jednostajnym prostoliniowym. Zasada ta definiuje pojęcie układu inercjalnego – czyli takiego układu odniesienia, w którym to prawo jest spełnione. Własność ciał polegająca na dążeniu do zachowania swojego stanu ruchu nazywamy bezwładnością, a miarą bezwładności jest masa ciała ($m$).

\textbf{II Zasada Dynamiki:} Zmiana pędu ciała w czasie jest wprost proporcjonalna do wypadkowej siły działającej na to ciało i odbywa się w kierunku działania tej siły. W postaci uogólnionej równanie to przyjmuje formę:
\[
\vec{F}_{wyp} = \frac{d\vec{p}}{dt}
\]
W przypadku klasycznym, dla ciał o stałej masie, równanie to upraszcza się do znanej postaci $\vec{F}_{wyp} = m\vec{a}$, gdzie $\vec{a}$ to przyspieszenie ciała. Postać z pochodną pędu jest bardziej ogólna – sprawdza się również w układach o zmiennej masie (np. ruch rakiety) oraz w mechanice relatywistycznej.

\textbf{III Zasada Dynamiki (Zasada akcji i reakcji):} Oddziaływania ciał są zawsze wzajemne. Jeżeli ciało A działa na ciało B siłą $\vec{F}_{AB}$, to ciało B działa na ciało A siłą $\vec{F}_{BA}$ o takiej samej wartości, tym samym kierunku, lecz przeciwnym zwrocie ($\vec{F}_{AB} = -\vec{F}_{BA}$). Siły te przyłożone są do dwóch różnych ciał, dlatego się nie równoważą.

\textbf{Pęd i jego zasada zachowania:}
Pęd (oznaczany jako $\vec{p}$) jest wielkością wektorową definiowaną jako iloczyn masy ciała i jego prędkości:
\[
\vec{p} = m\vec{v}
\]
Z II zasady dynamiki wynika fundamentalne prawo przyrody: \textbf{Zasada zachowania pędu}. Mówi ona, że jeżeli układ ciał jest izolowany (czyli wypadkowa sił zewnętrznych działających na układ wynosi zero, $\vec{F}_{zewn} = 0$), to całkowity pęd tego układu pozostaje stały w czasie ($\frac{d\vec{p}_{ukl}}{dt} = 0$, co implikuje $\vec{p}_{ukl} = const$). Zasada ta ma kluczowe znaczenie przy analizie zderzeń (sprężystych i niesprężystych) oraz rozpadów cząstek. Działa w skali makroskopowej i mikroskopowej, stanowiąc konsekwencję jednorodności przestrzeni.

\section{Inercjalne i nieinercjalne układy odniesienia. Siły bezwładności.}

W mechanice opis ruchu zależy od wybranego układu odniesienia. Zasadniczo dzielimy je na dwie kategorie: inercjalne i nieinercjalne, co ma fundamentalne konsekwencje dla postaci równań ruchu.

\textbf{Inercjalne układy odniesienia} to takie, w których spełniona jest I zasada dynamiki Newtona. Jeśli ciało nie podlega żadnym oddziaływaniom, spoczywa lub porusza się ruchem jednostajnym prostoliniowym względem tego układu. Ważną własnością jest to, że każdy układ odniesienia poruszający się ruchem jednostajnym prostoliniowym względem układu inercjalnego jest również układem inercjalnym. W mechanice klasycznej przejście między takimi układami opisują \textbf{transformacje Galileusza}, które zakładają absolutność czasu i przestrzeni. We wszystkich układach inercjalnych prawa mechaniki klasycznej mają identyczną postać.

\textbf{Nieinercjalne układy odniesienia} to układy, które poruszają się z pewnym przyspieszeniem względem układów inercjalnych (np. ruszający pociąg, karuzela, obracająca się Ziemia). W takich układach I zasada dynamiki wydaje się być łamana – ciała mogą doznawać przyspieszenia bez działania żadnych fizycznych (pochodzących od oddziaływań) sił.

Aby zachować strukturę II zasady dynamiki Newtona w układach nieinercjalnych ($\vec{F} = m\vec{a}_{wzg}$, gdzie $\vec{a}_{wzg}$ to przyspieszenie w tym układzie), wprowadza się pojęcie \textbf{sił bezwładności} (pozornych). Nie są one wynikiem oddziaływań ciał, lecz samego faktu przyspieszania układu odniesienia.

W ogólnym przypadku układu obracającego się z prędkością kątową $\vec{\omega}$ i poruszającego się ruchem postępowym z przyspieszeniem $\vec{a}_0$, na ciało o masie $m$ poruszające się z prędkością $\vec{v}_{wzg}$ działają następujące siły bezwładności:
\begin{itemize}
    \item \textbf{Siła bezwładności ruchu postępowego (unoszenia):} $\vec{F}_u = -m\vec{a}_0$. Działa przeciwnie do wektora przyspieszenia układu (np. wgniatanie w fotel w przyspieszającym samochodzie).
    \item \textbf{Siła odśrodkowa:} $\vec{F}_{od} = m\omega^2\vec{r}$ (gdzie $\vec{r}$ to wektor promienia wodzącego prostopadły do osi obrotu). Występuje w układach obracających się, kierowana zawsze na zewnątrz od osi obrotu. Jest odpowiedzialna np. za spłaszczenie Ziemi na biegunach.
    \item \textbf{Siła Coriolisa:} $\vec{F}_C = -2m(\vec{\omega} \times \vec{v}_{wzg})$. Zależy od prędkości ciała względem obracającego się układu odniesienia. Działa prostopadle do wektora prędkości i osi obrotu. Na Ziemi odpowiada m.in. za kierunek wirowania cyklonów oraz podmywanie prawych brzegów rzek na półkuli północnej.
    \item \textbf{Siła bezwładności Eulera:} $\vec{F}_E = -m(\frac{d\vec{\omega}}{dt} \times \vec{r})$. Występuje tylko wtedy, gdy prędkość kątowa układu zmienia się w czasie (ruch obrotowy niejednostajny).
\end{itemize}
Siły bezwładności są całkowicie realne z punktu widzenia obserwatora znajdującego się w układzie nieinercjalnym i mogą wywoływać mierzalne skutki fizyczne.

\section{Praca i energia. Energia kinetyczna i potencjalna. Siły zachowawcze. Zasada zachowania energii mechanicznej.}

Pojęcia pracy i energii stanowią alternatywny (wobec analizy wektorowej opartej na siłach) skalarowy sposób opisu układów fizycznych. Zastosowanie metod energetycznych bardzo często upraszcza rozwiązywanie złożonych problemów mechanicznych.

\textbf{Praca (W)} w sensie fizycznym jest wykonywana, gdy na ciało działa siła, a ciało to ulega przemieszczeniu, przy czym wektor siły posiada niezerową składową zgodną z kierunkiem ruchu. Dla siły stałej wektorowo praca jest iloczynem skalarnym: $W = \vec{F} \cdot \Delta\vec{r} = F \Delta r \cos\alpha$. W przypadku ogólnym, gdy siła zmienia się wzdłuż toru krzywoliniowego $L$, praca jest definiowana jako całka krzywoliniowa:
\[
W = \int_L \vec{F}(\vec{r}) \cdot d\vec{r}
\]

\textbf{Twierdzenie o pracy i energii kinetycznej:} 
Energia kinetyczna ($E_k$) to energia związana z ruchem ciała, określona dla punktu materialnego wzorem $E_k = \frac{1}{2}mv^2$. Całkowita praca wykonana przez wszystkie siły (wypadkową siłę) działające na ciało jest dokładnie równa zmianie jego energii kinetycznej:
\[
W = \Delta E_k = E_{k, konc} - E_{k, pocz}
\]

\textbf{Siły zachowawcze i energia potencjalna:}
Siłę nazywamy \textbf{zachowawczą}, jeżeli praca wykonywana przez nią nad ciałem przemieszczającym się między dwoma punktami zależy wyłącznie od położenia tych punktów, a nie od kształtu wybranej drogi. Równoważnie: praca siły zachowawczej na drodze zamkniętej wynosi zero ($\oint \vec{F} \cdot d\vec{r} = 0$). Z matematycznego punktu widzenia pole sił zachowawczych jest bezwirowe ($\nabla \times \vec{F} = 0$). Przykłady to siła grawitacji, siła sprężystości czy siła elektrostatyczna (siłą niezachowawczą jest np. siła tarcia).

Dla sił zachowawczych można zdefiniować \textbf{energię potencjalną ($E_p$)}, która jest funkcją położenia. Zmiana energii potencjalnej przy przemieszczeniu ciała równa jest pracy wykonanej przez siłę zachowawczą wziętej ze znakiem minus:
\[
W_{zach} = -\Delta E_p \quad \text{oraz} \quad \vec{F} = -\nabla E_p
\]
Energia potencjalna określa możliwość wykonania pracy przez układ dzięki wzajemnemu położeniu jego elementów (np. $E_p = mgh$ w jednorodnym polu grawitacyjnym, $E_p = \frac{1}{2}kx^2$ dla sprężyny).

\textbf{Zasada zachowania energii mechanicznej:}
Całkowita energia mechaniczna układu to suma jego energii kinetycznej i potencjalnej ($E = E_k + E_p$). Z twierdzenia o pracy i energii kinetycznej wynika, że jeżeli w izolowanym układzie ciał działają wyłącznie siły zachowawcze, to całkowita energia mechaniczna układu jest wielkością stałą:
\[
\Delta E = \Delta E_k + \Delta E_p = 0 \implies E_k + E_p = const
\]
Gdy w układzie działają również siły niezachowawcze (np. tarcie, opór powietrza), całkowita energia mechaniczna nie jest zachowana, a jej ubytek jest równy pracy wykonanej przez te siły rozpraszające (często zamieniając się w energię termiczną, zgodnie z I zasadą termodynamiki).

\section{Dynamika bryły sztywnej. Zasada zachowania momentu pędu.}

\textbf{Bryła sztywna} to wyidealizowany model ciała rozciągłego, w którym odległości między dowolnymi dwoma jego punktami pozostają stałe, niezależnie od działających sił (ciało, które nie ulega odkształceniom). Dowolny ruch bryły sztywnej można rozłożyć na ruch postępowy jej środka masy oraz ruch obrotowy wokół osi przechodzącej przez środek masy.

Kluczową wielkością opisującą rozkład masy bryły względem danej osi obrotu jest \textbf{moment bezwładności (I)}. Odgrywa on w ruchu obrotowym rolę analogiczną do masy w ruchu postępowym – jest miarą bezwładności ciała na zmianę jego prędkości kątowej. Dla układu punktów materialnych $I = \sum m_i r_i^2$, a w ujęciu ciągłym:
\[
I = \int r^2 dm
\]
gdzie $r$ to odległość elementu masy $dm$ od osi obrotu. Obliczanie momentów bezwładności ułatwia \textbf{Twierdzenie Steinera}: $I = I_0 + md^2$, gdzie $I_0$ to moment bezwładności względem osi przechodzącej przez środek masy, a $d$ to odległość nowej osi od osi środka masy.

\textbf{Równanie dynamiki ruchu obrotowego (II zasada dynamiki dla ruchu obrotowego):}
Zmiana wektora momentu pędu bryły w czasie jest równa sumie momentów sił zewnętrznych działających na tę bryłę:
\[
\vec{M}_{wyp} = \frac{d\vec{L}}{dt}
\]
Gdzie $\vec{M}$ to wektor momentu siły ($\vec{M} = \vec{r} \times \vec{F}$). Dla bryły sztywnej obracającej się wokół stałej osi symetrii równanie to upraszcza się do postaci $M_{wyp} = I\varepsilon$, gdzie $\varepsilon$ to przyspieszenie kątowe (odpowiednik $F = ma$). Energia kinetyczna ruchu obrotowego wynosi $E_{k, obr} = \frac{1}{2}I\omega^2$.

\textbf{Moment pędu i jego zasada zachowania:}
Moment pędu punktu materialnego względem bieguna to wektor $\vec{L} = \vec{r} \times \vec{p}$. Dla bryły sztywnej obracającej się wokół ustalonej osi moment pędu powiązany jest z prędkością kątową relacją $\vec{L} = I\vec{\omega}$.
Z równania dynamiki bryły sztywnej wynika \textbf{Zasada zachowania momentu pędu}: 
Jeżeli wypadkowy moment sił zewnętrznych działających na układ wynosi zero ($\vec{M}_{zewn} = 0$), to całkowity moment pędu układu pozostaje stały:
\[
\frac{d\vec{L}}{dt} = 0 \implies \vec{L} = const \quad \text{czyli} \quad I\omega = const
\]
Zasada ta tłumaczy powszechne zjawiska, takie jak zwiększenie prędkości obrotowej łyżwiarza wykonującego piruet po przyciagnięciu ramion do tułowia (zmniejszenie $I$ powoduje wzrost $\omega$ w celu zachowania iloczynu $I\omega$). Jest to fundamentalne prawo fizyki mające zastosowanie od układów atomowych po galaktyki.

\section{Oscylator harmoniczny i oscylator z tłumieniem. Drgania wymuszone i rezonans.}

Drgania stanowią jeden z najbardziej uniwersalnych rodzajów ruchu w fizyce. Dowolny układ fizyczny wyprowadzony ze stanu trwałej równowagi i pozostawiony samemu sobie będzie wykonywał drgania wokół tego położenia.

\textbf{Oscylator harmoniczny prosty (drgania niewymuszone, nietłumione):}
Występuje wtedy, gdy jedyną siłą działającą na ciało jest siła przywracająca (kierująca do położenia równowagi), proporcjonalna do wychylenia $x$: $F = -kx$. Zgodnie z II zasadą dynamiki Newtona otrzymujemy równanie różniczkowe:
\[
m\ddot{x} + kx = 0 \implies \ddot{x} + \omega_0^2 x = 0
\]
gdzie $\omega_0 = \sqrt{k/m}$ to częstość kołowa drgań własnych. Rozwiązaniem jest funkcja sinusoidalna: $x(t) = A \cos(\omega_0 t + \phi_0)$, gdzie $A$ to amplituda, a $\phi_0$ to faza początkowa. Energia całkowita oscylatora jest stała ($E = \frac{1}{2}kA^2$) i cyklicznie przechodzi z potencjalnej w kinetyczną i odwrotnie.

\textbf{Oscylator z tłumieniem:}
W rzeczywistości zawsze obecne są siły oporu ośrodka. W modelu tłumienia lepkiego siła oporu jest proporcjonalna do prędkości: $F_{op} = -b\dot{x}$. Równanie ruchu przyjmuje postać:
\[
m\ddot{x} + b\dot{x} + kx = 0 \implies \ddot{x} + 2\beta\dot{x} + \omega_0^2 x = 0
\]
gdzie $\beta = b/(2m)$ to współczynnik tłumienia. W zależności od wartości $\beta$, rozwiązania dzielą się na trzy przypadki:
\begin{itemize}
    \item Tłumienie słabe ($\beta < \omega_0$): ruch jest wciąż oscylacyjny, ale amplituda maleje eksponencjalnie w czasie $x(t) = A_0 e^{-\beta t}\cos(\omega t + \phi)$, gdzie $\omega = \sqrt{\omega_0^2 - \beta^2}$.
    \item Tłumienie krytyczne ($\beta = \omega_0$): układ powraca do położenia równowagi w najkrótszym możliwym czasie bez oscylacji (zastosowanie w amortyzatorach).
    \item Tłumienie silne (aperiodyczne) ($\beta > \omega_0$): układ powraca powoli do zera ruchem pełzającym, bez przechodzenia przez położenie równowagi.
\end{itemize}

\textbf{Drgania wymuszone i rezonans:}
Jeżeli do tłumionego oscylatora przyłożymy zewnętrzną siłę harmoniczną $F(t) = F_0 \cos(\omega t)$, otrzymamy drgania wymuszone. Po pewnym czasie zanikają stany przejściowe i układ drga z częstością wymuszającą $\omega$, a nie swoją częstością własną. Amplituda tych drgań zależy od $\omega$:
\[
A(\omega) = \frac{F_0/m}{\sqrt{(\omega_0^2 - \omega^2)^2 + 4\beta^2\omega^2}}
\]
\textbf{Zjawisko rezonansu} zachodzi, gdy częstość siły wymuszającej $\omega$ zbliża się do częstości własnej układu $\omega_0$ (dokładniej do wartości $\sqrt{\omega_0^2 - 2\beta^2}$). Wtedy mianownik we wzorze na amplitudę osiąga minimum, a sama amplituda $A$ osiąga gwałtowne maksimum. Im mniejsze tłumienie $\beta$, tym wyższy i węższy jest pik rezonansowy. Zjawisko to pozwala na efektywny transfer energii do układu, ma kluczowe znaczenie w akustyce, elektronice i konstrukcji maszyn, gdzie niekontrolowany rezonans (np. mostu) może prowadzić do katastrofy.

\section{Grawitacja: prawo grawitacji Newtona, prawa Keplera, grawitacyjna energia potencjalna.}

\textbf{Prawo powszechnego ciążenia Newtona:}
Isaac Newton sformułował fundamentalne prawo natury głoszące, że każde dwa ciała obdarzone masą przyciągają się wzajemnie siłą grawitacyjną. Wartość tej siły jest wprost proporcjonalna do iloczynu ich mas ($m_1$, $m_2$) i odwrotnie proporcjonalna do kwadratu odległości między ich środkami ($r$):
\[
\vec{F} = -G \frac{m_1 m_2}{r^2} \hat{r}
\]
gdzie $G \approx 6.674 \times 10^{-11} \, \text{N}\cdot\text{m}^2/\text{kg}^2$ to stała grawitacji, a $\hat{r}$ to wektor jednostkowy kierunku łączącego masy. Znak minus oznacza, że siła ma charakter przyciągający. Prawo to, opracowane początkowo dla mas punktowych, ze względu na kulistą symetrię i twierdzenie Gaussa, jest dokładnie spełnione dla makroskopowych ciał o sferycznym rozkładzie gęstości (jak gwiazdy i planety).

\textbf{Grawitacyjna energia potencjalna:}
Ponieważ pole grawitacyjne jest polem sił zachowawczych, można dla niego zdefiniować energię potencjalną. Jako punkt odniesienia (zero energii) przyjmuje się nieskończoność ($r \to \infty$). Wówczas energia potencjalna układu dwóch mas odległych o $r$ wyraża się wzorem:
\[
E_p(r) = -G \frac{m_1 m_2}{r}
\]
Znak ujemny oznacza, że układ dwóch mas jest związany (aby je oddalić do nieskończoności należy wykonać dodatnią pracę). Z tego pojęcia wynika kluczowe pojęcie \textbf{pierwszej i drugiej prędkości kosmicznej}. Druga prędkość kosmiczna (prędkość ucieczki) jest wyznaczana z zasady zachowania energii ($E_k + E_p = 0$), dając dla Ziemi wynik $v_{II} = \sqrt{2GM/R} \approx 11.2$ km/s.

\textbf{Prawa Keplera:}
Prawa empiryczne sformułowane na podstawie obserwacji astronomicznych przez Johannesa Keplera, które można analitycznie wyprowadzić z zasad dynamiki i prawa grawitacji Newtona:
\begin{itemize}
    \item \textbf{I Prawo Keplera:} Każda planeta Układu Słonecznego porusza się po orbicie w kształcie elipsy, w której jednym z ognisk znajduje się Słońce. Jest to bezpośredni dowód na to, że siła zależy od $1/r^2$.
    \item \textbf{II Prawo Keplera (Prawo pól):} Promień wodzący planety (odcinek łączący Słońce z planetą) w równych odstępach czasu zakreśla równe pola powierzchni. Oznacza to, że prędkość liniowa planety jest największa w peryhelium (blisko Słońca), a najmniejsza w aphelium. Wynika to wprost z zasady zachowania momentu pędu planety.
    \item \textbf{III Prawo Keplera:} Kwadrat okresu obiegu planety wokół Słońca ($T$) jest wprost proporcjonalny do sześcianu wielkiej półosi jej orbity ($a$):
    \[
    \frac{T^2}{a^3} = \frac{4\pi^2}{GM} = const
    \]
    gdzie $M$ to masa Słońca. Prawo to umożliwia określanie mas obiektów centralnych (np. Słońca, Jowisza) na podstawie obserwacji ich satelitów.
\end{itemize}

\section{Statyka i dynamika płynów: prawo Pascala, prawo Archimedesa, równanie ciągłości strugi, równanie Bernoulliego.}

Płyny (ciecze i gazy) to substancje, które nie stawiają trwałego oporu przy zmianie kształtu (nie przenoszą naprężeń ścinających w spoczynku). Z tego względu do ich opisu stosuje się pojęcia gęstości ($\rho = m/V$) oraz ciśnienia ($p = F/S$), a nie siły punktowe.

\textbf{Prawo Pascala (Statyka płynów):}
Zewnętrzne ciśnienie wywierane na zamknięty płyn jest przekazywane bez zmian na każdą część płynu oraz na ścianki naczynia. Oznacza to, że zmiana ciśnienia $\Delta p$ w jednym punkcie zamkniętego układu hydraulicznego powoduje identyczną zmianę ciśnienia w każdym innym jego punkcie. Prawo to jest podstawą działania maszyn hydraulicznych (np. prasy, hamulców), gdzie mała siła działająca na mały tłok generuje odpowiednio większą siłę na dużym tłoku: $F_1/S_1 = F_2/S_2$.

\textbf{Prawo Archimedesa (Statyka płynów):}
Na każde ciało zanurzone w płynie (całkowicie lub częściowo) działa skierowana ku górze siła wyporu. Wartość tej siły jest równa ciężarowi płynu wypartego przez to ciało:
\[
F_{wyp} = \rho_{plynu} \cdot g \cdot V_{zanurzone}
\]
gdzie $\rho_{plynu}$ to gęstość płynu, a $V_{zanurzone}$ to objętość wypartego płynu. Zjawisko to wynika z różnicy ciśnień hydrostatycznych działających na dolną i górną powierzchnię zanurzonego ciała ($p = \rho g h$). Prawo to decyduje o pływalności ciał (ciało tonie, gdy $\rho_{ciala} > \rho_{plynu}$, a pływa, gdy $\rho_{ciala} < \rho_{plynu}$).

\textbf{Równanie ciągłości strugi (Dynamika płynów):}
W dynamice płynów idealnych (nieściśliwych i pozbawionych lepkości) ważnym prawem jest zasada zachowania masy. Dla ustalonego przepływu przez rurę o zmiennym przekroju $S$, masa płynu przepływająca w jednostce czasu przez każdy przekrój musi być stała. Ponieważ płyn jest nieściśliwy ($\rho = const$), stałe jest natężenie przepływu objętościowego:
\[
S_1 \cdot v_1 = S_2 \cdot v_2 = const
\]
gdzie $v$ to prędkość przepływu. Oznacza to, że w przewężeniach płyn musi przyspieszać.

\textbf{Równanie Bernoulliego:}
Stanowi ono matematyczny zapis zasady zachowania energii dla płynącego idealnego płynu. Wzdłuż dowolnej linii prądu suma ciśnienia statycznego ($p$), ciśnienia dynamicznego (związanego z energią kinetyczną objętości właściwej, $\frac{1}{2}\rho v^2$) oraz ciśnienia hydrostatycznego (związanego z energią potencjalną, $\rho g h$) jest stała:
\[
p + \frac{1}{2}\rho v^2 + \rho g h = const
\]
Z równania tego wynika szereg ważnych zjawisk: np. paradoks hydrodynamiczny (spadek ciśnienia w miejscu przewężenia rury, gdzie rośnie prędkość) – wykorzystywany m.in. w palnikach Bunsena, gaźnikach, atomizerach oraz (w pewnym przybliżeniu) tłumaczący powstawanie siły nośnej na skrzydle samolotu.

\section{Transformacje Galileusza. Transformacja Lorentza. Relatywistyczne transformacje długości, prędkości i czasu. Relatywistyczny pęd i energia.}

\textbf{Transformacje Galileusza i ich ograniczenia:}
W mechanice klasycznej przejście od jednego układu inercjalnego ($O$) do drugiego ($O'$), poruszającego się względem pierwszego ze stałą prędkością $v$ wzdłuż osi $x$, opisują transformacje Galileusza: $x' = x - vt$, $y' = y$, $z' = z$, $t' = t$. Zakładają one absolutność czasu ($t'=t$) i przestrzeni (długość obiektu jest taka sama w każdym układzie). Prowadzą one również do klasycznego prawa dodawania prędkości ($u' = u - v$). Zakończyły się one jednak niepowodzeniem przy próbie opisu zjawisk elektromagnetycznych (równania Maxwella nie są względem nich niezmiennicze), a doświadczenie Michelsona-Morleya wykazało, że prędkość światła $c$ jest stała niezależnie od ruchu obserwatora.

\textbf{Transformacja Lorentza:}
Albert Einstein oparł Szczególną Teorię Względności (STW) na dwóch postulatach: wszystkie prawa fizyki (w tym elektromagnetyzmu) są jednakowe w każdym inercjalnym układzie odniesienia, a prędkość światła w próżni $c$ jest absolutna i niezależna od układu. Poprawnym matematycznie przejściem między układami inercjalnymi stały się transformacje Lorentza (przyjmując ruch układu $O'$ z prędkością $v$ wzdłuż osi $x$):
\[
x' = \gamma(x - vt), \quad y' = y, \quad z' = z, \quad t' = \gamma\left(t - \frac{vx}{c^2}\right)
\]
gdzie $\gamma = \frac{1}{\sqrt{1 - v^2/c^2}}$ to czynnik Lorentza. Z transformacji tych wynikają fascynujące, całkowicie nieintuicyjne dla małych prędkości (gdzie $\gamma \approx 1$ i wzory redukują się do Galileusza) zjawiska.

\textbf{Relatywistyczne efekty czasoprzestrzenne:}
\begin{itemize}
    \item \textbf{Dylatacja czasu:} Zegar poruszający się z prędkością $v$ względem obserwatora odmierza czas wolniej niż zegar spoczywający względem niego: $\Delta t = \gamma \Delta t_0$. $\Delta t_0$ to czas własny.
    \item \textbf{Kontrakcja (skrócenie) długości:} Przedmiot poruszający się z prędkością $v$ jest skrócony w kierunku ruchu w oczach obserwatora stacjonarnego: $L = L_0 / \gamma$, gdzie $L_0$ to długość spoczynkowa.
    \item \textbf{Relatywistyczne składanie prędkości:} Nie możemy przekroczyć $c$. Jeśli obiekt porusza się z prędkością $u'$ w układzie $O'$, to w układzie stacjonarnym $O$ jego prędkość wynosi: $u = \frac{u' + v}{1 + \frac{u' v}{c^2}}$.
\end{itemize}

\textbf{Relatywistyczny pęd i energia:}
Aby zasada zachowania pędu obowiązywała we wszystkich układach inercjalnych zgodnie z transformacją Lorentza, definicja pędu musiała ulec modyfikacji. Relatywistyczny pęd ciała o masie $m$ (zawsze masie spoczynkowej) i prędkości $v$ wynosi:
\[
\vec{p} = \gamma m \vec{v} = \frac{m \vec{v}}{\sqrt{1 - v^2/c^2}}
\]
Praca wykonana nad ciałem przez siłę powoduje wzrost jego energii całkowitej, która w ujęciu relatywistycznym definiowana jest wzorem $E = \gamma m c^2$. Energia całkowita składa się z energii spoczynkowej ($E_0 = mc^2$) oraz energii kinetycznej ($E_k$):
\[
E = E_k + E_0 \implies E_k = E - E_0 = (\gamma - 1)mc^2 = mc^2 \left(\frac{1}{\sqrt{1 - v^2/c^2}} - 1\right)
\]
Dla małych prędkości ($v \ll c$) rozwinięcie tego wzoru w szereg Taylora daje znajome klasyczne $E_k \approx \frac{1}{2}mv^2$. Ważnym niezmiennikiem relatywistycznym łączącym energię i pęd jest relacja: $E^2 = (pc)^2 + (mc^2)^2$.

\section{Gaz doskonały, jego przemiany i równanie stanu.}

\textbf{Model gazu doskonałego:}
Gaz doskonały to termodynamiczny model teoretyczny upraszczający zachowanie rzeczywistych gazów. Opiera się na następujących mikroskopowych postulatach kinetyczno-molekularnych:
\begin{itemize}
    \item Gaz składa się z bardzo dużej liczby cząsteczek, które są traktowane jako punkty materialne (ich objętość własna jest zaniedbywalnie mała w porównaniu z objętością naczynia).
    \item Cząsteczki poruszają się w sposób chaotyczny (ruchy Browna, rozkład Maxwella), zgodnie z zasadami dynamiki Newtona.
    \item Pomiędzy cząsteczkami nie występują żadne siły oddziaływania na odległość (brak energii potencjalnej oddziaływań międzycząsteczkowych).
    \item Zderzenia cząsteczek między sobą oraz ze ściankami naczynia są idealnie sprężyste (zachowana jest zarówno energia kinetyczna, jak i pęd układu).
\end{itemize}
Większość rzeczywistych gazów (np. azot, tlen, hel) zachowuje się w sposób bardzo zbliżony do gazu doskonałego w warunkach stosunkowo niskiego ciśnienia i wysokiej temperatury (z dala od punktu skraplania).

\textbf{Równanie stanu gazu doskonałego (Równanie Clapeyrona):}
Stan danej masy gazu doskonałego jest w pełni określony przez trzy makroskopowe parametry termodynamiczne: ciśnienie ($p$), objętość ($V$) i temperaturę bezwzględną ($T$). Równanie wiążące te wielkości ma postać:
\[
pV = nRT \quad \text{lub alternatywnie} \quad pV = NkT
\]
gdzie $n$ to liczba moli gazu, $R \approx 8.314 \, \text{J/(mol}\cdot\text{K)}$ to uniwersalna stała gazowa, $N$ to całkowita liczba cząsteczek, a $k \approx 1.38 \times 10^{-23} \, \text{J/K}$ to stała Boltzmanna ($R = N_A k$, gdzie $N_A$ to liczba Avogadra).

\textbf{Przemiany termodynamiczne gazu doskonałego:}
Z równania Clapeyrona można wyprowadzić specyficzne prawa dla izo-przemian (procesów, w których jeden z parametrów jest stały):
\begin{itemize}
    \item \textbf{Przemiana izotermiczna} ($T = const$): Prawo Boyle'a-Mariotte'a. Jeżeli temperatura jest stała, ciśnienie gazu jest odwrotnie proporcjonalne do jego objętości: $pV = const$. Na wykresie $p(V)$ wykresem jest hiperbola (izoterma). Praca wykonywana przez gaz równa jest dostarczonemu ciepłu, a energia wewnętrzna nie ulega zmianie.
    \item \textbf{Przemiana izobaryczna} ($p = const$): Prawo Gay-Lussaca. Objętość gazu jest wprost proporcjonalna do jego temperatury bezwzględnej: $V/T = const$. Wykres na płaszczyźnie $V(T)$ to półprosta wychodząca z początku układu współrzędnych. Praca wyraża się prostym wzorem $W = p\Delta V$.
    \item \textbf{Przemiana izochoryczna} ($V = const$): Prawo Charles'a. Ciśnienie gazu jest wprost proporcjonalne do temperatury: $p/T = const$. Ponieważ objętość się nie zmienia, gaz nie wykonuje pracy ($W = 0$), a całe dostarczone ciepło zamienia się na wzrost energii wewnętrznej ($Q = \Delta U$).
    \item \textbf{Przemiana adiabatyczna} ($Q = 0$): Przemiana zachodząca bez wymiany ciepła z otoczeniem (np. w izolowanym naczyniu lub bardzo szybkim procesie). Rządzi się równaniem Poissona: $pV^\kappa = const$, gdzie $\kappa = C_p/C_v$ to stosunek ciepła właściwego przy stałym ciśnieniu do ciepła właściwego przy stałej objętości. Na wykresie $p(V)$ adiabata opada bardziej stromo niż izoterma. Zmiana energii wewnętrznej odbywa się wyłącznie kosztem pracy zewnętrznej ($\Delta U = -W$).
\end{itemize}

\section{Pierwsza zasada termodynamiki. Zasada ekwipartycji energii.}

\textbf{Pierwsza zasada termodynamiki:}
Jest to adaptacja ogólnej zasady zachowania energii do układów termodynamicznych. W układach tych zjawiska mechaniczne i termiczne przenikają się nawzajem. Stan wewnętrzny każdego układu charakteryzuje jednoznacznie funkcja stanu zwana \textbf{energią wewnętrzną ($U$)}. Dla gazu doskonałego jest to suma energii kinetycznych wszystkich jego cząsteczek.
I zasada termodynamiki mówi, że energię wewnętrzną układu można zmienić na dwa sposoby: poprzez wymianę ciepła z otoczeniem ($Q$) lub poprzez wykonanie pracy ($W$). Matematycznie zapisujemy to (w konwencji, gdzie $W$ to praca wykonana \textit{nad} układem):
\[
\Delta U = Q + W
\]
W zapisie różniczkowym dla procesów elementarnych (quasistatycznych): $dU = \delta Q + \delta W$. Warto zauważyć notację: $dU$ jest różniczką zupełną (zmiana zależy tylko od stanu początkowego i końcowego), podczas gdy $\delta Q$ i $\delta W$ są małymi porcjami przekazanej energii (ich całkowita wartość zależy od drogi, czyli rodzaju procesu). Praca objętościowa wyraża się całką $W = -\int p dV$ (jeśli gaz się rozpręża, $dV > 0$, układ wykonuje pracę nad otoczeniem, więc praca wykonana nad układem jest ujemna, zasilając rozprężanie kosztem $U$).

W zastosowaniu do przemian gazowych:
- w procesie izochorycznym $W=0$, więc $\Delta U = Q$.
- w procesie adiabatycznym $Q=0$, więc $\Delta U = W$ (gaz rozprężając się ruchem adiabatycznym ulega ochłodzeniu, bo traci energię wewnętrzną).
- w cyklu zamkniętym (powrót do stanu początkowego) $\Delta U = 0$, co oznacza sumarycznie $Q = -W$ (ciepło zamienia się w pracę, ale nie jest możliwe stworzenie \textit{perpetuum mobile} pierwszego rodzaju – maszyny wykonującej pracę bez dostarczania energii).

\textbf{Zasada ekwipartycji energii:}
Twierdzenie fizyki statystycznej sformułowane w mechanice klasycznej przez Maxwella i Boltzmanna. Mówi ono, że w stanie równowagi termodynamicznej, całkowita energia kinetyczna układu dzieli się równo po średnio $\frac{1}{2}kT$ (gdzie $k$ to stała Boltzmanna, a $T$ temperatura bezwzględna) na każdy niezależny \textbf{stopień swobody} cząsteczki.

Liczba stopni swobody ($f$) zależy od budowy cząsteczki:
\begin{itemize}
    \item Gaz jednoatomowy (np. Hel, Argon): Cząsteczki traktujemy jak punkty. Mają one $f=3$ stopnie swobody związane z ruchem postępowym w trzech wymiarach ($x, y, z$).
    \item Gaz dwuatomowy (np. $O_2$, $N_2$): Cząsteczki przypominają hantle. Oprócz trzech stopni postępowych, mają $f_{obr}=2$ istotne stopnie swobody związane z ruchem obrotowym (obrót wokół własnej osi wiązania ma zaniedbywalny moment bezwładności). Zatem $f = 3 + 2 = 5$. (W bardzo wysokich temperaturach mogą aktywować się drgania wzdłuż osi, dając kolejne stopnie).
    \item Gazy wieloatomowe (np. $CO_2$, $H_2O$): $f=6$ (3 postępowe, 3 obrotowe).
\end{itemize}

Z zasady tej wynika bezpośrednio wzór na całkowitą energię wewnętrzną jednego mola gazu doskonałego: $U = \frac{f}{2}RT$. Pozwala to na przewidywanie teoretycznych wartości ciepła właściwego gazów przy stałej objętości $C_v = \frac{f}{2}R$ oraz przy stałym ciśnieniu $C_p = C_v + R = \frac{f+2}{2}R$ (z relacji Mayera). Wyniki te bardzo dobrze zgadzają się z doświadczeniem dla umiarkowanych temperatur, jednak pełne zrozumienie zamarzania stopni swobody w niskich temperaturach wymagało już mechaniki kwantowej.

\section{Druga zasada termodynamiki. Ciepło zredukowane, entropia.}

Druga zasada termodynamiki określa kierunek zachodzenia zjawisk termodynamicznych i nakłada ograniczenia na wzajemne przemiany ciepła i pracy, których pierwsza zasada (będąca tylko bilansem energii) nie precyzuje. Wyjaśnia zjawisko nieodwracalności procesów naturalnych.

\textbf{Sformułowania Drugiej Zasady Termodynamiki:}
\begin{itemize}
    \item \textbf{Sformułowanie Kelvina-Plancka:} Niemożliwe jest zbudowanie cyklicznie działającej maszyny cieplnej, której jedynym rezultatem byłoby całkowite zamienienie ciepła pobranego ze zbiornika na pracę mechaniczną. Innymi słowy, każda maszyna cieplna musi oddawać część pobranego ciepła do chłodnicy (nie istnieje \textit{perpetuum mobile} drugiego rodzaju o sprawności $\eta = 100\%$). Najwyższą teoretyczną sprawność posiada idealny cykl Carnota: $\eta = 1 - T_{chl}/T_{grz}$.
    \item \textbf{Sformułowanie Clausiusa:} Niemożliwe jest przejście ciepła od ciała o temperaturze niższej do ciała o temperaturze wyższej, bez jednoczesnego zachodzenia innych zmian w układzie (bez wykonania pracy z zewnątrz). Proces samorzutny zawsze prowadzi do wyrównywania temperatur.
\end{itemize}
Zasadniczo oba sformułowania są równoważne – wykazano, że złamanie jednego pociąga za sobą natychmiastowe złamanie drugiego.

\textbf{Ciepło zredukowane i Entropia w ujęciu makroskopowym:}
Rozpatrując odwracalne cykle termodynamiczne, wprowadzono pojęcie \textbf{ciepła zredukowanego}, definiowanego jako stosunek porcji wymienionego ciepła $\delta Q$ do temperatury bezwzględnej $T$, w której ta wymiana nastąpiła ($\delta Q/T$). Twierdzenie Clausiusa mówi, że dla dowolnego procesu odwracalnego, suma całkowa ciepeł zredukowanych wokół pełnego cyklu jest równa zero: $\oint \frac{\delta Q_{odwr}}{T} = 0$. 

To matematycznie dowodzi, że całka od stanu A do B nie zależy od drogi przemiany odwracalnej, co pozwoliło Clausiusowi na zdefiniowanie nowej wielkości fizycznej będącej funkcją stanu – \textbf{entropii ($S$)}. Różniczka entropii wynosi:
\[
dS = \frac{\delta Q_{odwr}}{T} \implies \Delta S = \int_A^B \frac{\delta Q_{odwr}}{T}
\]
W procesach rzeczywistych (zawsze nieodwracalnych – np. rozprężanie gazu do próżni, swobodny przepływ ciepła, procesy z tarciem) układ izolowany zawsze zwiększa swoją entropię. Ostatecznie Druga Zasada Termodynamiki może być sformułowana najzwięźlej: \textbf{Entropia każdego układu izolowanego nie może maleć w czasie ($\Delta S \ge 0$)}. Procesy zachodzą samorzutnie tylko w kierunku rosnącej entropii układu wszechświata. Stan równowagi termodynamicznej to stan o maksymalnej możliwej entropii.

\textbf{Statystyczny sens entropii:}
Podczas gdy fenomenologia łączyła entropię z ciepłem, Ludwig Boltzmann ukazał jej znacznie głębszy, fundamentalny sens wynikający z fizyki statystycznej. Entropia jest miarą nieuporządkowania układu, a ściślej – miarą prawdopodobieństwa termodynamicznego danego makrostanu.
Wzór Boltzmanna, wyryty zresztą na jego nagrobku, głosi:
\[
S = k \ln \Omega
\]
gdzie $k$ to stała Boltzmanna, a $\Omega$ to liczba mikrostanów (różnych rozkładów położeń i pędów cząsteczek) realizujących dany, obserwowalny makrostan. Wzrost entropii wynika czysto z rachunku prawdopodobieństwa – układy dążą do stanów najbardziej prawdopodobnych, a stany równomierne i zdezorganizowane mogą być zrealizowane na znacznie więcej mikroskopowych sposobów niż stany wysoce uporządkowane.

\section{Siły spójności: napięcie powierzchniowe, włoskowatość. Ruchy Browna.}

\textbf{Siły spójności (kohezji):}
Płyny (głównie ciecze) utrzymują swoją gęstość dzięki istnieniu sił przyciągających między cząsteczkami tej samej substancji, zwanych siłami spójności (kohezji). Mają one krótki zasięg (rzędu rozmiarów cząsteczek) i pochodzenie elektromagnetyczne (oddziaływania van der Waalsa, wiązania wodorowe). Siły działające między cząsteczkami cieczy a ściankami naczynia nazywamy z kolei siłami przylegania (adhezji). Relacja między tymi dwoma siłami decyduje o wielu zjawiskach powierzchniowych (np. czy ciecz zwilża powierzchnię, czy nie).

\textbf{Napięcie powierzchniowe:}
Wewnątrz objętości cieczy, na wybraną cząsteczkę siły spójności od sąsiednich cząsteczek działają symetrycznie we wszystkich kierunkach – ich wypadkowa jest zerowa. Jednak cząsteczka znajdująca się na powierzchni cieczy nie ma nad sobą innych cząsteczek cieczy. W efekcie wypadkowa siła skierowana jest w głąb cieczy. Powoduje to, że ciecz dąży do minimalizacji swojej pola powierzchni (dlatego swobodne krople przyjmują kształt kuli, która ma najmniejszą powierzchnię ze wszystkich brył o tej samej objętości).
Aby powiększyć powierzchnię cieczy o $\Delta S$, należy wykonać pracę z zewnątrz, przeciwko siłom spójności. \textbf{Napięcie powierzchniowe ($\sigma$)} definiuje się jako stosunek tej pracy do przyrostu powierzchni:
\[
\sigma = \frac{\Delta W}{\Delta S}
\]
Można je także zdefiniować mechanicznie jako siłę styczną działającą na jednostkę długości brzegu powierzchni cieczy: $\sigma = F/l$ [N/m]. To ono sprawia, że na wodzie mogą utrzymywać się lekkie przedmioty (np. owady, żyletka) o gęstości większej niż woda, pod warunkiem, że nie przerwą "błony" powierzchniowej.

\textbf{Włoskowatość (Kapilarność):}
Jest to zjawisko wznoszenia się lub opadania poziomu cieczy w bardzo wąskich rurkach (kapilarach), w porównaniu do swobodnego poziomu w szerokim naczyniu. Wypływa ono z faktu, że jeśli ciecz zwilża rurkę (siły adhezji do szkła są silniejsze od sił kohezji w wodzie), tworzy się wklęsły menisk. Napięcie powierzchniowe na zakrzywionej powłoce generuje różnicę ciśnień po obu jej stronach (tzw. ciśnienie Laplace'a: $\Delta p = 2\sigma / R$, gdzie $R$ to promień krzywizny menisku). 
Wywołuje to "podsysanie" cieczy w górę rurki, do momentu zrównoważenia tego ciśnienia przez ciśnienie hydrostatyczne słupa podniesionej cieczy. Wysokość podniesienia cieczy w kapilarze opisuje wzór Jurina:
\[
h = \frac{2\sigma \cos\theta}{\rho g r}
\]
gdzie $\theta$ to kąt zwilżania, $\rho$ gęstość cieczy, a $r$ promień kapilary. Jeśli ciecz nie zwilża naczynia (np. rtęć w szkle), $\theta > 90^\circ$, $\cos\theta$ jest ujemny, menisk jest wypukły i poziom cieczy w kapilarze opada. Zjawisko to ma kolosalne znaczenie dla transportu wody w glebie oraz soków w tkankach roślinnych.

\textbf{Ruchy Browna:}
Zjawisko to, odkryte w 1827 roku przez botanika Roberta Browna, polega na ciągłym, chaotycznym i zygzakowatym ruchu bardzo drobnych (ale makroskopowych, mierzonych w mikrometrach) cząstek zawieszonych w płynie (cieczy lub gazie) – np. drobin pyłku kwiatowego w wodzie lub dymu w powietrzu. 

Początkowo traktowane jako zjawisko biologiczne, zostało ostatecznie wyjaśnione na gruncie fizyki molekularnej m.in. przez Alberta Einsteina i Mariana Smoluchowskiego w 1905 roku. Ruch ten jest bezpośrednim i naocznym dowodem słuszności kinetyczno-molekularnej teorii budowy materii. Jest on wywoływany przez ciągłe, asymetryczne bombardowanie pyłku ze wszystkich stron przez niewidoczne pod mikroskopem świetlnym, znajdujące się w nieustannym ruchu termicznym cząsteczki płynu. Ze względu na wahania statystyczne (fluktuacje), siły działające w danej ułamku sekundy z różnych stron nie równoważą się idealnie, co nadaje drobinie przypadkowy pęd. Zgodnie z teorią Einsteina-Smoluchowskiego, średni kwadrat przesunięcia cząstki w ruchach Browna jest wprost proporcjonalny do czasu obserwacji ($\langle x^2 \rangle \propto t$) oraz temperatury płynu. Analiza tych ruchów posłużyła m.in. Jeanowi Perrinowi do precyzyjnego wyznaczenia liczby Avogadra.

\section{Zjawiska transportu w gazach: dyfuzja, przewodnictwo cieplne, lepkość, przewodnictwo elektryczne.}

Układ termodynamiczny pozostawiony samemu sobie zawsze dąży do stanu równowagi (maksymalizacji entropii). Jeżeli gaz znajduje się w stanie nierównowagowym (np. występują w nim gradienty stężenia, temperatury czy prędkości makroskopowej), wewnątrz gazu samorzutnie powstają przepływy materii, energii lub pędu mające na celu zniwelowanie tych różnic. Zjawiska te nazywamy zjawiskami transportu. Ze względu na budowę mikroskopową gazów, mechanizmem transportu są zderzenia cząsteczek. Kluczowym parametrem mikroskopowym jest \textbf{średnia swobodna droga ($\lambda$)}, czyli średnia odległość, jaką cząsteczka przebywa pomiędzy dwoma kolejnymi zderzeniami.

\textbf{Dyfuzja (transport masy):}
Zachodzi, gdy w gazie istnieje gradient koncentracji (stężenia) cząsteczek. Chaotyczny ruch cieplny powoduje, że statystycznie więcej cząsteczek przemieszcza się z obszaru o wyższym stężeniu do obszaru o stężeniu niższym niż odwrotnie. Makroskopowo zjawisko to opisuje \textbf{prawo Ficka}:
\[
\vec{J} = -D \nabla n
\]
gdzie $\vec{J}$ to gęstość strumienia cząstek (liczba cząstek przechodząca przez jednostkową powierzchnię w jednostce czasu), $\nabla n$ to gradient koncentracji, a $D$ to współczynnik dyfuzji. Z teorii kinetyczno-molekularnej wynika, że $D = \frac{1}{3} \bar{v} \lambda$, gdzie $\bar{v}$ to średnia prędkość termiczna cząsteczek.

\textbf{Przewodnictwo cieplne (transport energii):}
Gdy w gazie występuje różnica temperatur, cząsteczki w cieplejszych obszarach posiadają większą średnią energię kinetyczną. Poprzez zderzenia, przekazują one część tej energii cząsteczkom wolniejszym, co makroskopowo obserwujemy jako przepływ ciepła. Opisuje to \textbf{prawo Fouriera}:
\[
\vec{q} = -\kappa \nabla T
\]
gdzie $\vec{q}$ to gęstość strumienia ciepła, a $\kappa$ to współczynnik przewodnictwa cieplnego. Mikroskopowo $\kappa = \frac{1}{3} \rho \bar{v} \lambda c_v$, gdzie $\rho$ to gęstość gazu, a $c_v$ to ciepło właściwe. Warto zauważyć, że ponieważ w gazach $\lambda \propto 1/\rho$, przewodnictwo cieplne gazów jest w szerokim zakresie niezależne od ciśnienia (co wykorzystuje się m.in. w termosach, gdzie trzeba wypompować gaz do bardzo głębokiej próżni, aby $\lambda$ stało się większe od rozmiaru naczynia).

\textbf{Lepkość / Tarcie wewnętrzne (transport pędu):}
Jeżeli warstwy gazu poruszają się z różnymi prędkościami makroskopowymi, cząsteczki przeskakujące między warstwami (dzięki ruchom termicznym) przenoszą makroskopowy pęd, dążąc do wyrównania prędkości. Efektem jest powstawanie sił tarcia wewnętrznego pomiędzy warstwami. Opisuje to \textbf{prawo Newtona dla lepkości}:
\[
\tau = \eta \frac{dv}{dx}
\]
gdzie $\tau$ to naprężenie ścinające (siła na jednostkę powierzchni), $dv/dx$ to gradient prędkości poprzecznej, a $\eta$ to współczynnik lepkości dynamicznej. Teoretycznie $\eta = \frac{1}{3} \rho \bar{v} \lambda$. W odróżnieniu od cieczy, lepkość gazów rośnie wraz ze wzrostem temperatury (bo rośnie średnia prędkość cieplna $\bar{v}$).

\textbf{Przewodnictwo elektryczne w gazach:}
W normalnych warunkach gazy są doskonałymi izolatorami (dielektrykami), ponieważ składają się z obojętnych atomów i cząsteczek. Aby gaz zaczął przewodzić prąd, musi dojść do \textbf{jonizacji} – oderwania elektronów od atomów, co generuje nośniki ładunku (elektrony swobodne i jony dodatnie).
\begin{itemize}
    \item \textbf{Wyładowania niesamoistne:} Wymagają ciągłego działania zewnętrznego czynnika jonizującego (np. promieniowania UV, rentgenowskiego, płomienia). Po usunięciu czynnika prąd zanika na skutek rekombinacji jonów.
    \item \textbf{Wyładowania samoistne:} Powstają w bardzo silnym polu elektrycznym. Rozpędzone przez pole elektrony zyskują na drodze swobodnej energię wystarczającą do jonizacji kolejnych atomów przy zderzeniach (jonizacja zderzeniowa). Następuje lawinowe namnażanie nośników ładunku (np. iskra, piorun, łuk elektryczny, jarzenie w świetlówkach).
\end{itemize}

\section{Potencjał elektryczny, natężenie pola elektrycznego. Strumień pola elektrycznego. Prawo Gaussa.}

Oddziaływanie ładunków elektrycznych (opisywane prawem Coulomba) najwygodniej jest tłumaczyć koncepcją pola fizycznego. Każdy ładunek modyfikuje przestrzeń wokół siebie, tworząc pole elektryczne, które z kolei działa siłą na inne ładunki wprowadzone w tę przestrzeń.

\textbf{Natężenie pola elektrycznego ($\vec{E}$):}
Jest to wielkość wektorowa opisująca lokalną "siłę" pola elektrycznego. Definiujemy ją jako stosunek siły elektrostatycznej ($\vec{F}$), jakiej doznaje dodatni ładunek próbny ($q_0$) umieszczony w danym punkcie, do wartości tego ładunku (przy założeniu $q_0 \to 0$, aby ładunek próbny nie zaburzał badanego układu):
\[
\vec{E} = \frac{\vec{F}}{q_0}
\]
Jednostką natężenia jest niuton na kulomb (N/C) lub wolt na metr (V/m). Zgodnie z prawem Coulomba, natężenie pola wokół ładunku punktowego $Q$ wynosi $\vec{E} = \frac{1}{4\pi\varepsilon_0} \frac{Q}{r^2} \hat{r}$, gdzie $\hat{r}$ to wektor jednostkowy wskazujący od ładunku źródłowego. Pole elektryczne spełnia zasadę superpozycji – całkowite natężenie od układu ładunków jest sumą wektorową natężeń od ładunków indywidualnych.

\textbf{Potencjał elektryczny ($V$):}
Pole elektrostatyczne (stałe w czasie) jest polem sił zachowawczych, tzn. praca przeniesienia ładunku między dwoma punktami zależy tylko od ich położeń, a nie od drogi. Pozwala to na zdefiniowanie energii potencjalnej ($E_p$). \textbf{Potencjał elektryczny} definiuje się jako energię potencjalną przypadającą na jednostkowy ładunek próbny:
\[
V = \frac{E_p}{q}
\]
Potencjał jest wielkością skalarną (jednostka to wolt, V = J/C), co znacznie ułatwia obliczenia (sumowanie algebraiczne zamiast wektorowego). Różnica potencjałów między dwoma punktami to \textbf{napięcie elektryczne ($U$)}, które określa pracę pola: $W_{A\to B} = q(V_A - V_B) = qU$. Natężenie i potencjał są ze sobą ściśle powiązane gradientem:
\[
\vec{E} = -\nabla V
\]
Wektor natężenia pola wskazuje kierunek najszybszego spadku potencjału i jest zawsze prostopadły do powierzchni ekwipotencjalnych.

\textbf{Strumień pola elektrycznego ($\Phi_E$):}
Z matematycznego punktu widzenia, strumień wektora przez zadaną powierzchnię informuje, "jak dużo" pola przebija tę powierzchnię. Jest on zdefiniowany jako całka powierzchniowa:
\[
\Phi_E = \int_S \vec{E} \cdot d\vec{S}
\]
Gdzie $d\vec{S}$ to wektor powierzchniowy (skierowany normalnie do powierzchni, o wartości równej polu elementu powierzchni). Optycznie, reprezentuje on liczbę linii pola sił przenikających przez daną powierzchnię.

\textbf{Prawo Gaussa:}
To fundamentalne prawo elektrostatyki, stanowiące jedno z czterech równań Maxwella. Łączy ono całkowity strumień pola elektrycznego przenikający przez dowolną zamykającą się powierzchnię (powierzchnię Gaussa) z całkowitym ładunkiem zamkniętym w jej wnętrzu:
\[
\oint_S \vec{E} \cdot d\vec{S} = \frac{Q_{wew}}{\varepsilon_0}
\]
gdzie $\varepsilon_0$ to przenikalność elektryczna próżni, a $Q_{wew}$ to algebraiczna suma ładunków wewnątrz.
Prawo Gaussa jest zawsze prawdziwe, ale jako narzędzie obliczeniowe staje się niezwykle użyteczne przy wyznaczaniu natężenia $\vec{E}$ dla rozkładów ładunków o wysokiej symetrii (np. naładowana nieskończona płaszczyzna, długa nić, naładowana sfera). Dzięki symetrii wektor $\vec{E}$ jest stały co do wartości na powierzchni Gaussa, co pozwala na trywialne obliczenie całki i wyznaczenie pola. Postać różniczkowa prawa Gaussa to $\nabla \cdot \vec{E} = \frac{\rho}{\varepsilon_0}$ (rozbieżność pola zależy od lokalnej gęstości ładunku).

\section{Dielektryk w polu elektrycznym. Polaryzacja. Elektrostrykcja i piezoelektryczność. Ferroelektryki. Pojemność.}

\textbf{Dielektryk w polu elektrycznym i mechanizmy polaryzacji:}
Dielektryki to izolatory, czyli materiały, w których elektrony są silnie związane z jądrami atomowymi i brak w nich nośników ładunku zdolnych do przemieszczania się na makroskopowe odległości. Choć dielektryk nie przewodzi prądu, silnie reaguje na zewnętrzne pole elektryczne, ulegając \textbf{polaryzacji}. Polega ona na mikroskopowym rozsunięciu (przesunięciu przestrzennym) ładunków dodatnich i ujemnych w materiale. Wektorową miarą tego zjawiska jest wektor polaryzacji $\vec{P}$ (moment dipolowy jednostki objętości). Wyróżniamy polaryzację:
\begin{itemize}
    \item \textbf{Elektronową/jonową (indukowaną):} pod wpływem zewnętrznego pola jądro atomowe i chmura elektronowa przemieszczają się lekko w przeciwnych kierunkach, indukując moment dipolowy (występuje w atomach niepolarnych).
    \item \textbf{Orientacyjną (dipolową):} występuje w dielektrykach zbudowanych z cząsteczek polarnych (np. woda), posiadających własne, stałe momenty dipolowe. Chaotycznie rozrzucone w wyniku ruchów termicznych dipole ulegają częściowemu uporządkowaniu zgodnie z liniami zewnętrznego pola. Zjawisko to silnie zależy od temperatury.
\end{itemize}
Polaryzacja sprawia, że na powierzchniach (lub w objętości przy niejednorodnym rozkładzie) dielektryka pojawiają się ładunki polaryzacyjne (związane). Wytwarzają one własne, przeciwne pole elektryczne, które \textbf{osłabia} pole zewnętrzne we wnętrzu materiału. Pole wewnątrz dielektryka $E$ jest $\varepsilon_r$ razy mniejsze od pola zewnętrznego $E_0$ ($E = E_0/\varepsilon_r$, gdzie $\varepsilon_r$ to przenikalność względna).

\textbf{Zjawiska elektrostrykcji i piezoelektryczności:}
\begin{itemize}
    \item \textbf{Elektrostrykcja} to niewielkie, odwracalne odkształcenie (zmiana wymiarów) każdego dielektryka po umieszczeniu go w zewnętrznym polu elektrycznym. Odkształcenie to jest proporcjonalne do kwadratu natężenia pola ($E^2$), a więc zachowuje znak niezależnie od kierunku pola.
    \item \textbf{Piezoelektryczność} to zjawisko polegające na pojawieniu się różnicy potencjałów (ładunku powierzchniowego) na kryształach pod wpływem naprężenia mechanicznego. Zjawisko występuje tylko w kryształach o określonej symetrii (nie posiadających środka symetrii, np. kwarc). Występuje tu też \textbf{odwrotny efekt piezoelektryczny}: kryształ poddany działaniu zewnętrznego napięcia zmienia swoje wymiary. Zależność jest w nim liniowa względem $E$, czyli odwrócenie pola odwraca kierunek odkształcenia (ściskanie zamienia się w rozciąganie).
\end{itemize}

\textbf{Ferroelektryki:}
To specyficzna, mała klasa dielektryków (np. tytanian baru), które wykazują własność spontanicznej polaryzacji nawet przy braku zewnętrznego pola elektrycznego (istnieją w nich makroskopowe obszary jednorodnie spolaryzowane zwane domenami). Posiadają nieliniową zależność polaryzacji od pola zewnętrznego oraz wykazują zjawisko histerezy dielektrycznej. Ich przenikalność elektryczna ($\varepsilon_r$) jest gigantyczna (często rzędu $10^3 - 10^4$). Cechy ferroelektryczne zanikają powyżej charakterystycznej \textbf{temperatury Curie}, po której materiał staje się zwykłym paraelektrykiem.

\textbf{Pojemność kondensatora z dielektrykiem:}
Kondensator to układ przewodników zdolny do gromadzenia ładunku, charakteryzujący się pojemnością $C = Q/U$. Wprowadzenie dielektryka w całą przestrzeń między okładkami kondensatora powoduje (dzięki osłabieniu pola elektrycznego przez ładunki polaryzacyjne) obniżenie różnicy potencjałów (napięcia $U$) przy zachowaniu tego samego ładunku na okładkach $Q$. W efekcie, pojemność kondensatora wzrasta $\varepsilon_r$ razy względem pojemności w próżni:
\[
C = \varepsilon_r C_0
\]
Dlatego stałe materiały dielektryczne (np. ceramika, tantal, tworzywa sztuczne) są niezbędne w konstrukcji elementów elektronicznych, aby zminiaturyzować ich rozmiary przy dużej pojemności i jednoczesnym podniesieniu wytrzymałości na przebicie elektryczne.

\section{Klasyczny model przewodnictwa elektrycznego.}

Zaproponowany przez Paula Drudego na początku XX wieku (a następnie rozwijany przez Lorentza) klasyczny model gazu elektronowego stanowił pierwszą próbę zrozumienia zjawiska przewodnictwa prądu w metalach w oparciu o mikroskopową strukturę materii.

\textbf{Założenia modelu Drudego:}
Model ten zakłada, że w objętości metalu wyodrębnić można nieruchomą sieć krystaliczną zbudowaną z ciężkich, dodatnich jonów oraz "gaz" uwolnionych elektronów walencyjnych (tzw. elektronów swobodnych). Drude przyjął dla tego gazu prawa termodynamiki klasycznej, identyczne jak dla gazu doskonałego:
\begin{itemize}
    \item Elektrony traktowane są jako klasyczne punkty materialne, a ich wzajemne oddziaływania są zaniedbane.
    \item Elektrony ulegają zderzeniom tylko ze stacjonarnymi jonami sieci (które traktowane są jako nieprzenikliwe kule).
    \item Pomiędzy zderzeniami elektrony poruszają się swobodnie (po liniach prostych przy braku sił zewnętrznych). Czas między dwoma zderzeniami charakteryzowany jest parametrem zwanym \textbf{czasem relaksacji ($\tau$)}.
    \item Prędkość zyskiwana z oddziaływań termicznych (rozład Maxwella) jest rzędu $10^5$ m/s, a kierunek ruchu elektronu tuż po zderzeniu z jonem staje się całkowicie losowy.
\end{itemize}

\textbf{Mechanizm przewodnictwa:}
Przy braku pola elektrycznego chaotyczny ruch elektronów we wszystkich kierunkach sprawia, że wypadkowy makroskopowy prąd wynosi zero. Sytuacja zmienia się, gdy wewnątrz przewodnika przyłożymy pole elektryczne $\vec{E}$. Wówczas, podczas ruchu między zderzeniami (zgodnie z II zasadą dynamiki Newtona: $m_e \vec{a} = -e\vec{E}$), elektrony zyskują stałe, ukierunkowane przyspieszenie. W ten sposób, na potężną i chaotyczną prędkość termiczną nakłada się niewielka (często rzędu ułamków milimetra na sekundę), jednostajna **prędkość unoszenia ($\vec{v}_d$)**.
Średnia prędkość, jaką elektron uzyskuje przed zderzeniem trwającym przeciętnie czas $\tau$, to:
\[
\vec{v}_d = -\frac{e\tau}{m_e} \vec{E}
\]
Prąd gęstości ładunku swobodnego definiuje się jako $\vec{j} = -n e \vec{v}_d$, gdzie $n$ to koncentracja elektronów. Podstawiając wyrażenie na $\vec{v}_d$ otrzymujemy \textbf{prawo Ohma w postaci mikroskopowej}:
\[
\vec{j} = \left(\frac{n e^2 \tau}{m_e}\right) \vec{E} = \sigma \vec{E}
\]
gdzie $\sigma = \frac{n e^2 \tau}{m_e}$ to wyznaczona analitycznie \textbf{konduktywność (przewodność właściwa)} materiału.

\textbf{Sukcesy i porażki klasycznego modelu:}
Model Drudego odniósł wielki sukces, zgrabnie wyprowadzając od podstaw mikroskopowych empiryczne prawo Ohma (wykazując, że prąd jest wprost proporcjonalny do pola). Ponadto zdołał poprawnie uzasadnić prawo Wiedemanna-Franza (stałość stosunku przewodnictwa cieplnego do elektrycznego u wszystkich metali), wiążąc transport ładunku z transportem energii kinetycznej.

Z czasem jednak zaobserwowano fundamentalne porażki modelu:
1) \textbf{Problem ciepła właściwego:} Wg zasady ekwipartycji, klasyczny gaz elektronowy powinien drastycznie (o ok. $50\%$) zwiększać ciepło właściwe metali względem izolatorów. Doświadczenia wykazały jednak, że elektronowy wkład do ciepła właściwego jest niezwykle mały.
2) \textbf{Zależność oporu od temperatury:} Model klasyczny przewiduje, że prędkość termiczna $\bar{v} \propto \sqrt{T}$, co prowadzi do oporu elektrycznego proporcjonalnego do $\sqrt{T}$. Eksperymenty pokazywały u metali zależność ściśle liniową ($R \propto T$).
Powyższe problemy klasycznego modelu zostały rozwiązane dopiero z nastaniem fizyki kwantowej – zastąpiono statystykę klasyczną statystyką Fermiego-Diraca (gdzie w transporcie bierze udział tylko garstka elektronów w pobliżu powierzchni Fermiego) oraz wprowadzono falową naturę elektronu rozpraszającego się na drganiach sieci krystalicznej (fononach).

\section{Prawo Ohma. Siła elektromotoryczna. Reguły Kirchhoffa. Praca i moc prądu.}

\textbf{Prawo Ohma:}
W sformułowaniu makroskopowym prawo to określa relację między napięciem przyłożonym do przewodnika a prądem przez niego płynącym. Georg Ohm stwierdził, że w stałej temperaturze, natężenie prądu ($I$) płynącego przez przewodnik (tzw. opornik omowy) jest wprost proporcjonalne do napięcia ($U$) na jego końcach. Opór elektryczny ($R$) zdefiniowano jako stałą proporcjonalności w tym równaniu:
\[
I = \frac{U}{R} \quad \implies \quad U = I R
\]
Opór elektryczny kawałka przewodnika (mierzony w omach, $\Omega$) zależy od jego geometrii i materiału: $R = \rho \frac{l}{S}$, gdzie $l$ to długość, $S$ pole przekroju, a $\rho$ to rezystywność (opór właściwy) materiału zależny od temperatury.

\textbf{Siła elektromotoryczna (SEM):}
Przepływ prądu stałego w zamkniętym obwodzie powoduje ciągłe straty energii układu ładunków na cieplne opory w przewodniku (spadek potencjału). Aby utrzymać ciągły przepływ, obwód musi posiadać źródło ładujące te elektrony nową energią i pompujące je "pod górę" potencjału, wbrew elektrostatyce. Wykonują to siły nieelektryczne (np. chemiczne w bateriach, indukcyjne w prądnicach, termiczne w termoparach).
\textbf{Siła elektromotoryczna ($\mathcal{E}$)} to praca wykonana przez siły nieelektryczne we wnętrzu źródła przy przeniesieniu jednostkowego dodatniego ładunku z bieguna ujemnego na biegun dodatni. Mierzy się ją w woltach. Obwód zamknięty, zawierający źródło o własnym oporze wewnętrznym $r_w$, opisuje uogólnione prawo Ohma: $I = \frac{\mathcal{E}}{R + r_w}$.

\textbf{Reguły Kirchhoffa:}
Służą do analizy układów elektrycznych o złożonej topologii (wielu oczkach), których nie da się łatwo zwinąć do pojedynczego opornika. Wynikają bezpośrednio z fundamentalnych zasad fizyki klasycznej.
\begin{itemize}
    \item \textbf{I Prawa Kirchhoffa (reguła węzłów):} W obwodzie ze stałym prądem ładunek nie może się nigdzie gromadzić ani znikać. Jest to zapis \textbf{zasady zachowania ładunku}. Suma natężeń prądów wpływających do dowolnego węzła musi równać się sumie prądów wypływających z tego węzła ($\sum I_{wp} = \sum I_{wyp}$).
    \item \textbf{II Prawo Kirchhoffa (reguła oczek):} Praca sił zewnętrznych po drodze zamkniętej równa się zyskowi potencjału od źródeł, ponieważ całkowita zmiana potencjału przy powrocie do punktu wyjścia wynosi zero. Jest to zapis \textbf{zasady zachowania energii}. W zamkniętym oczku obwodu, suma spadków napięć na wszystkich elementach biernych jest równa sumie algebraicznej sił elektromotorycznych działających w tym oczku ($\sum I_k R_k = \sum \mathcal{E}_i$). Należy przy tym ściśle przestrzegać konwencji znaków ustalając kierunek obiegu.
\end{itemize}

\textbf{Praca i moc prądu elektrycznego:}
Praca wykonywana przez siły pola elektrycznego przemieszczające ładunek $q$ przez spadek napięcia $U$ wynosi $W = qU$. Biorąc pod uwagę, że prąd $I = q/t$, praca wykonana w czasie $t$ wyraża się wzorem:
\[
W = U \cdot I \cdot t
\]
Związana z tym moc (szybkość wykonywania pracy) to $P = W/t = U \cdot I$. Korzystając z Prawa Ohma ($U=IR$), możemy wyrazić moc wydzielaną w elemencie oporowym w różnych, alternatywnych formach:
\[
P = U \cdot I = I^2 R = \frac{U^2}{R}
\]
Energia opisanego wyżej prądu zamieniana jest podczas przechodzenia elektronów przez siatkę krystaliczną na wzbudzenie drgań jonów (ogrzewanie przewodnika) – jest to tak zwane \textbf{Ciepło Joule'a-Lenza}. Wzór ten uzasadnia, dlaczego prąd z elektrowni przesyłany jest liniami wysokiego napięcia – mniejszy natężenie prądu ($I$) to diametralnie (w kwadracie) mniejsze straty energii na oporach samych linii transmisyjnych.

\section{Prawo Ampere'a; oddziaływanie dwóch przewodników z prądem. Indukcja elektromagnetyczna; prawo Faradaya.}

\textbf{Prawo Ampere'a:}
Fundamentalne prawo opisujące pole magnetyczne wytwarzane przez stacjonarne przepływy ładunków elektrycznych. Wiąże ono krążenie (cyrkulację) wektora indukcji magnetycznej ($\vec{B}$) wzdłuż dowolnego, zamkniętego konturu (krzywej matematycznej $C$) z całkowitym algebraicznym prądem przenikającym przez powierzchnię rozpiętą na tym konturze. W postaci całkowej (w próżni):
\[
\oint_C \vec{B} \cdot d\vec{l} = \mu_0 \sum I_{wew}
\]
gdzie $\mu_0$ to przenikalność magnetyczna próżni. Równanie to jest potężnym narzędziem (podobnie jak prawo Gaussa w elektrostatyce) do obliczania pola magnetycznego $\vec{B}$ z wysoce symetrycznych układów (np. prostoliniowy, nieskończenie długi przewodnik; długi zwojnik z prądem). Postać różniczkowa, stanowiąca ramy równań Maxwella, to $\nabla \times \vec{B} = \mu_0 \vec{j}$.

\textbf{Oddziaływanie dwóch przewodników z prądem:}
Jeżeli w zewnętrznym polu magnetycznym znajdzie się przewodnik z prądem, pole działa na płynące w nim ładunki, objawiając się w skali makro siłą elektrodynamiczną (makroskopowy skutek siły Lorentza). Wynosi ona $\vec{F} = I \vec{l} \times \vec{B}$.
Zjawisko to jest najdobitniej widoczne przy oddziaływaniu na siebie dwóch nieskończenie długich, równoległych, prostoliniowych przewodników z prądem, oddalonych o dystans $d$. Jeden przewodnik generuje wirowe pole magnetyczne (wynikające z prawa Ampere'a o natężeniu $B = \frac{\mu_0 I_1}{2\pi d}$), wewnątrz którego znajduje się drugi przewodnik doznający siły:
\[
\frac{F}{l} = \frac{\mu_0 I_1 I_2}{2\pi d}
\]
gdzie $F/l$ to siła przypadająca na jednostkę długości drutu. Reguły wektorowe dowodzą, że prądy płynące w tę \textbf{samą stronę przyciągają się}, a w strony \textbf{przeciwne – odpychają się} ze wspomnianą siłą. Przez dekady zjawisko to posłużyło jako definicja 1 Ampera w układzie SI.

\textbf{Indukcja elektromagnetyczna; prawo Faradaya:}
W 1831 roku Michael Faraday wykazał kluczowe zjawisko odwrotne do prawa Ampere'a – udowodnił, że zmienne pole magnetyczne potrafi generować prąd. Jeżeli w obszarze przewodzącego oczka zmienia się przenikający przezeń strumień indukcji magnetycznej ($\Phi_B = \int \vec{B} \cdot d\vec{S}$), w obwodzie wytwarzana jest tak zwana indukowana siła elektromotoryczna. Zjawisko to, nazywane indukcją elektromagnetyczną, zmatematyzowane jest \textbf{prawem Faradaya}:
\[
\mathcal{E}_{ind} = -\frac{d\Phi_B}{dt}
\]
Strumień pola magnetycznego $\Phi_B = B \cdot S \cdot \cos(\alpha)$ może zmieniać się z upływem czasu na trzy zasadnicze sposoby: poprzez obracanie cewki w stałym polu (zmiana kąta $\alpha$ – podstawa działania światowych elektrowni produkujących prąd przemienny prądnicami), poprzez zmianę pola powierzchni $S$ (rozszerzanie lub wyciąganie przewodnika z pola), lub też poprzez fizyczną zmianę wartości lokalnego pola $\vec{B}$ (np. zbliżanie magnesu).

Obecny we wzorze \textbf{znak minus} jest bezpośrednią konsekwencją zasady zachowania energii, nazwaną na cześć twórcy \textbf{regułą Lenza} (regułą przekory). Mówi ona, że prąd zaindukowany w obwodzie przez indukcję płynie zawsze w takim kierunku, by wytwarzane przez niego (indukowane) pole magnetyczne przeciwstawiało się zmianie strumienia zewnętrznego pola magnetycznego, która wywołała powstanie tego prądu.

\section{Materia w polu magnetycznym. Diamagnetyki, paramagnetyki, ferromagnetyki.}

Obecność materii makroskopowej zasadniczo modyfikuje pole magnetyczne w stosunku do pola w próżni. Wektor indukcji magnetycznej $\vec{B}$ wewnątrz ośrodka zależy od zewnętrznego natężenia pola magnetycznego $\vec{H}$ oraz od namagnesowania samego ośrodka $\vec{M}$ (wektora polaryzacji magnetycznej, czyli sumarycznego momentu magnetycznego jednostki objętości). Związek ten wyraża równanie:
\[
\vec{B} = \mu_0 (\vec{H} + \vec{M})
\]
W większości materiałów izotropowych wektor namagnesowania jest proporcjonalny do natężenia pola: $\vec{M} = \chi \vec{H}$, gdzie $\chi$ to bezwymiarowa podatność magnetyczna. Wprowadzając przenikalność magnetyczną względną $\mu_r = 1 + \chi$, otrzymujemy znaną relację: $\vec{B} = \mu_0 \mu_r \vec{H}$. W zależności od wartości i znaku podatności $\chi$, a także od struktury domenowej, materię dzielimy na trzy główne klasy.

\textbf{Diamagnetyki ($\chi < 0$, $\mu_r < 1$):}
Podatność magnetyczna diamagnetyków jest ujemna i bardzo mała (rzędu $-10^{-5}$). Oznacza to, że materiał ten nieznacznie osłabia pole magnetyczne w swoim wnętrzu. Na poziomie mikroskopowym atomy diamagnetyka (z zamkniętymi powłokami elektronowymi) nie posiadają własnego stałego momentu magnetycznego. Diamagnetyzm jest uniwersalnym zjawiskiem, będącym bezpośrednią konsekwencją reguły Lenza na poziomie atomowym. Włączenie zewnętrznego pola magnetycznego indukuje w atomach dodatkowy, makroskopowy moment obrotowy elektronów wokół jądra (tzw. precesja Larmora), który wytwarza wtórne pole skierowane zawsze przeciwnie do pola wymuszającego. Makroskopowym efektem jest bardzo słabe odpychanie (wypychanie) diamagnetyków z obszarów silniejszego pola magnetycznego. Przykłady: woda, miedź, bizmut, szkło.

\textbf{Paramagnetyki ($\chi > 0$, $\mu_r > 1$):}
Podatność paramagnetyków jest dodatnia i również niewielka ($10^{-5}$ do $10^{-3}$). Atomy (lub cząsteczki) tych substancji posiadają nieskompensowane spiny elektronowe, a więc trwałe, własne momenty magnetyczne. W braku pola zewnętrznego momenty te, na skutek chaotycznych ruchów cieplnych, są zorientowane losowo, a całkowite namagnesowanie wynosi zero. Zewnętrzne pole $\vec{B}$ wymusza na tych dipolach moment sił ($\vec{M} = \vec{\mu} \times \vec{B}$), dążąc do ich ustawienia wzdłuż linii pola, co wzmacnia pole wewnątrz materiału. Efekt ten rywalizuje jednak z niszczącym uporządkowanie ruchem termicznym, dlatego podatność paramagnetyczna silnie zależy od temperatury, co opisuje \textbf{prawo Curie}: $\chi = C/T$ (gdzie $C$ to stała materiałowa). Przykłady: aluminium, platyna, ciekły tlen.

\textbf{Ferromagnetyki ($\chi \gg 1$, $\mu_r$ rzędu $10^3 - 10^5$):}
Materiały te potrafią potężnie wzmocnić pole magnetyczne i wykazują zdolność do trwałego namagnesowania. Mechanizm nie ma klasycznego odpowiednika – wynika z kwantowo-mechanicznego oddziaływania wymiennego między spinami sąsiednich atomów, które zmusza je do równoległego ustawienia (nawet przy braku pola zewnętrznego). Prowadzi to do powstawania \textbf{domen magnetycznych} (mikroskopowych obszarów o 100\% uporządkowaniu). Rozrost domen pod wpływem zewnętrznego pola jest procesem nieliniowym i wykazuje zjawisko \textbf{histerezy magnetycznej} (opóźnienia namagnesowania względem pola $\vec{H}$). Po wyłączeniu pola zewnętrznego ferromagnetyk zachowuje \textbf{pozostałość magnetyczną} (staje się magnesem trwałym). Do jego rozmagnesowania niezbędne jest przyłożenie przeciwnego \textbf{pola koercji}. Własności ferromagnetyczne zanikają całkowicie powyżej charakterystycznej temperatury Curie (powyżej niej materiał staje się zwykłym paramagnetykiem, gdyż energia termiczna rozrywa wiązania wymienne). Przykłady: żelazo, kobalt, nikiel i ich stopy.

\section{Drgania własne i wymuszone w obwodach LC i RLC. Rezonans obwodu.}

Obwody elektryczne zawierające cewki (L) i kondensatory (C) są elektromagnetycznym odpowiednikiem mechanicznych oscylatorów. Pozwalają one na magazynowanie energii naprzemiennie w polu elektrycznym kondensatora ($E_E = q^2/2C$) oraz w polu magnetycznym cewki ($E_B = Li^2/2$).

\textbf{Drgania własne nietłumione w obwodzie LC:}
Rozpatrzmy wyidealizowany obwód z naładowanym kondensatorem $C$ zwartym przez cewkę $L$ o zerowym oporze czynnym ($R=0$). Z II prawa Kirchhoffa suma spadków napięć równa jest zeru:
\[
U_L + U_C = 0 \implies L \frac{di}{dt} + \frac{q}{C} = 0
\]
Ponieważ natężenie prądu $i = \frac{dq}{dt}$, otrzymujemy równanie różniczkowe drgań harmonicznych:
\[
\frac{d^2q}{dt^2} + \omega_0^2 q = 0
\]
gdzie $\omega_0 = \frac{1}{\sqrt{LC}}$ to tak zwana \textbf{częstość własna obwodu LC} (wzór Thomsona). Ładunek na kondensatorze, jak i prąd w obwodzie, oscylują sinusoidalnie z tą częstością $q(t) = Q_0 \cos(\omega_0 t)$. Całkowita energia układu $E = E_E + E_B$ jest w tym procesie ściśle zachowana.

\textbf{Drgania własne tłumione w obwodzie RLC:}
W rzeczywistości każdy obwód posiada rezystancję $R$ (opór drutu cewki, przewodów). Obecność rezystora powoduje dyssypację (rozpraszanie) energii elektromagnetycznej w postaci ciepła Joule'a-Lenza ($P = i^2 R$). Równanie obwodu modyfikuje się do:
\[
L \frac{d^2q}{dt^2} + R \frac{dq}{dt} + \frac{q}{C} = 0
\]
Rozwiązaniem jest funkcja wykazująca \textbf{drgania tłumione}: $q(t) = Q_0 e^{-\beta t} \cos(\omega t + \phi)$, gdzie $\beta = R/(2L)$ to współczynnik tłumienia, a częstość drgań tłumionych wynosi $\omega = \sqrt{\omega_0^2 - \beta^2}$. Przy bardzo dużym oporze ($R > 2\sqrt{L/C}$) układ traci zdolność do oscylacji – następuje pełzające, aperiodyczne rozładowanie kondensatora.

\textbf{Drgania wymuszone i rezonans obwodu:}
Aby podtrzymać drgania w obwodzie RLC, należy włączyć do niego źródło zmiennego napięcia (siły elektromotorycznej), $\mathcal{E}(t) = \mathcal{E}_0 \cos(\omega t)$. Po krótkim czasie trwania stanów przejściowych (które szybko gasną wskutek tłumienia), obwód zacznie drgać z ustaloną częstością źródła wymuszającego $\omega$. Analizując obwód z użyciem rachunku wskazówkowego lub liczb zespolonych, wyznacza się całkowity opór obwodu dla prądu przemiennego, zwany \textbf{impedancją ($Z$)}:
\[
Z = \sqrt{R^2 + \left(\omega L - \frac{1}{\omega C}\right)^2}
\]
Amplituda prądu ustalonego wynosi $I_0 = \mathcal{E}_0 / Z$.
Zjawisko \textbf{rezonansu napięć} (dla obwodu szeregowego) następuje, gdy reaktancja indukcyjna i pojemnościowa wzajemnie się znoszą:
\[
\omega L = \frac{1}{\omega C} \implies \omega = \frac{1}{\sqrt{LC}} = \omega_0
\]
W rezonansie zachodzą spektakularne zjawiska:
1) Impedancja obwodu staje się minimalna, czysto rzeczywista i równa $Z = R$.
2) Zasilające napięcie i płynący prąd stają się w fazie ($\phi = 0$).
3) Amplituda prądu osiąga w obwodzie maksimum. W konsekwencji napięcia na cewce i na kondensatorze mogą wielokrotnie (tzw. dobroć obwodu razy) przekraczać amplitudę napięcia zasilającego. Stan ten zwany jest \textbf{przepięciem rezonansowym} i z jednej strony bywa destrukcyjny w liniach przesyłowych, z drugiej zaś jest podstawowym mechanizmem wybierania pasma radiowego w strojeniu odbiorników RTV.

\section{Zasada superpozycji fal; fale stojące, dudnienia.}

\textbf{Zasada superpozycji fal:}
Dla ogromnej większości zjawisk falowych w przyrodzie (fale na wodzie, dźwięk, fale elektromagnetyczne w próżni i środowiskach liniowych) spełniona jest empiryczna zasada superpozycji. Mówi ona, że jeżeli w ośrodku rozchodzą się dwie (lub więcej) fale, to w miejscu ich nakładania się na siebie, całkowite wychylenie ośrodka (lub natężenie pola) w danym punkcie w danym czasie jest równe wektorowej (dla fal polaryzowanych, elektromagnetycznych) lub algebraicznej (dla fal skalarnych, ciśnieniowych) sumie wychyleń powodowanych przez każdą z fal osobno.
\[
y_{wyp}(x,t) = y_1(x,t) + y_2(x,t) + \dots + y_n(x,t)
\]
Matematycznie zasada ta wynika z faktu, że równanie falowe $\Delta \psi = \frac{1}{v^2} \frac{\partial^2 \psi}{\partial t^2}$ jest równaniem różniczkowym cząstkowym liniowym, a zatem suma jego rozwiązań również jest jego rozwiązaniem. Fale po przeniknięciu przez siebie zachowują swój pierwotny kształt, amplitudę i kierunek bez zmian.

\textbf{Fale stojące:}
Stanowią najprostszy, ale fundamentalnie ważny przypadek superpozycji dwóch idealnie dopasowanych fal harmonicznych. Fala stojąca powstaje w wyniku nałożenia się dwóch fal biegnących o jednakowych amplitudach i jednakowych częstościach, rozchodzących się w \textbf{przeciwnych} kierunkach. Sytuacja ta zachodzi najczęściej podczas nakładania się fali padającej i fali odbitej prostopadle od przeszkody. Suma dwóch takich równań falowych to:
\[
y(x,t) = A\sin(\omega t - kx) + A\sin(\omega t + kx) = [2A\cos(kx)] \cdot \sin(\omega t)
\]
Powstałe równanie \textbf{nie jest już falą biegnącą} (nie zawiera wyrażenia typu $x \pm vt$). Funkcja posiada wyseparowaną część przestrzenną i czasową. Charakteryzuje się występowaniem stacjonarnych maksimów, zwanych \textbf{strzałkami} (miejsca, dla których $|\cos(kx)| = 1$, mają amplitudę $2A$), oraz punktów zerowych, zwanych \textbf{węzłami} (gdzie $\cos(kx) = 0$, amplituda na stałe wynosi zero). Węzły oddalone są od siebie co $\lambda/2$. Kluczową cechą fal stojących jest to, że \textbf{nie przenoszą one energii} w przestrzeni; energia zmagazynowana między węzłami po prostu oscyluje między energią kinetyczną a potencjalną sprężystości (lub odpowiednio polem E i B).

\textbf{Dudnienia:}
Zjawisko to powstaje jako wynik superpozycji dwóch fal biegnących w tym \textbf{samym kierunku}, różniących się zaledwie nieznacznie swoją częstotliwością ($f_1 \approx f_2$). Wykorzystując tożsamości trygonometryczne do sumy fal o częstościach $\omega_1$ i $\omega_2$, otrzymujemy falę, której wyrażenie opisuje "nośną" szybko drgającą z częstością średnią $\bar{\omega} = (\omega_1 + \omega_2)/2$, której amplituda makroskopowo i wolno pulsuje (jest obwiedziona) funkcją modulującą o częstości $\Delta\omega = (\omega_1 - \omega_2)/2$.
Ucho ludzkie lub miernik rejestruje wówczas natężenie sygnału proporcjonalne do kwadratu tej amplitudy. Ze względu na trygonometryczną własność podwajania kąta przy kwadracie kosinusa, częstość tych modulacji głośności nosi nazwę \textbf{częstości dudnień} i wyraża się wprost różnicą:
\[
f_{dud} = |f_1 - f_2|
\]
Zjawisko dudnień jest powszechnie stosowane m.in. w akustyce do strojenia instrumentów "na zero" (zbliżając do siebie częstości zmniejszamy powoli częstość dudnień aż do jej całkowitego zaniku), oraz w radiofonii w układach superheterodynowych do transpozycji sygnałów na użyteczne częstotliwości pośrednie.

\section{Interferencja fal z dwóch identycznych źródeł punktowych. Doświadczenie Younga. Interferometr Michelsona.}

\textbf{Interferencja fal i warunki spójności:}
Interferencja to wysoce szczególny przypadek superpozycji fal, polegający na powstawaniu stabilnego, przestrzennego rozkładu maksimów (wzmocnień) i minimów (wygaszeń) natężenia fali. Fale emitowane przez niezależne zwykłe źródła (np. dwie żarówki) nakładają się losowo, dając po uśrednieniu w czasie po prostu sumę swych natężeń (jasności). Aby zaszła obserwowalna interferencja, fale muszą być \textbf{spójne (koherentne)}, co oznacza, że zachowują stałą, niezmienną w czasie różnicę faz pomiędzy sobą. Najprościej osiąga się to rozdzielając jedną falę z jednego pierwotnego źródła na dwie części (tzw. rozszczepienie amplitudy lub czoła fali) i nakładając je na siebie z powrotem.

Warunek na \textbf{wzmocnienie (maksimum) i wygaszenie (minimum)} fal wynika z różnicy dróg optycznych ($\Delta r$) przebytych przez obie spójne fale. Ponieważ pełna długość fali $\lambda$ odpowiada przesunięciu fazy o $2\pi$:
\begin{itemize}
    \item \textbf{Wzmocnienie:} następuje, gdy fale spotykają się "w fazie" (grzbiet pokrywa grzbiet), co matematycznie oznacza różnicę dróg stanowiącą wielokrotność całkowitą długości fali: $\Delta r = n\lambda$ (gdzie $n = 0, 1, 2, \dots$).
    \item \textbf{Wygaszenie:} następuje, gdy fale spotykają się "w przeciwfazie" (grzbiet pokrywa dolinę): $\Delta r = (n + \frac{1}{2})\lambda$. Wtedy natężenie wypadkowe spada do zera.
\end{itemize}

\textbf{Doświadczenie Younga (1801):}
Ten epokowy eksperyment stanowił koronny, bezdyskusyjny dowód na falową naturę światła, obalając dominującą korpuskularną teorię Newtona. Thomas Young przepuścił światło ze spójnego, punktowego źródła przez barierę z dwoma bliskimi siebie (odległymi o $d$) szczelinami $S_1$ i $S_2$. Zgodnie z zasadą Huygensa, szczeliny stały się nowymi, idealnie spójnymi (w przestrzeni) emiterami fal kulistych. Za barierą fale nałożyły się, rzutując na oddalony o odległość $L$ ekran układ na przemian jasnych i ciemnych prążków. 
Jeśli $L \gg d$, to różnica dróg do punktu na ekranie (odległego o $y$ od środka układu) wynosi w przybliżeniu $\Delta r = d \sin\theta \approx d \frac{y}{L}$. Podstawiając to do równania na wzmocnienie, Young wyznaczył położenie jasnych prążków: $y_n = n \frac{\lambda L}{d}$. Otrzymany w tym doświadczeniu rozkład natężeń na ekranie opisuje wzór $I(y) = I_0 \cos^2\left(\frac{\pi d y}{\lambda L}\right)$, jednoznacznie pozwalając na pierwszy w historii pomiar długości świetlnej fali $\lambda$.

\textbf{Interferometr Michelsona:}
Podczas gdy Young wykorzystał rozszczepienie czoła fali, w 1881 roku A.A. Michelson skonstruował interferometr bazujący na koncepcji \textbf{rozszczepienia amplitudy}. Światło (spójność czasowa) kierowane jest na półprzepuszczalne zwierciadło płytkowe pod kątem $45^\circ$, gdzie połowa fali ulega odbiciu w ramię pionowe, a połowa przechodzi w ramię poziome. Na końcach ramion zwierciadła odbijają fotony z powrotem, gdzie spotykają się one i ponownie łączą. Umożliwia to badanie z nieprawdopodobną wprost precyzją przesunięć fazowych wywołanych ułamkowymi (na poziomie nanometrów) zmianami długości dróg między dwoma prostopadłymi ramionami. To właśnie tym urządzeniem udowodniono doświadczalnie (w słynnym eksperymencie Michelsona-Morleya) całkowity brak tak zwanego "wiatru eteru", kładąc podwaliny, na których Einstein oparł Szczególną Teorię Względności.

\section{Zjawisko Dopplera dla fal mechanicznych i elektromagnetycznych; fala uderzeniowa.}

\textbf{Zjawisko Dopplera (fale mechaniczne, np. dźwięk):}
Polega ono na rejestrowaniu przez obserwatora innej częstotliwości (i długości) fali niż ta, którą faktycznie generuje źródło, co jest spowodowane ruchem źródła bądź obserwatora względem sprężystego ośrodka propagacji.
Dla fal mechanicznych kluczowy jest fakt istnienia wyróżnionego układu odniesienia – samego ośrodka (np. powietrza). Oznacza to, że matematycznie i fizycznie należy odróżnić ruch źródła od ruchu obserwatora:
\begin{itemize}
    \item \textbf{Źródło porusza się (obserwator spoczywa):} Zmienia się obiektywna długość fali uformowanej w ośrodku (grzbiety w kierunku ruchu nakładają się "gęściej"). Nowa długość fali $\lambda' = \lambda - v_z T$, co daje rejestrowaną częstotliwość $f' = f \frac{v}{v \pm v_z}$ ($v$ to prędkość fali w ośrodku, minus dla zbliżania).
    \item \textbf{Obserwator porusza się (źródło spoczywa):} Długość fali w ośrodku jest niezmienna ($\lambda = v/f$), ale obserwator wyłapuje szybciej mijające go grzbiety, zatem zmienia się relatywna prędkość czoła fali uderzającej w detektor: $v_{wzg} = v \pm v_o$. Wtedy $f' = \frac{v_{wzg}}{\lambda} = f \frac{v \pm v_o}{v}$ (plus dla zbliżania).
\end{itemize}
Uogólniony wzór klasyczny: $f' = f \left(\frac{v \pm v_o}{v \mp v_z}\right)$. Górne znaki dotyczą zbliżania. Zjawisko manifestuje się w życiu codziennym m.in. drastycznym spadkiem wysokości syreny pogotowia w momencie mijania słuchacza na ulicy.

\textbf{Zjawisko Dopplera dla fal elektromagnetycznych:}
Fale elektromagnetyczne w próżni nie wymagają ośrodka. Zgodnie z postulatami Einsteina prędkość światła $c$ jest identyczna dla każdego obserwatora, zatem próżnia jako ośrodek nie istnieje – odróżnienie ruchu "obserwatora" od "źródła" nie ma sensu fizycznego. Wyprowadzenie wzoru bazuje na transformacji Lorentza (dylatacja czasu powoduje, że okres zegara w tykającym, poruszającym się źródle wydłuża się z punktu widzenia obserwatora). Relatywistyczny wzór zależy tylko od jednej \textbf{prędkości względnej $v$}:
\[
f' = f \sqrt{\frac{1 \pm v/c}{1 \mp v/c}}
\]
(plus w liczniku dla zbliżania się układów). Zjawisko to ma kolosalne zastosowanie w astronomii; m.in. pozwala na pomiar radialnych prędkości oddalania się galaktyk na podstawie systematycznego przesuwania się ich linii widmowych w stronę fal dłuższych, czyli tak zwanego przesunięcia ku czerwieni (z ang. \textit{redshift}), będącego koronnym dowodem na inflację i ewolucję Wszechświata (prawo Hubble'a), a także jest fundamentem działania radarów do pomiaru prędkości pojazdów drogowych.

\textbf{Fala uderzeniowa i stożek Macha:}
To szczególne, wysoce gwałtowne zjawisko mechaniczne ma miejsce, gdy prędkość źródła fali w ośrodku ($v_z$) przewyższa prędkość rozchodzenia się fal w tymże ośrodku ($v$), czyli $v_z > v$ (liczba Macha $M > 1$). W tej sytuacji źródło "wyprzedza" wytwarzane przez siebie fale. Poszczególne fronty fal kulistych nieustannie nakładają się na siebie w fazie wzdłuż wspólnej powierzchni stycznej tworzącej potężne zbicie materii. Powstaje tzw. \textbf{stożek Macha}. Kąt rozwarcia $\alpha$ tego stożka (liczony od osi toru ruchu) jest tym ostrzejszy, im szybszy jest obiekt:
\[
\sin\alpha = \frac{vt}{v_z t} = \frac{v}{v_z} = \frac{1}{M}
\]
Przejściu stożka fali uderzeniowej nad obserwatorem towarzyszy skokowy, ekstremalny wzrost ciśnienia i gęstości rzędu milionowych części sekundy – objawia się to ogromnym, nierzadko destrukcyjnym dla szyb w oknach i błon bębenkowych, grzmotem powszechnie określanym jako grom dźwiękowy przelatującego samolotu naddźwiękowego. Optycznym analogiem tego w elektromagnetyzmie, gdy cząstka naładowana porusza się w medium (np. reaktorze z wodą) z szybkością przewyższającą tamtejszą fazową prędkość światła $v = c/n$, jest promieniowanie Czerenkowa – dające charakterystyczną chłodno błękitną poświatę rdzenia jądrowego.

\section{Przegląd widma fal elektromagnetycznych; najważniejsze cechy poszczególnych rodzajów promieniowania.}

Z punktu widzenia elektrodynamiki klasycznej Maxwella, zjonizowana atmosfera Słońca, cewka radiowa w smartfonie oraz klisza rentgenowska operują zjawiskami fizycznie tego samego rodzaju. Promieniowanie elektromagnetyczne to rozchodzące się na wskroś przestrzeni falowe, nawzajem prostopadłe oscylacje sprzężonego pola elektrycznego i magnetycznego, poruszające się w próżni ze stałą prędkością $c$. Poszczególne zakresy różnią się makroskopowo jedynie długością fali ($\lambda$) i wynikającą z relacji Plancka-Einsteina energią pojedynczego fotonu ($E = hf = hc/\lambda$). Klasyczny podział widma w kierunku malejącej długości fali i rosnącej niszczącej energii (od niskich do wysokich częstotliwości) prezentuje się następująco:

\textbf{1. Fale radiowe ($\lambda > 10^{-1}$ m):}
Posiadają najniższą częstość i potrafią osiągać długości fal rzędu tysięcy kilometrów. Powstają w wyniku przepływu prądu szybkozmiennego w antenach radiowych (obwodach LC) według opisu zjawisk drgań ładunku. Fale radiowe nie są pochłaniane przez przeszkody niemetalowe ani chmury, dzięki czemu ich ranga komunikacyjna zrewolucjonizowała świat i przoduje w telefonii, radiofonii, sygnałach GPS oraz Wi-Fi. Krótkie fale radiowe odbijają się od jonosfery ziemskiej, potrafiąc objąć swym zasięgiem całą kulę ziemską (radio amatorskie).

\textbf{2. Mikrofale ($10^{-3}$ m do $10^{-1}$ m):}
Są generowane w specjalistycznych lampach próżniowych jak klistrony czy magnetrony (w oparciu o silne wirowanie i hamowanie elektronów w magnetronie). Mają krytyczne zastosowania radarowe (bo odbijają się rykoszetem od dużych, fizycznych ciał metalowych bez wielkiego rozpraszania, tzw. echolokacja radarowa). Mikrofale w okolicach 2.45 GHz rezonują (wywołują rotacyjne stopnie swobody) stricte z dipolowymi cząsteczkami węgla i wody (co jest podwaliną potężnego ogrzewania komórek wewnątrz mikrofalówek gastronomicznych). Są też istotą Kosmicznego Promieniowania Tła ($T = 2.7$ K) pozostałego z termodynamiki po Wielkim Wybuchu.

\textbf{3. Podczerwień (IR) ($700$ nm do $1$ mm):}
Emitowana przez absolutnie każde ciało posiadające temperaturę większą od zera bezwzględnego (jest to po prostu manifestacja promieniowania cieplnego ciała doskonale czarnego wynikająca z równań Plancka i Wiena). Źródłem są termiczne drgania molekularne i cząsteczkowe na małą odległość. Stosowana w termowizji, laserach cięcia materiałowego, spektroskopii molekularnej oraz do transmisji bliskiej (piloty RTV, dawniej IrDA).

\textbf{4. Światło widzialne (VIS) ($400$ nm do $700$ nm):}
Wąski wycinek całego gigantycznego tortu elektromagnetycznego. Promieniowanie emitowane przede wszystkim przez skoki elektronów powłokowych (walencyjnych) podczas relaksacji oraz przejść optycznych w rozgrzanych zjonizowanych atomach metali. Jest jedynym promieniowaniem rejestrowalnym na receptory czopków i pręcików siatkówki ludzkiego oka (największa czułość optymalizowana ewolucyjnie przypada na pasmo 550 nm, zielone piki profilu Słońca G2V). Przenosi akurat taką energię (rzędu eV), która jest zdolna dysocjować chemicznie powiązania węglowodorowe bez szkodliwej i natychmiastowej jonizacji jądrowej – uruchamia proces fotosyntezy na Ziemi.

\textbf{5. Ultrafiolet (UV) ($10$ nm do $400$ nm):}
Od tego pasma rozpoczynają się zagrożenia radiologiczne. Fotony pasma UV są na tyle energetyczne, by móc z łatwością przenikać na płytką glębokość skórną, wywoływać defekty poparzeniowe i bezpośrednio modyfikować, mutować molekuły spiralnego kwasu RNA, uderzając wiązania między parami zasad. Posiadają potężne zdolności fotojonizacyjne na zewnątrz. Słońce uderza tym w Ziemię obficie, z czego ponad 97\% chłonie wysoko na pułapie pasmo jonizującego ozonu O3. Ultrafioletem powszechnie sterylizuje się chemię organiczną (zabijanie bakterii szpitalnych, utwardzanie polimerowe plomb stomatologicznych pod światłem).

\textbf{6. Promieniowanie rentgenowskie (X) ($0.01$ nm do $10$ nm):}
Wytwarzane podczas potężnego wyhamowywania elektronów rzucanych napięciem kilkudziesięciu kilowoltów w cel ciężki (tzw. zjawisko Bremsstrahlung na anodzie z wolframu - promieniowanie hamowania) oraz jako przeskoki najgłębszych, uwięzionych K-powłokowych nośników (promieniowanie charakterystyczne tarczy). Jest wysoko przenikliwe i bez trudu przechodzi przez zjonizowaną materię miękką tkanek, jednak za sprawą zjawiska ugięcia i absorpcji fotoelektrycznej blokuje się na cięższych atomach potasu i wapnia budujących układ kostny. Otworzyło to nową epokę w krystalografii strukturalnej ciał stałych.

\textbf{7. Promieniowanie gamma ($\gamma$) ($< 0.01$ nm):}
Pochodzenie to stanowi koronę ekstremów energii fotonów ($E > 100$ keV na pojedynczy foton!). Fale gamma nie są w żadnym ujęciu tworzone przez elektrony ani powłoki, pochodzą ze strikte spontanicznych, gwałtownych przegrupowań cząstek nuklearnych ukrytych w jądrach niestabilnych izotopów ciężkich (np. Uranu czy Radu), podczas anihilacji par (elektron-pozyton), rozpadów obcych i reakcji termojądrowych w supernowych. Mają niesamowitą, głęboką przenikliwość przez materię (przenika przez ołów i żelbet). Mają najbardziej niszczące działanie dla wszelkiego życia biologicznego na planecie (choroba popromienna). Wąskie dawki skupione bywają zbawienne jako nuklearny skalpel potrafiący utrącać selektywnie głęboko osadzone komórki rakowe, gdzie operacja fizyczna chirurgiczna jest ryzykowna – w radio-onkologii.

\section{Dyfrakcja wiązki świetlnej; zasada Huygensa-Fresnela, dyfrakcja Fraunhofera na szczelinie, siatka dyfrakcyjna.}

\textbf{Zasada Huygensa-Fresnela:}
Zjawisko dyfrakcji (ugięcia fal) to odstępstwo od prostoliniowego rozchodzenia się fal świetlnych w pobliżu krawędzi nieprzezroczystych przeszkód. Światło "zagląda" w obszar cienia geometrycznego. Zjawisko to, podobnie jak interferencję, tłumaczy zasada Huygensa-Fresnela, będąca fundamentem optyki falowej.
Zasada ta głosi, że \textbf{każdy punkt ośrodka, do którego dociera czoło fali świetlnej, staje się źródłem nowej, kulistej fali elementarnej} (postulat Huygensa). Amplituda fali w dowolnym późniejszym punkcie przestrzeni jest wynikiem \textbf{interferencji} (superpozycji z uwzględnieniem różnic faz i amplitud) wszystkich tych fal elementarnych (uzupełnienie Fresnela).

\textbf{Dyfrakcja Fraunhofera na pojedynczej szczelinie:}
Dyfrakcja Fraunhofera to szczególny, graniczny przypadek dyfrakcji (zwany dyfrakcją w polu dalekim), w którym zarówno źródło światła, jak i ekran obserwacyjny znajdują się w nieskończoności (w praktyce osiąga się to stosując równoległą wiązkę światła z kolimatora i obserwując obraz przez soczewkę skupiającą).
Rozpatrzmy szczelinę o szerokości $a$, oświetloną prostopadle płaską, monochromatyczną falą o długości $\lambda$. Zgodnie z zasadą Huygensa, szerokość $a$ możemy traktować jako nieskończony ciąg punktowych źródeł fal. Fale te uginają się pod kątem $\theta$. Warunek na \textbf{minimum dyfrakcyjne} (całkowite wygaszenie) wynika z geometrycznego parowania fal z górnej i dolnej połowy szczeliny (różnica dróg wynosi $\lambda/2$). Wzór na minima dyfrakcyjne to:
\[
a \sin\theta_n = n\lambda \quad \text{dla } n = \pm 1, \pm 2, \dots \text{ (ale } n \neq 0\text{)}
\]
Pomiędzy minimami powstają jasne prążki (maksima wtórne), jednak ich natężenie gwałtownie maleje. Zdecydowana większość energii skupiona jest w potężnym prążku centralnym (dla $n=0$, $\theta=0$), który jest dwukrotnie szerszy od prążków bocznych.

\textbf{Siatka dyfrakcyjna:}
To przyrząd optyczny w postaci układu wielkiej liczby ($N$) identycznych, bardzo wąskich i równoległych do siebie szczelin (lub nacięć na szkle), oddalonych od siebie o stałą odległość zwaną \textbf{stałą siatki ($d$)}.
W przypadku siatki zachodzi potężna, wielokrotna interferencja fal z tysięcy szczelin. Aby nastąpiło \textbf{główne maksimum interferencyjne}, fale z dwóch \textit{sąsiednich} szczelin muszą wzmacniać się w fazie (różnica dróg optycznych musi być całkowitą wielokrotnością długości fali). Równanie siatki dyfrakcyjnej na maksima główne:
\[
d \sin\theta_n = n\lambda \quad \text{dla } n = 0, \pm 1, \pm 2, \dots
\]
Gdzie $n$ to rząd widma. Ponieważ szczelin jest $N \gg 1$, maksima te są niezwykle ostre i wąskie (pomiędzy nimi występuje $N-1$ minimów i $N-2$ bardzo słabych maksimów wtórnych). Siatka dyfrakcyjna jest fundamentalnym narzędziem \textbf{spektroskopii} – pozwala na niezwykle precyzyjne rozdzielenie fal o różnej długości $\lambda$ (każda barwa ugina się pod nieco innym kątem).

\section{Optyka geometryczna jako graniczny przypadek optyki falowej. Podstawowe przyrządy optyczne: lupa, luneta, mikroskop.}

\textbf{Optyka geometryczna jako granica optyki falowej:}
Rozchodzenie się światła ma z definicji naturę falową (opisywaną równaniami Maxwella). Optyka geometryczna to zbiór potężnych metod przybliżonych, polegających na całkowitym odrzuceniu pojęcia fali, jej fazy, interferencji i dyfrakcji, na rzecz \textbf{promienia świetlnego}. Promień ten to matematyczna linia (ortogonalna do powierzchni falowych), wzdłuż której transportowana jest energia w postaci światła. 
Przybliżenie optyki geometrycznej jest w pełni słuszne (i staje się asymptotycznie dokładne) wyłącznie w granicy, gdy \textbf{długość fali światła dąży do zera ($\lambda \to 0$)} lub mówiąc praktycznie – gdy długość fali światła jest mikroskopijnie mała w porównaniu z rozmiarami przyrządów optycznych (soczewek, źrenic, szczelin), przez które światło przechodzi ($a \gg \lambda$). Gdy przeszkody są rzędu $\lambda$, krzywoliniowe efekty dyfrakcyjne rozmywają pojęcie promienia.

\textbf{Podstawowe przyrządy optyczne:}
Przyrządy te służą poprawie zdolności rozdzielczej i powiększeniu kąta widzenia ludzkiego oka. Wykorzystują załamanie światła w soczewkach skupiających (wypukłych) i rozpraszających (wklęsłych). Opierają się na równaniu soczewki cienkiej: $\frac{1}{f} = \frac{1}{x} + \frac{1}{y}$, gdzie $f$ to ogniskowa, $x$ - odległość przedmiotu, $y$ - odległość obrazu.

\begin{itemize}
    \item \textbf{Lupa:} 
    Najprostszy przyrząd optyczny zbudowany z pojedynczej soczewki skupiającej o małej ogniskowej $f$. Przedmiot umieszcza się blisko soczewki, w odległości mniejszej niż ogniskowa ($x < f$). W efekcie po stronie obserwatora powstaje obraz pozorny, prosty i powiększony. Celem lupy jest wytworzenie tego obrazu w odległości dobrego widzenia dla oka ludzkiego (zwykle definiowanej jako $D \approx 25$ cm). Powiększenie kątowe lupy wyraża się wzorem:
    \[
    p = \frac{D}{f}
    \]
    (wyprowadzone przy założeniu powstawania obrazu w odległości dobrego widzenia, nie w nieskończoności). Typowe powiększenia lupy wynoszą od 2x do 10x.

    \item \textbf{Mikroskop:}
    Służy do ekstremalnego powiększania bardzo małych przedmiotów położonych w małej odległości. Zbudowany z dwóch układów optycznych – soczewki o wybitnie małej ogniskowej zwanej \textbf{obiektywem} ($f_{ob}$) oraz soczewki pełniącej rolę lupy zwanej \textbf{okularem} ($f_{ok}$). Przedmiot umieszcza się minimalnie za ogniskową obiektywu. Obiektyw tworzy w tubusie (o długości optycznej $\Delta$) powiększony obraz rzeczywisty. Obraz ten staje się z kolei przedmiotem dla okularu, przez który oko ogląda finalny obraz (pozorny, odwrócony). Całkowite powiększenie to iloczyn powiększeń obu elementów:
    \[
    p = p_{ob} \cdot p_{ok} \approx \left(\frac{\Delta}{f_{ob}}\right) \left(\frac{D}{f_{ok}}\right)
    \]

    \item \textbf{Luneta (Teleskop refrakcyjny):}
    Przyrząd służący do powiększania kąta widzenia przedmiotów położonych w nieskończoności (astronomicznych) lub bardzo daleko. Różnica w budowie względem mikroskopu polega na tym, że obiektyw lunety ma \textbf{bardzo dużą} ogniskową ($f_{ob}$). Światło z odległej gwiazdy (wiązka równoległa) wpada do obiektywu, który tworzy jej rzeczywisty obraz w swojej płaszczyźnie ogniskowej. Płaszczyzna ta pokrywa się z przednim ogniskiem okularu (krótkiego). Powiększenie kątowe lunety (tzw. lunety Keplera, odwracającej obraz) wynosi:
    \[
    p = \frac{f_{ob}}{f_{ok}}
    \]
\end{itemize}

\section{Oddziaływanie światła z ośrodkiem; prawa odbicia i załamania, całkowite wewnętrzne odbicie, dyspersja normalna i anomalna.}

Kiedy wiązka światła (fala elektromagnetyczna) pada na granicę dwóch ośrodków materialnych o różnej gęstości optycznej, następuje redystrybucja jej energii: część ulega odbiciu z powrotem do ośrodka pierwszego, część przechodzi (załamuje się) do ośrodka drugiego, a część jest absorbowana. Gęstość optyczną charakteryzuje bezwymiarowy współczynnik załamania światła $n = c/v$, gdzie $v$ to prędkość światła w danym ośrodku.

\textbf{Prawo Odbicia:}
Promień padający, promień odbity oraz normalna (prostopadła) do powierzchni granicznej wystawiona w punkcie padania, leżą w jednej płaszczyźnie. Dodatkowo:
\[
\text{Kąt padania } (\alpha) = \text{Kąt odbicia } (\alpha')
\]
Kąty te mierzymy zawsze od normalnej. Fala po odbiciu zachowuje swoją długość i prędkość.

\textbf{Prawo Załamania (Prawo Snella):}
Przechodząc do ośrodka, w którym światło porusza się wolniej (większe $n$), czoło fali "zwalnia", co geometrycznie wymusza złamanie (skrzywienie) toru promienia względem granicy ośrodków. Stosunek sinusa kąta padania ($\alpha$) do sinusa kąta załamania ($\beta$) jest równy stosunkowi prędkości światła w tych ośrodkach, czyli odwrotnemu stosunkowi ich współczynników załamania:
\[
\frac{\sin\alpha}{\sin\beta} = \frac{v_1}{v_2} = \frac{n_2}{n_1} \implies n_1 \sin\alpha = n_2 \sin\beta
\]
Fala załamana zmienia swoją długość fali ($\lambda = \lambda_0 / n$), ale jej częstotliwość ($f$) pozostaje bezwzględnie stała.

\textbf{Całkowite wewnętrzne odbicie (CWO):}
Zachodzi uwarunkowane przejściem światła z ośrodka optycznie gęstszego ($n_1$) do rzadszego ($n_2$), czyli gdy $n_1 > n_2$ (np. ze szkła do powietrza). W takim układzie kąt załamania jest większy niż kąt padania ($\beta > \alpha$). Zwiększając kąt padania, w pewnym momencie kąt załamania osiągnie wartość graniczną $90^\circ$ (promień "śliźnie" się po powierzchni). Kąt padania, dla którego tak się dzieje, nazywamy \textbf{kątem granicznym ($\alpha_{gr}$)}:
\[
\sin\alpha_{gr} = \frac{n_2}{n_1}
\]
Gdy fala padnie pod kątem $\alpha > \alpha_{gr}$, zjawisko załamania przestaje fizycznie istnieć. Cała (100\%) energia świetlna zostaje rykoszetem odbita w głąb ośrodka pierwszego. Zjawisko to przewyższa jakością najlepsze metaliczne lustra (które zawsze nieco absorbują energię) i jest fundamentem działania \textbf{światłowodów} dla celów globalnego Internetu.

\textbf{Dyspersja normalna i anomalna:}
Dyspersja to zjawisko polegające na zależności prędkości fali w ośrodku materialnym (a więc i współczynnika załamania $n$) od jej częstotliwości (długości fali $\lambda$). Z tego powodu pryzmat rozszczepia światło białe na widmo kolorów (każdy załamuje się pod nieco innym kątem).
\begin{itemize}
    \item \textbf{Dyspersja normalna:} zachodzi w obszarach przezroczystości dielektryka (daleko od jego częstości rezonansowych). Współczynnik załamania \textbf{maleje} wraz ze wzrostem długości fali ($\frac{dn}{d\lambda} < 0$). Oznacza to, że światło fioletowe (krótsze) jest w szkle załamywane silniej niż czerwone (dłuższe).
    \item \textbf{Dyspersja anomalna:} występuje w otoczeniu pasm silnej absorpcji (rezonansów wewnątrz atomów). W tym bardzo wąskim przedziale długości fal współczynnik załamania \textbf{rośnie} wraz ze wzrostem długości fali ($\frac{dn}{d\lambda} > 0$). Powoduje to odwrócenie kolejności barw w widmie i silne pochłanianie energii przez ośrodek.
\end{itemize}

\section{Polaryzacja fal: prawo Malusa, kąt Brewstera, dwójłomność, aktywność optyczna.}

Zjawisko polaryzacji jest bezwzględnym dowodem na to, że światło jest \textbf{falą poprzeczną} (wektory drgań $\vec{E}$ i $\vec{B}$ są prostopadłe do kierunku rozchodzenia się wiązki). Naturalne światło spójne (ze Słońca, żarówki) jest niespolaryzowane – emitowane drgania wektora $\vec{E}$ zachodzą we wszystkich możliwych płaszczyznach prostopadłych do promienia. Gdy wektor natężenia pola elektrycznego drga tylko w jednej ustalonej płaszczyźnie, falę nazywamy spolaryzowaną liniowo.

\textbf{Prawo Malusa (Polaryzatory):}
Światło naturalne można spolaryzować, przepuszczając je przez polaroid (np. długołańcuchowe polimery). Idealny polaryzator przepuści $I = \frac{1}{2}I_0$ natężenia naturalnego. Jeżeli to już spolaryzowane światło trafi na drugi polaryzator (analizator), którego oś przepuszczania jest skręcona o kąt $\theta$ względem pierwszego, to do analizatora przejdzie tylko rzut składowej pola na jego oś ($E = E_0 \cos\theta$). Ponieważ natężenie fali jest proporcjonalne do kwadratu amplitudy pola ($I \propto E^2$), otrzymujemy \textbf{prawo Malusa}:
\[
I = I_{max} \cos^2\theta
\]
Skrzyżowanie polaryzatorów ($\theta = 90^\circ$) prowadzi do całkowitego wygaszenia ($I = 0$).

\textbf{Kąt Brewstera (Polaryzacja przez odbicie):}
Kiedy niespolaryzowane światło pada na dielektryk (np. taflę wody), odbita fala staje się częściowo spolaryzowana (w płaszczyźnie równoległej do tafli). Szkocki fizyk David Brewster odkrył, że istnieje taki jeden charakterystyczny kąt padania, przy którym światło odbite jest \textbf{w pełni, 100\% spolaryzowane liniowo}. Następuje to w geometrycznym warunku, gdy promień odbity i promień załamany tworzą ze sobą kąt prosty ($90^\circ$). Z prawa Snella wynika wtedy równanie na ten specyficzny \textbf{kąt Brewstera ($\alpha_B$)}:
\[
\tan\alpha_B = \frac{n_2}{n_1}
\]
Zjawisko to jest powszechnie wykorzystywane w fotografii i okularach polaryzacyjnych, które "odcinają" oślepiające spolaryzowane refleksy np. od powierzchni mokrej szosy.

\textbf{Dwójłomność kryształów:}
Większość cieczy i amorficznych ciał stałych to materiały izotropowe (własności w każdym kierunku są takie same). Istnieją jednak kryształy anizotropowe (np. szpat islandzki – węglan wapnia), w których gęstość elektronowa zależy od osi sieci. Pojedynczy promień światła niespolaryzowanego wchodzący do kryształu rozszczepia się na dwa promienie: 
- \textbf{Zwyczajny} (spełniający prawo Snella, $n_o = const$),
- \textbf{Nadzwyczajny} (którego prędkość i współczynnik $n_e$ zależą od kierunku padania, co łamie klasyczne prawo załamania).
Oba wyjściowe promienie są całkowicie i wzajemnie prostopadle spolaryzowane. Przesunięcie fazowe między nimi jest podstawą budowy tzw. płytek ćwierćfalowych.

\textbf{Dwójłomność wymuszona:}
Izotropowe ciało (np. szkło, pleksi) poddane \textbf{naprężeniom mechanicznym} traci swoją symetrię i staje się w tym miejscu kryształem anizotropowym, wykazując dwójłomność. Przepuszczenie światła przez dany detal umieszczony między dwoma skrzyżowanymi polaryzatorami ujawnia piękne, kolorowe prążki interferencyjne dokładnie w miejscach kumulacji naprężeń (zjawisko elastooptyki). Podobny anizotropowy efekt (efekt Kerra) można wywołać, przykładając do cieczy (np. nitrobenzenu) bardzo silne pole elektryczne.

\textbf{Aktywność optyczna ciał (Skręcenie płaszczyzny polaryzacji):}
Niektóre substancje organiczne i nieorganiczne (tzw. ośrodki chiralne, pozbawione środka symetrii cząsteczek, np. wodny roztwór cukru, kwarc lewoskrętny i prawoskrętny) obracają płaszczyznę polaryzacji przechodzącego przez nie liniowo spolaryzowanego światła. Kąt skręcenia jest wprost proporcjonalny do grubości warstwy substancji oraz, w przypadku roztworów, do jej stężenia. Zjawisko polarymetrii jest gigantycznym wsparciem w chemii analitycznej (badanie czystości i stężeń cukrów).

\section{Naturalna szerokość linii widmowej - czas życia stanów wzbudzonych. Stany metatrwałe.}

Zgodnie z klasycznym modelem atomu Bohra, różnica energii fotonu wyemitowanego przy przeskoku elektronu między orbitą wyższą ($E_2$) a niższą ($E_1$) jest stała. Zatem długość emitowanej fali $\lambda = \frac{hc}{E_2 - E_1}$ powinna stanowić nieskończenie cienką kreskę w widmie spektroskopowym (funkcję delta Diraca). 

W rzeczywistości nawet w ideale zerowej temperatury i absolutnej próżni żadna linia widmowa nie jest nieskończenie cienka. Obserwowane w spektroskopach przejścia posiadają minimalne, niemożliwe do usunięcia "rozmycie" energetyczne zwane \textbf{naturalną szerokością linii widmowej}. Oprócz tego na szerokość linii wpływa temperatura (\textit{poszerzenie dopplerowskie}) oraz uderzenia atomów (\textit{poszerzenie zderzeniowe}), ale szerokość naturalna wynika bezpośrednio z najgłębszych praw mechaniki kwantowej.

\textbf{Zasada nieoznaczoności a szerokość linii:}
Powodem poszerzenia jest słynna \textbf{Zasada nieoznaczoności Heisenberga} dla czasu i energii:
\[
\Delta E \cdot \Delta t \ge \frac{\hbar}{2}
\]
Oznacza ona, że natura zabrania jednoczesnego z całkowitą dokładnością wyznaczenia energii stanu i czasu trwania tego stanu. Jeśli elektron pozostaje w stanie podstawowym ($E_1$) w nieskończoność ($\Delta t \to \infty$), to jego energia jest doskonale określona ($\Delta E \to 0$). Zjawisko to jednak dotyczy elektronów na powłokach wyższych.

\textbf{Czas życia stanów wzbudzonych ($\tau$):}
Elektron wzbudzony do wyższego stanu ($E_2$) nie przebywa tam wiecznie. Podlega spontanicznej emisji (deekscytacji). Akt rozpadu rządzony jest prawdopodobieństwem i jest nieprzewidywalny dla konkretnego atomu, jednak makroskopowo opisuje go \textbf{średni czas życia stanu wzbudzonego ($\tau$)}, będący odpowiednikiem pojęcia z rozpadów promieniotwórczych (zwykle dla przejść dozwolonych optycznie tzw. dipolowych czas ten jest rzędu zaledwie $\tau \sim 10^{-8}$ s).
Podstawiając go do zasady Heisenberga ($\Delta t \approx \tau$), uzyskujemy nieoznaczoność energii stanu wzbudzonego:
\[
\Delta E \approx \frac{\hbar}{\tau}
\]
Ponieważ energia stanu ulega rozmyciu w paśmie $\Delta E$, także wyemitowany foton nie będzie rygorystycznie "sztywny". Jego częstotliwość rozłoży się wokół częstotliwości centralnej w paśmie $\Delta \nu \approx \frac{1}{2\pi\tau}$. Rozkład natężenia tego profilu ma matematyczną formę krzywej Lorentza. Jest to właśnie naturalna szerokość linii widmowej, i wynosi ona zazwyczaj zaledwie ułamki nanometra.

\textbf{Stany metatrwałe:}
Powyższe fakty opisywały typowe stany wzbudzone, z których przejście jest dozwolone przez reguły wyboru kwantowo-mechanicznych liczb (wynikające z zachowania pędu i parzystości, np. zasada Laporte'a zakazująca zmian, gdy $\Delta l = 0$).
Istnieją jednak wysokie stany atomowe, dla których z powodów kwantowych powrót na niższy poziom optycznym powrotem dipolowym jest fizycznie zabroniony. Taki elektron utknął na górze. Oczywiście wróci, używając innych mechanizmów, uderzeń termicznych, lub powrotów wielopolowych (wzbronionych), jednak jest to potwornie nieprawdopodobne. 
Stany takie nazywamy \textbf{metatrwałymi}. Ich średni czas życia jest gigantycznie dłuższy od normalnego, może wynosić ułamki sekund, minuty (zamiast $10^{-8}$ s). Stany te mają kluczowe, rewolucyjne zastosowania:
\begin{itemize}
    \item \textbf{Fosforencja:} elektrony powoli wymykają się z metatrwałej "pułapki", powodując, że zegarki fluoryzujące po naświetleniu żarówką świecą w ciemności przez długie godziny.
    \item \textbf{Budowa LASERÓW:} wydłużony czas życia stanu pozwala optycznie "napompować" materię, uzyskując tak zwane obsadzenie inwersyjne (więcej elektronów wysoko niż nisko) konieczne do zapłonu zjawiska uwięzionej Akcji Laserowej.
\end{itemize}

\section{Doświadczenie Rutherforda - wpływ rezultatów na model budowy atomu.}

Na początku XX wieku fizycy znali już istnienie lekkich, ujemnych elektronów dzięki J.J. Thomsonowi. Wiedząc, że z racji stabilności atomy są obojętne, Thomson zaproponował w 1904 roku tzw. model "ciasta z rodzynkami". Twierdził w nim, że cała objętość atomu wypełniona jest "rozmazanym", jednorodnym ładunkiem dodatnim stanowiącym chmurę bazową, w której tkwią unieruchomione elektrony (jak rodzynki w rzadszym cieście). 

\textbf{Doświadczenie (Eksperyment Geigera-Marsdena, 1909):}
Ernest Rutherford, wraz ze swymi badaczami Hansem Geigerem i Ernestem Marsdenem, postanowił zweryfikować prawdziwość tego modelu za pomocą potężnych mechanicznych rykoszetów nuklearnych z cząstek alfa ($\alpha$ – ciężkie i dodatnio naładowane jądra helu powstające w rozpadach promieniotwórczych). Strumień niezwykle szybkich cząstek $\alpha$ skierowano z ołowianego kolimatora na bardzo cienką folię ze złota (grubość zaledwie ułamków mikrometra – około tysiąca warstw atomowych). Za folią, na okrągłym ekranie scyntylacyjnym rejestrowano pod mikroskopem mikrobłyski uderzających cząstek. 

Zgodnie z modelem Thomsona, cząstki $\alpha$ stanowiące masywny i bardzo gęsty taran miały całkowicie lekceważyć dodatnie "ciasto" poprzez jego małą gęstość i penetrować folię przelatując praktycznie po prostej. Maksymalne spodziewane kąty rozproszenia szacowano na zaledwie ułamki jednego stopnia.

\textbf{Rezultaty doświadczenia:}
Wyniki eksperymentu wstrząsnęły klasyczną mechaniką i fizyką statystyczną:
1) Istotnie, znakomita większość cząstek $\alpha$ przeleciała przez złotą blachę w idealnej linii prostej bez drgnięcia lub z drgnięciem odchylającym w milistopniach.
2) Niemniej jednak \textbf{co około 8000. cząstka doznawała zjawisk wysoce "niefizycznych" według modeli statystycznych: odchylała się o potężne kąty} ($> 90^\circ$), a czasami potrafiła wręcz uderzyć frontalnie i odbić się prosto w twarz badaczy, o kąt niemal $180^\circ$.
Sam Rutherford wspominał, że "było to równie niewiarygodne jak wystrzelenie pocisku 15-calowego z armaty w chusteczkę higieniczną, od której on by rykoszetem odbił w ciebie."

\textbf{Powstanie planetarnego modelu budowy atomu (1911):}
Fakt, że większość cząstek przeniknęła, oznaczał, iż atom składa się z wprost zatrważającej \textbf{pustki} przestrzennej.
Z faktu natomiast zaobserwowania tak dramatycznych rozproszeń tylnych (aby zawrócić niezwykle ciężką cząstkę alfa, należy oprzeć na niej potężną siłę elektrostatyczną z niesamowicie bliskiego uderzenia $F \propto 1/r^2$), wysnuto, iż obalony został model gładkiego ciasta. 

Rutherford, na kanwie eksperymentów, sformułował słynny \textbf{planetarny model atomu}:
\begin{itemize}
    \item Zdecydowana, gigantyczna część objętości (i rozmiaru optycznego $\sim 10^{-10}$ m) atomu to bezwzględna przestrzeń próżna wypełniona jedynie przez chmurę krążących po stacjonarnych i kołowych orbitach elektronów o nikłej wadze.
    \item Cały dodatni ładunek atomu oraz uderzające 99,95\% jego pełnej masy, nie jest rozsmarowane po całej objętości, lecz upchane zostało do potwornie gęstego, punktowego wręcz \textbf{jądra atomowego} rezydującego precyzyjnie w centrum. Rozmiary tego jądra (wynikające z rzutowania równań rozpraszających fizyki kulombowskiej) szacowano na obiekty o promieniach niewiele więcej niż gigantyczne skale rzędu femtometra ($10^{-15}$ m) – a więc 100 000 razy mniej objętościowe w rozmiarze od wymiarów obłoków świetlnych atomu!
\end{itemize}
Mimo bezapelacyjnego tryumfu na polu rozpraszania, Model Rutherforda zawierał jedną drastyczną słabość z perspektywy klasycznej elektrodynamiki: elektron jako ładunek krążący po łuku z przyspieszeniem dośrodkowym musi bez ustanku promieniować falę elektro-magnetyczną i drastycznie tracąc energię w czasie ok. $10^{-8}$ sekund zawalić się spiralnie wprost w kulombowskie ramiona nuklearnego jądra. Brak wyjaśnienia faktu stabilności materii na ziemi stał się punktem wyjścia, od którego wkrótce Niels Bohr zapoczątkował model i potęgę Mechaniki Kwantowej w nauce.

\section{Widmo promieniowania rentgenowskiego; ogólna charakterystyka, promieniowanie charakterystyczne i hamowania.}

Promieniowanie rentgenowskie (odkryte przez Wilhelma Röntgena w 1895 r.) to wysokoenergetyczne fale elektromagnetyczne o długościach z zakresu od ok. $0.01$ nm do $10$ nm, powstające wskutek uderzania strumieniem szybkich elektronów w metalową tarczę (anodę) w rurze próżniowej. Widmo emitowanego w ten sposób promieniowania składa się z dwóch niezależnie powstających, nakładających się na siebie części: widma ciągłego (promieniowania hamowania) oraz dyskretnego (promieniowania charakterystycznego).

\textbf{Promieniowanie hamowania (Bremsstrahlung) - widmo ciągłe:}
Mechanizm ten ma naturę czysto klasyczną elektrodynamiczną (choć z istotną granicą kwantową). Padające z katody rozpędzone elektrony (o energii kinetycznej $E_k = eU$, gdzie $U$ to napięcie przyspieszające lampy) wnikają głęboko w tarczę z ciężkiego metalu (np. wolframu, molibdenu). Zbliżając się do masywnych, wysoce naładowanych jąder atomowych, elektrony są potężnie odchylane (zakrzywiają tor lotu) i gwałtownie hamowane w silnym polu kulombowskim jądra.
Zgodnie z elektrodynamiką klasyczną, każdy ładunek poruszający się z przyspieszeniem (lub opóźnieniem) emituje falę elektromagnetyczną. Ponieważ akty hamowania są w każdym przypadku inne i losowe (elektron może oddać dowolną część swej energii na wyemitowanie fotonu), powstaje całe \textbf{widmo ciągłe} fal, na którym zarejestrowane jest tło (garb) natężenia.
Widmo ciągłe posiada jednak \textbf{ostrą granicę krótkofalową ($\lambda_{min}$)}. Wynika ona bezpośrednio z faktu kwantowego – największą energię będzie miał foton wytworzony w pojedynczym zdarzeniu absolutnie całkowitego wyhamowania elektronu "do zera". Wtedy cała jego energia $eU$ idzie w jeden foton ($hc/\lambda_{min}$). Stąd słynny wzór Duane'a-Hunta na granicę obcięcia widma ciągłego:
\[
\lambda_{min} = \frac{hc}{eU}
\]

\textbf{Promieniowanie charakterystyczne - widmo liniowe:}
Drugi rodzaj zjawisk pojawia się, gdy napięcie przyspieszające w lampie przekroczy próg krytyczny dla danego materiału tarczy. Mechanizm tutaj uwarunkowany jest całkowicie wewnętrzną, kwantową budową atomów metalu i generuje prążki nakładające się w postaci ostrych i potężnych pików intensywności (np. serie $K_\alpha, K_\beta$) na tło ciągłe.
Jeżeli uderzający z zewnątrz niezwykle szybki elektron wpadnie z odpowiednim pędem głęboko do wnętrza atomu w anodzie i trafi tam (wybije) inny elektron mocno uwięziony na jednej z najgłębszych powłok (np. najniższej powłoce K, czyli głównej liczbie kwantowej $n=1$), pozostawi w tarczy po sobie "puste miejsce" energetyczne – powłokę z dziurą. Stan ten jest bardzo niestabilny. Niemal natychmiast elektron rezydujący wyżej (z powłoki L, M itd.) z impetem "spada" w dół, na ową powłokę K, zatykając dziurę.
Różnica energii obu precyzyjnie określonych przez orbitale kwantowe powłok jest emitowana na zewnątrz w postaci pojedynczego fotonu X o \textbf{dyskretnej (ściśle określonej)} dla danego atomu długości. Henry Moseley odkrył, że częstotliwości promieniowania w danej serii charakterystycznej ($\nu$) rosną proporcjonalnie do kwadratu ładunku jądra atomowego ($Z$) wg prawa Moseleya ($\sqrt{\nu} \propto (Z-\sigma)$). Fakt ten pozwala po dziś dzień używać analizatorów widma rentgenowskiego jako niepodważalnej, bezbłędnej wizytówki (identyfikatora) składu minerałów.

\section{Wzbudzanie i deekscytacja jader i atomowych powłok elektronowych. Zjawisko Mössbauera.}

Wzbudzenia atomów lub ich potężnych, małych centrów nuklearnych pociągają za sobą zjawiska deekscytacji na wiele dróg sprzęgających elektrony powłokowe i samo jądro.

\textbf{Emisja Kwantów Gamma ($\gamma$) i rentgenowskich (X):}
Jądro atomowe (podobnie jak powłoka elektronowa) jest układem kwantowym o rygorystycznie dyskretnych dozwolonych poziomach energii. Po procesie np. rozpadu radioaktywnego beta, powołane do życia jądro córka zazwyczaj znajduje się w wysokim stanie wzbudzonym. Powrót nukleonów na głęboki stan podstawowy pociąga za sobą wyemitowanie na zewnątrz różnicy ogromnej energii w postaci potężnego fotonu nazywanego \textbf{kwantem gamma ($\gamma$)} o świetlnej i destrukcyjnej wręcz sile MeV (megaelektronowoltów). Kwanty rentgenowskie o niższych, na poziomie kilkudziesięciu keV, rzędach, generowane są przy przegrupowaniach z głębokich powłok K, jak to ujęto wcześniej.

\textbf{Elektrony konwersji wewnętrznej:}
To bez-fotonowy, potężny mechanizm alternatywny dla emisji kwantu $\gamma$. Wzbudzone jądro nie emituje fotonu, ale przekazuje całą swoją nadmiarową energię bezpośrednio jednemu ze swoich własnych, bardzo głęboko krążących elektronów (najczęściej z powłoki K - które ze względu na mechanikę falową potrafią "odwiedzać" bliskie im jądro fizycznie). Elektron ten, za sprawą gigantycznego strzału energetycznego zostaje wyrwany z układu (całkowicie z atomu) w przestrzeń. Elektron ów unosi pęd oraz ma energię kinetyczną ściśle określoną widmowo jako ($E_{kin} = \Delta E_{jadra} - E_{wiazania}$). Zjawisko konwersji zostawia jednak za sobą silnie zjonizowany na głębokościach atom.

\textbf{Elektrony Augera:}
Gdy głęboka powłoka zostanie pusta (wskutek uderzenia rentgenowskiego, bądź podanej przed sekundą \textit{konwersji wewnętrznej}), elektron z wyższej orbity zrzuca się tam, aby zająć lukę. Emitowana przy tym energia najczęściej ulatuje jako foton charakterystyczny X. \textbf{Efekt Augera (Pierre Auger, 1925)} jest z kolei potężnym, alternatywnym bez-fotonowym aktem rozładowania atomu. Energia wyzwolona z przeskoku górnego elektronu zostaje pochłonięta od razu bezpośrednio przez zupełnie \textbf{inny (trzeci) elektron} na obrzeżach. Nabiera on impetu na opuszczenie całego układu, jako dyskretny i bezfotonowy strzał, z określoną w widmach siłą (znów, co do ułamków, zależną od gatunku chemicznego jako sygnatura uwięzienia powłok z węgla, tlenu).

\textbf{Zjawisko Mössbauera (1958):}
Szalenie precyzyjne wręcz, unikalne i fundamentalne doświadczenie w rygorystycznej mechanice kwantowej. W wolnych atomach emisja oraz absorpcja promieniowania tak gwałtownego jak kwanty gamma cierpi z powodu \textbf{odrzutu} (zgodnego z klasycznym zasobem zachowania pędu jądra przy wystrzale z ogromną silą relatywistycznego kwantu), podobnie jak działa strzelba. Energia fotonu zostaje pomniejszona o energię odrzutu jądra atomu. Gdy taki słabszy foton dojdzie do bliźniaczego atomu obok, brakuje mu energii by rozbudzić dokładnie rezonans – a szpice jądrowe rezonansowe mają w tym przypadku wręcz śmiertelną ostrość z rzędu $10^{-13}$ w stosunku do mocy z powodu długiego życia metatrwałego stanu jądrowego (Ir 191).
Rudolf Mössbauer odkrył, że uwięziwszy owe jądra rygorystycznie bardzo sztywno w chłodnej i twardej matrycy (sieci) krystalicznej i sprowadzając do niskich temperatur by zatrzymać szumy statystyczne, moment gwałtownego odrzutu ułamkowego uderzenia jądra pobierany jest zbiorczo (masywnie) przez miliony jąder wokół naraz złączonych ze spoiwem chemicznym. Ponieważ ciężar całego tak zaatakowanego kryształu (M) dąży do niebotycznej wartości względem samego elektronu, ubytek pędu i straty rozbicia drgań zanika (efektywnie do perfekcyjnego, czystego ZERA). Rezultatem stała się fizyka bezodrzutowej absorpcji gamma z fantastyczną rozdzielczością dla weryfikowania przesunięć czasoprzestrzennych z Ogólnej Względności.

\section{Jądro atomowe, jego odkrycie, budowa, własności.}

\textbf{Odkrycie Jądra Atomowego:}
Dokonało się to za sprawą historycznego doświadczenia \textbf{Geigera-Marsdena (w roku 1909)} analizowanego w laboratorium i matematycznie okiełznanego i upublicznionego dwa lata później na wskroś autorytetem samego ich szefa – \textbf{Ernesta Rutherforda (1911)}. Miarowe i regularne "odbicia piłeczek tenisowych" (cząstek $\alpha$) wystrzeliwanych strumieniem z polonu wprost o złotą i potwornie cienką taflę mikrometrową odrzucano na bok jako niemożliwe dla teorii śliwkowego puddingu J.J Thomsona. Tylko upchanie absolutnie całej pozytywnie ładującej się gęstości masowej atomu u samej jego podstawy (średnicy zaledwie rzędu femtometru) uratowało równania elektrostatycznego powrotnego odepchnięcia krzywych torów rozproszeń dla $\alpha$.

\textbf{Budowa jądra (Nukleony):}
Jądro to miniaturowe (o potwornych i niewyobrażalnych w świecie makroskopowym gęstościach rzędu $\sim 10^{17}$ kg/m$^3$) drobiny zlokalizowane pośrodku atomów, składające się z cząstek oddziałujących \textbf{silnie} zwanych łącznie \textbf{nukleonami}. Nukleony formują potężną nuklearną bryłę zbudowaną z dwóch fermionów o potężnej, zbliżonej na skali ziaren masie:
\begin{itemize}
    \item \textbf{Protony (p):} Cząstki obarczone ładunkiem dodatnim ($+1e$). Liczba protonów oznaczona historycznie \textbf{Z} decyduje dożywotnio i chemicznie o tożsamości izotopowej w Układzie Mendelejewa (jest równoznaczna z nazwą gazu szlachetnego czy węgla czy rtęci) i rygorystycznie paruje powłoki elektronowe.
    \item \textbf{Neutrony (n):} Odnalezione późno z rąk ucznia samego Rutherforda, \textbf{Jamesa Chadwicka (1932)}, w toku precyzyjnych i wybitnie inteligentnie ułożonych zderzeń parafiny pod ciosami uderzeń promieni z radu. Są fundamentalnie niewrażliwe na pole Coulomba i nieładują elektrycznie (0e). Liczbę neutronów w ułamkach podaje się jako $N$.
\end{itemize}
Całość budulców masy i gęstości spina rygorystyczna \textbf{Liczba Masowa A ($A= Z+N$)}. Izotop z racji budowy (na ogół ułamek węgla ${}^{12}_{6}C$ czy niestabilny węgiel patologiczny ${}^{14}_{6}C$) zmienia objętość (liczbę $N$), zatajając dla chemii resztę zjawisk, z kolei fizycznie modyfikując jądro i spin diametralnie.

\textbf{Własności fizyczne jąder atomowych (Oddziaływania i Energia):}
Zgromadzenie potwornej puli silnie odpychających się dodatnich protonów z ładunkami na bardzo ograniczonym blisko rozmiarze $10^{-15}$ m musi rygorystycznie prowadzić (poprzez prawa Maxwella $F = kq^2/r^2$) do nieobliczalnej nuklearnej eksplozji Coulomba jądra już podczas pierwszych milisekund Wszechświata (jądra nie mogłyby trwać we wspólnocie i przyrodzie same ze sobą jako twory).
Fakt oparcia stabilności zawdzięcza istnieniu w Naturze zjawiska nieuchwytnego na miarę skali ziemskiej – czyli tzw. \textbf{Oddziaływań Silnych} nuklearnego gniazda. Występuje tu pole sił przyciągających się wzajemnie neutronów z protonami i odwrotnie, ale wysoce \textbf{krótkozasięgowych} na długość nie dalej i nie szerzej niż zaledwie około 2 do 3 fm (femtometrów). Gdy zbliżyć nukleony pod promień rzędu $1 fm$, siły "Sklejenia Silnego" miażdżąco tryumfują siłą potęgi nad powierzchownym i elektrostatycznym rozerwaniem w obłoki (około stokrotnie!), stąd ciężkie jądra radzą sobie tak idealnie, i dlatego u cięższych transuranów jądrowych trzeba coraz gęściej nasycać proporcje neutronowej mazi buforowej ponad proporcje samych protonów (aż na ułamek 1.5 dla stabilnego Bizmutu czy Ołowiu!).

Owa praca i tytaniczny pot ulepienia Nukleonów z silnego rwania to tzw. potęga \textbf{Defektu Masy i Energii Wiązania ($E = \Delta m c^2$)} jako spadek ułamka całości w relacji o masę mniejszej przy złączeniu względem odważenia atomów "solo". Najstabilniejszym pikiem oddziaływań upakowania we wszechświecie po chłodnych gwiazdach jest stabilne wiązanie ${}^{56}Fe$ uwięzione 8 MeV na każdy pojedynczy sklejony nukleon żelaza.

\section{Modele jądra atomowego – kroplowy, powłokowy.}

Budowa jądra jest tak dalece oddalona kwantowo i silnie zagęszczona nieliniowo z dynamiką obłoku gluonowego oraz zderzeń pędowych, że nie potrafiono nigdy wyprodukować matematycznej "Zunifikowanej" potęgi praw, objaśniającej idealnie jądro w zwięzłej czystej teorii bytu równań całkujących, co ma miejsce chociaż w prozaicznym elektro-gęstym atomie. Dlatego w fizyce od lat ratowano spójne dane rygorystycznymi, niezależnie dobieranymi fenomenologicznymi pod kątem i profilem zjawisk \textbf{Modelami}, uogólnionymi na podstawie prób statystycznych. Dwa najważniejsze uwarunkowane zjawiskowo profile to kropla wodna o dużej gęstości oraz układy kwantowo-powłokowe do ujęć rotacji izotopowej i magicznej magii fermionów spinowych.

\textbf{Model Kroplowy Jądra (N. Bohr i C.F. von Weizsäcker, 1935):}
Skupiony rygorystycznie wokół oddania ujęcia \textbf{makroskopowego (kolektywnego)}. Zauważono bezbłędnie fenomenologię rygorystycznych danych uśredniania zjawisk silnych pędów jądrowych. Podobnież i fizycznie spójnie jak w skali u kropli mikroskopowej mazi cieczy (rzucanej i kołyszącej się wody z racji gęstości bezwzględnej):
\begin{itemize}
    \item Jądra różnych potęg i różnych pierwiastków (jak krople dużej miary morskiej i mniejszej) nie ściskają swoich pętli gęstniejąc! \textbf{Gęstość jądrowej materii jest u absolutnie wszystkich jąder bezkompromisowo STAŁA} - jądro uranu ma potężniejszą objętość niż węgiel czy ołów rygorystycznie pod wielokrotność izotopu A ($V \propto A \implies R = R_0 \sqrt[3]{A}$). Zbliżenie zjawisk i materii jest w uścisku absolutnie ślepe i zwarte "od spodu".
    \item Występuje tam oddziaływanie rygorystycznie nasyceniowe spoiw "do najbliższych rąk u boku", toteż wewnętrzny proton łączy uścisk mocno z kolegą do promienia rzędu około 2 fm, totalnie i z góry lekceważąc obcy proton czy neutron rezydujący gdzieś po drugiej ułamkowej połówce. Zatem potęga jądra nie mnoży się pod silnymi rygorami, lecz w sumie rośnie do izotopu.
\end{itemize}
Owocem i wybitnym matematycznie ujęciem wziętym z opisu "Kropli Nukleonowej wody" ujętym m.in napięciem objętościowym i powierzchniowym był podany genialny empiryczny tzw. \textbf{Wzór Weizsäckera}. Służył do niezwykle rygorystycznej (w ogólnym ogonie gładkiego łuku danych) fenomenologii obliczeń całkowitej "Energii Wiązania", tłumacząc u zarania gwałtowne potęgowanie ubytków energii spoiwa jako dowód gwałtownych mechanicznych skoków na ewolucyjne cięcie się jądra u procesów ROZSZCZEPIANIA się "kropli na dwie małe kropelki odkształcając asymetrycznie jądro w gruszkowatą bańkę nuklearną" co miało i zaowocowało projektem reaktorowym o potężnej wydajności (U-235).

\textbf{Model Powłokowy (M. Goeppert-Mayer, J.H.D. Jensen, 1949):}
Gdzie Model kroplowy jest głuchy na gładkim grzbiecie (jako gładka linia danych o rozpadzie materii i spajaniu w potężną bańkę jądrową) fizycy uświadczyli jednak wybitne zęby na wykreślnym wykresie potężnej energii! Otóż statystyka i stabilność niektórych absolutnie silnie skrojonych nuklearnie izotopów rwała się wybitnie silnie, uciekając gładkim formom Weizsäckera, zachowując zjawiskową i bezsprzeczną formę i moc przy potwornych naporach i zderzeniach izotopowych. Te ukoronowane izotopy z potężnymi wiązaniami odznaczono specjalną gładką \textbf{liczbą MAGICZNĄ} protonów i neutronów (np. 2, 8, 20, 28, 50, 82, 126 potężnie trzymały twarde nukleony np. Hel, Tlen-16, czy Żelazo-56).
Postawiono bezapelacyjny i genialny kwantowy \textbf{Model powłokowy rygorystyczny}. Jądro nuklearne z racji praw Pauli'ego traktowane tak i w ten sposób musi wybudować sobie wewnętrzną (podobną elektronowej ale oddziaływującym twardym silnym gniazdem jądrowym spin-orbita ujęciem u spin-orbit-coupling) głęboką własną "Nuklearną Powłokę (orbitale jądrowe)". Protony jak i Neutrony wpadając potężnie, obsadzają dolne gniazda energii kwantowo i po brzeg, zamykając skorupy szczelnie w stabilność rygorystyczną, twardą, i magiczną powłokę bez reaktywności wyrywającej izotop u błędu na brzeg. Model z uderzeniem bezwzględnym rozwiązał mity wybitnego ułożenia spinu oraz anomalie magii w obronie stabilnych rzędów ciał układu ziemi, kładąc za ujęcie sprzężeń spin-orbita w roku 1963 nagrodę Nobla.

\section{Spontaniczne przemiany jądrowe – rodzaje, teoria.}

Część z pośród rygorystycznych wielkich ciał naładowanego izotopowo atomu pozostawiona samoistnie i luźno nie utrzyma we wspólnocie potęgi swoich jądrowych układów przy niebilansowej puli budulca sił neutronów do protonów w potężnych ciałach Coulomba. Nukleony chcące zrzucić i podważyć wysoki spiętrzony bilans, osuwają (transformują spontanicznie u praw rygorów probabilistyki jądrowej z czasem zaniku) na mniejszy zysk i niższą stabilną formę energetyczną u stóp "koryta energii powłok i mas", emitując spontaniczne i dyskretne zawirowania nazywane ogółem promieniotwórczym procesem potęg ciał (promieniotwórczość Marii Skłodowskiej). Zjawiska probabilistycznego statystycznego rozpuku nuklearnego podyktowano jako uniwersalne absolutne zasady zachowania praw \textbf{ładunku, masy, energii i potężnego pędu rygorystycznego w osiach jądrowych}:

\textbf{1. Przemiana ALFA ($\alpha$):}
Zjawisko dręczące potężnie jądra wybujałe rozmiarowo powyżej magicznych wielkich ołowianych złożeń. U tak dużego obwodu izotopu nukleon w jądrze mocując słabe krawędzie rozrywa spoiwa uciekając ze zwartą paczką jako jądro Helu-4 (tzw. zwarty na skraju i magicznie uformowany, potężny zamek $\alpha$ - o ładunku $+2e$, masie $4u$). Całkowity rygor zrzuca Jądro-córkę wybitnie o 2 oczka do tyłu z ładunku, staczając izotop $A$ w dół potężnie o mocną ułamkową formę $-4$ w dół na tabeli! Całe zjawisko pokonuje bezbłędny potężny mur energii wokół potężnego ładunku ułamkowego u krawędzi "za wysokiego progu" jako bezwzględny czysty probabilistyczny \textbf{kwantowy cud Efektu Tunelowego}, uwięziona i stłamszona paczka przenika powolnie z mroku muru stawiając potęgę prawdopodobieństw u Gamowa (twierdzenie tunelowania formuje regułę Geiger-Nuttall dla mocy Czasu-T i zjawisk wybijającej wielkiej energii Alfa!).
\[ {}^{A}_{Z}X \rightarrow {}^{A-4}_{Z-2}Y + {}^{4}_{2}He \]

\textbf{2. Przemiany i patologie z rodziny BETA ($\beta$):}
Rejestrujemy zjawisko rygorystycznie i całkowicie głuche w procesie na układy Sił Głębokich Silnych jądra u potężnych mas, jako iż za mechanizm szarpiący rozbiciem atomu na części ułamu neutronów za odpowiadają (ukryte fundamentalnie u cząstek pędu głębiej) \textbf{Oddziaływania fundamentalne Słabe (bozony wektorowe $W$ i $Z$ na ułamki rzędu od kwarków w jądro!)}. Mechanizmy regulują niefortunny przeładowany pęd izobary u rygoru krzywej, utrzymując stałe jądro pod mase całkowitą ($A=const$!) - ładunek skacze o jądro u góry i do tyłu rzutując nowym układem atomu o rygor mniejszy lub więcej protonu z pikiem o $Z \pm 1$ bez straty zrębów!:
\begin{itemize}
    \item \textbf{Rozpad $\beta^-$:} W izotopach "przebogatych" z nadmiarem luźnych stłoczonych zjawisk i nadmiaru neutronu rygor kwantowy rozłupuje wewnątrz i fundamentalnie owy obcy neutron i przeobraża na wspaniały gładki proton (jako wbudowany atomowy pęd stabilności jądra) przy tym dla oddania pędu u wyrzutu wypluwając rygorystycznie niesłychanie szybki Elektron i drobinę bez masy, ułamkowe pędowe rygorystyczne - \textbf{Antyneutrino elektronowe ($\bar{\nu_e}$)}. Rozpad izotopu C-14, nienawiść rygorystyczna od biologicznych cząstek na rzecz stabilności obcego powłokowego azotu i archeologia.
    \item \textbf{Rozpad $\beta^+$:} Złamanie na potęgi jąder w pule potężnie za dużego ładunku rygorystycznego "protonowego" nad masą atomu bez ustanku zmusza ujemny wyrzut dodatniej w mroku puli atomowej potęgi jako ulepienie (rozerwany atom pozytonów z kwarkami powłok p) i powstaje z protonu w środku "Słabego" i potężny rozbity nowy "neutron" a z mroku wykwita, potężnie do strzału antycząstka (dodatni kwant antymaterii) z pędem nazywana u pędu potęg jako rygor \textbf{Pozyton (dodatni elektron)} dla ucieczki a dla równowagi sił z potęg na rygor dodany pęd rzuca luźnym pędem głuche dla praw atomów zderzeń jądro to słynne - \textbf{Neutrino rygorystyczne elektronowe ($\nu_e$)}.
    \item \textbf{Wychwyt E-elektronowy powłoki K (EC):} U jądra tak dużych i wielkich, z wielkim niedoborem i z zjawiskiem braków neutronów atom połyka potężną gęstość probabilistyczną krążącego blisko (orbital $s$) małego elektronu z powłoki rygorystycznie-najbliższej \textbf{K} neutralizując protonowy zapas jądra by połączyć do zera układy bez zrzutu z potęgi! 
\end{itemize}

\textbf{3. Przemiana rygorystycznych pętli GAMMA ($\gamma$):}
Wybitny brak naruszeń masy i tożsamości ładunków jądrowych i pierwiastka mendelejewskiego. Izotop (Y*) z patologią i rygorystycznie za mocno pobudzonym (napompowanym u po wypadku beta rozbitej formacji pędem lub spacji powłoki) podrzuca nukleony staczając o moc stabilnego ucisku wypluwając różnice energetycznej potęgi bez ciał rygorystycznie fotonem kwantowym (100 keV - Mev!) rygorystycznego ułamka.

\section{Oddziaływanie promieniowania jądrowego z materią.}

Nuklearne odpady ciał (korpuskularne rozproszone cząstki alfa z rozpadu oraz elektronowa szarża z wściekłych rozpadów izotopów z cząstką beta, bądź twarda spójność energii elektromagnetycznej w postaci morderczych rygorystycznych fotonów ciał Gamma i X), kiedy to wkroczą przebiegiem w miękką chmurę lub ołowianą gąszcz zjawisk - oddziałują absolutnie diametralnie. W fizyce i zjawisku osłon dozymetrycznych klasyfikuje się na ułamek "Ciężkie ładunki mocne z polem, Lekkie ładunki elektronów i wibrujące mrokiem na zjawiskach elektromagnetycznych z neutralnym tłem fotonów promieni rygorów" w różnorodnych proporcjach hamowni praw fizycznych gęstości w materii atomów ośrodka!

\textbf{1. Promienie masywne korpuskularne wielkich i ciężkich ułamków ALFA ($\alpha$):}
Cząstka ta z rzędu jądra (ładunek dwukrotnie rygorystycznie wielki u Coulomba i moc +2e u helu i potężna masa kinetyczna u uderzenia 5 MeV u atomów). Przemieszcza i rozpycha się u gęstwy atomowej wybitnie gwałtownie u praw. Ze względu na rygor potężny i ładunek przykleja z impetem z sił do gąszczu obcych elektronów mazi ośrodka, miażdżąc wiązania organiczne i uderzając potężną jonizacją w chmurach niszcząc u rygorystycznej paczki, setki tysięcy wiązań par. Jej zasięg - z racji grubiańskiego pędu upakowanej straszliwej mocy niszczeń Coulomba, traci uderzeniowy rygorystyczny pęd momentalnie w drodze zaledwie milimetrowym (w szkle lub papierze rzędu gładkiej potęgi kartki czy obumarłego zrogowaciałego naskórka) by ukraść na rygor z atomów o obłok gęsty elektrony do "Helu" i zgasnąć atomowo i pokojowo!

\textbf{2. Lekkie ładunki rozpędów z ciał z elektronów BETA ($\beta^-$ i $\beta^+$):}
Ponad 7-tysiąc razy od ciał mniejszych od u uderzenia Alfa pędowym i lżejsza u rygorystycznie (ładunkiem u potęg tylko jednej potęgi). Droga jądrowa "Lekka i sprytna elektronowa gąszcz" pokonuje znacznie, rygorystycznie bardziej chaotyczną drastycznie i niekoherencyjną potęgę torów i wielokrotnych odskoków u elastycznych kulombowskich rygorach odbić o chmury i obłoki (nie traci wielkim u ułamkiem energii rygoru w mgnieniu do rygoru potęg, zasięgi bywają w milimetrach wody i obok do centymetrów ciała - blacha milimetrowa ze stopów lekkich aluminium!). Dodatkowym czynnikiem do uderzenia beta z energią obok lekkich mas hamujących jest gigantyczny efekt powstawania u potęg na cięższym ośrodku widma \textbf{Hamowania rygorystycznego na obłoku u uderzenia rygoru ciężkiego Coulomba promieni - Promieniowania hamowania (X)!} Cząstki dodatniego elektronu $\beta^+$ (jako anty-materii) na pętli o potęgi gęste i obłoki traca rygorystycznie do ciał elektronów atomowych by rozbijać do spacji o rozłup i rygor u unicestwień dwóch na promieniach wybitnie ostrych promieni - i wybuchu jako dwie cząstki kwantów anihilacji!

\textbf{3. Fotony u wysokoenergetycznych bez mas z rygoru kwantowych fal GAMMA ($\gamma$ i rentgena):}
Promieniowanie mrocznej siły wybitnie fotonowej fali wnika jako potęga bez masowa nie na skutkach oddziaływań ciał i kulombowskich przymiotów tarczy powłok zderzeń w oddal, lecz dyskretnie gasi gęstwy i łupie się w mgnieniu na potężnych u zjawisk atomu kwantowego tarczy w 3 głównych aktach rygorystycznego ułamka!
\begin{itemize}
    \item \textbf{1. Zjawisko rygorystycznie fotoelektryczne:} (dominacja energii niskiego fotonu z rzędu gładkiego keV- 0.5 MeV u potęg dla ułamka gęstwy jądra Z cięższego) Kwant $\gamma$ spotyka elektron u ciał w gnieździe atomu, znika fizycznie za nic u mroku (absorpcja totalna ciał fotonu z promieni atomów), wyrywając elektron jako pęd na obłoku i pędy energii kinetycznej elektronu z łupkiem atomowym na rozwałkę ciał. 
    \item \textbf{2. Zjawisko u obłokach efektów Comptona (rozproszenie klasyczne fal z kwantowym uderzeniem!):} (U szarż potęg o energiach średniej ułamkowej MeV potęg atomowych w chmurach 1 MeV u węgla czy ciała): Rygorystyczny Foton wybitnie wysoko i niszczycielsko odbija po skraju elastycznym u tarczy małego kwantowego obłoku u uderzenia u mroku na elektron "w swobodnym mroku rzutowym atomów" - łupiąc potężnie foton od ułamka elektronu z rozbiciem atomu. Elektron na kule rzutowo wypada (i w kłębek obłoku powłok), a potężnie poturbowany rygor z rzędu energii Fotonu po strzale i zderzeniu, zmienia rygorystycznie kierunek pętli kwantowych pędu i zmniejsza energie fotonową bez uśmiercenia w procesie o dłuższą od powłok u w mgnienia ułamkach fale potęg! (Wzór na w mroku rozbity kwant Comptona - powielenie rozbicia fotonów do osłabień w bloku i osłonie gęstej potęg)!
    \item \textbf{3. Produkcja mrocznych Par w Polach (kreacja elektron-pozyton rygorów):} (Powyżej niszczycielskiego i progu masy relatywistycznej 1,02 MeV i dominująca potęgami potwornych u energii o 5 MeV obłokach w tarczach Ołowianych): Rygorystyczny wnikający foton z rzędu pod powłokami spotyka u jądra straszliwe twarde potężne mroki sił jądrowych i Kulombowskich szpilek ciał - z mroku fotonu powstaje z niczego załamanie rygoru energii w ułamkowe Dwie pełnoprawne relatywistyczne potęgi gęstwy kwantowej ciał - Elektron (-e) i Pozyton (+e) porywając na masy u łupku.
\end{itemize}

\section{Rozszczepienie jąder atomowych, energetyka jądrowa.}

Zjawisko \textbf{rozszczepienia ciał rygorystycznych w ułamku z Jąder ciężkich i wybitnie chwiejnych (U-235, Pu-239, U-233)} ma u zarania fizyki fundamenty spiętrzonej potęgi niestabilnej u zarania ciał jako przesycone u ułamka Coulomba ciał u modeli o kropli jądrowej z energią upakowań z modelów od węgla na 8 MeV a potężnych z rzędu $7.6$ MeV w ołowiu na głębokich gniazdach. Zjawisko obnażyli 1938 Otto Hahn i Fritz Strassmann wraz w rygorze analiz u Lise Meitner w rygorystycznie promiennym zjawisku mroku naświetlania uranu z uderzeniem luźnym na Neutron o mrocznej z wolnych termicznych o pędzie ok. gładkiej formy kinetycznej w ułamek rzędu dziesiątych (0.025 eV).

\textbf{Mechanizm wielkiego mrocznego łupnięcia rygorystycznego powłok jądra atomu:}
Złowiony, bez żadnego czołowego wielkiego ładunku Coulomba Neutron mroczny termiczny, obija cichym splotem wielką barierę kropli jąder z izotopu Uranu i wnosi na rzecz jądra do mroku i ułamku ok. 6,5 MeV u pędu wybitnie energii "wiązania ukrytego potęg" do całego kotła u atomów kropli jako uran wzbudzonego mroku izotopu ${}^{236}_{92}U$ rygorystycznego na skraj. Na potężnych u zderzeń pętli od ułamka modelu z u kropli Weizsäckera, z silnie rygorystycznej z okrągłych sfer kropli - formuje się rozbudzona podłużna potężna bania i hantle z wielkich, potwornie przeładowanych z ładunku dodatniego z mroku ciał Coulomba, obcych i dwóch równych fragmentów.
Przy promieniach i wielkim rozstępie hantli (za sprawą kwadratu ułamków ładunku i z potęg na rozryw u rygoru 2 kul nukleonowych) potężne spoiwo i od mroku przyciąganie nuklearne z krótkiego ramienia rygorystycznych oddziaływań ciał - Silnych mroku załamuje się do sił z oddaleniem gładkich jąder w mroki rozpadu i z przewagą wygrywa ładunkowa potężna "eksplozja potęgi jądrowej na Coulomba"!

Wynikiem i rygorem u procesów ułamka jest z reguły (w rozpadach wielokrotnych i stochastycznych potęgi w izotop, wybitnych dla Ba, Kr lub innych rygorystycznych podłupanych wielkich z mroku u atomów):
\begin{itemize}
    \item 1) Dwa nowo-odrodzone z wielkim i lżejszym z potęg spoiw u atomów pędem rygorystycznie Jądra promieniotwórcze o mocnych ciałach do po mroku rozpadowych (z mroczną różną izobaryczną potęgą masy).
    \item 2) Na dodatek dla gęstego rygoru pędu u uranu rygorystycznego nad pętli od ułamka neutronów na rygor rzutu wyskakuje natychmiast mrokiem szybkim od ok. \textbf{2 do 3 ciał luźnych zrzuconych dla potęg w pęd ułamkowy Neutronów z pędu atomów!!} 
    \item 3) Ze spraw o defekcie masy u rozrzutu ciał luźnych w sumach rygoru rzutowych wyzwala się potężnie i makabrycznie gigantyczna u pędu kinetycznym odrzutowa energia ciał około 200 MeV ciał kwantowych mrokiem na zjawiska atomu dla spacji (w uderzeniach z węglowych procesów spalanych o tlen to 4 eV mrokiem! Stosunek spalania z węgla to u pędu około ułamka na miliony rzędów wielkości!).
\end{itemize}

\textbf{Reakcja łańcuchowa u energetyki rygoru atomowego i pętli:}
Obecność rygorystycznych nadmiarów 2-3 szybkich z rzutu ułamkowych z mroku neutronów (przy spacji od uderzenia jednego od zarania u mroku!) to otwarcie dróg do reaktora wielkich energii o bilansie dodatnim z łańcuchowej gęstwy rozpędzonych rzutowo w klatce potęgi powieleniowych $k>1$ w mroku.
Aby wywołać wybitnie u powieleniu łańcuch potęgi w reaktorze, trzeba rygorystycznie i powieleniem u potęg mroku zachować moc ułamka krytycznego bilansowań $k=1$. U pędu z uranu te narzucone wysokie na 2 MeV ciał neutrony porywu po mroku nie rzutują z zjawiska dla potęg atomowych z pędu powrotnych na U-235 (są głuche - brak na nich u ułamków izotopu wchłonień z ciał, o małym w pętli u przekroju mroku na procesie, w pętli do U-238 pęd uwięzi). \textbf{W reaktorach termicznych węglem chwytowych rygorystycznie i potęgą używa się zrębów Moderatorów wielkiej ciężkiej masy gęstwy i grafitu w pętli!} (np. woda pędem izotopu ciężkiej, mrok powłok węglowych grafitu u powielenia zderzeń!) celem do pętli do omijającego absorpcję w spacji rzutowej dla uranowej gęstwy na potężne wstrzymanie w spacji prędkich ułamków potęg neutronu rzutowego pędem w mroku w obłoki pędowych pętli (do wolnych termicznych potęg 0.025eV dla uranu), utrzymując dla pętli wybuchy w kłębie mrocznego w klatce mroku rdzenia w ryzach i regulując u ciał od powłoki przez zatapiane kadmu, potężne na z rzędu na pręty dla wychwytu "w gładkich potęg gniazd!" i bez gładkich bomb nuklearnych.

\section{Funkcja Lagrange'a i równania Lagrange'a na przykładzie oscylatora harmonicznego. Rozwiązania równania Lagrange'a.}

\textbf{Funkcja Lagrange'a (Lagranżjan):}
Mechanika analityczna w ujęciu Lagrange'a odchodzi od operowania siłami wektorowymi (jak w dynamice Newtona) na rzecz wielkości skalarnych: energii kinetycznej ($T$) i energii potencjalnej ($V$). Podstawową wielkością opisującą układ mechaniczny jest funkcja Lagrange'a, oznaczana jako $L$. Dla klasycznych układów zachowawczych (gdzie siły wynikają z potencjału niezależnego od prędkości) funkcja ta jest zdefiniowana jako różnica między całkowitą energią kinetyczną a całkowitą energią potencjalną układu:
\[
L(q_i, \dot{q}_i, t) = T - V
\]
gdzie $q_i$ to współrzędne uogólnione, $\dot{q}_i$ to prędkości uogólnione, a $t$ to czas. Zgodnie z zasadą najmniejszego działania Hamiltona, rzeczywisty ruch układu odbywa się w taki sposób, aby całka z funkcji Lagrange'a po czasie (tzw. działanie $S$) osiągała ekstremum.

\textbf{Równania Lagrange'a drugiego rodzaju:}
Z wariacji działania wyprowadza się fundamentalne równania ruchu, zwane równaniami Eulera-Lagrange'a. Dla układu o $s$ stopniach swobody układ ten tworzy zbiór $s$ równań różniczkowych zwyczajnych drugiego rzędu:
\[
\frac{d}{dt} \left( \frac{\partial L}{\partial \dot{q}_i} \right) - \frac{\partial L}{\partial q_i} = 0 \quad \text{dla } i = 1, 2, \dots, s
\]

\textbf{Przykład: Oscylator harmoniczny:}
Rozważmy jednowymiarowy, klasyczny oscylator harmoniczny – punkt materialny o masie $m$ poruszający się wzdłuż osi $x$, poddany działaniu siły sprężystej $F = -kx$. 
1. Współrzędną uogólnioną jest zwykłe położenie: $q_1 = x$.
2. Energia kinetyczna: $T = \frac{1}{2}m\dot{x}^2$.
3. Energia potencjalna: $V = \frac{1}{2}kx^2$.
4. Konstruujemy Lagranżjan:
\[
L(x, \dot{x}) = \frac{1}{2}m\dot{x}^2 - \frac{1}{2}kx^2
\]

Aby zapisać równanie ruchu, obliczamy poszczególne pochodne z równania Lagrange'a:
- Pochodna po prędkości (pęd uogólniony): $\frac{\partial L}{\partial \dot{x}} = m\dot{x}$
- Pochodna czasowa pędu uogólnionego: $\frac{d}{dt}(m\dot{x}) = m\ddot{x}$
- Pochodna po położeniu (tzw. siła uogólniona): $\frac{\partial L}{\partial x} = -kx$

Podstawiając do ogólnego równania Eulera-Lagrange'a:
\[
m\ddot{x} - (-kx) = 0 \implies m\ddot{x} + kx = 0
\]
Podzielenie przez $m$ prowadzi do powszechnie znanego równania różniczkowego oscylatora: $\ddot{x} + \omega_0^2 x = 0$, gdzie $\omega_0 = \sqrt{k/m}$ to częstość kołowa drgań własnych.

\textbf{Rozwiązania równania Lagrange'a:}
Rozwiązaniem powyższego równania różniczkowego liniowego drugiego rzędu o stałych współczynnikach jest funkcja harmoniczna. Zależność położenia od czasu wyraża się jako:
\[
x(t) = A \cos(\omega_0 t + \phi) \quad \text{lub alternatywnie} \quad x(t) = C_1 \cos(\omega_0 t) + C_2 \sin(\omega_0 t)
\]
gdzie stałe $A$ (amplituda) oraz $\phi$ (faza początkowa) lub stałe $C_1, C_2$ wyznacza się z warunków początkowych położenia $x(0)$ oraz prędkości $\dot{x}(0)$. Rozwiązanie to idealnie obrazuje fizykę drgań nietłumionych wokół położenia równowagi.

\section{Stopnie swobody, uogólnione położenia i prędkości na przykładzie swobodnego ruchu punktu materialnego po powierzchni kuli.}

\textbf{Stopnie swobody i więzy:}
Liczba stopni swobody ($f$) układu mechanicznego to najmniejsza liczba niezależnych współrzędnych potrzebnych do jednoznacznego określenia położenia wszystkich elementów układu w przestrzeni w dowolnej chwili. Dla pojedynczego punktu materialnego swobodnie poruszającego się w przestrzeni trójwymiarowej $f = 3$. 
Jeżeli na układ nałożone są więzy (geometryczne ograniczenia ruchu), liczba stopni swobody maleje. Jeżeli punkt jest zmuszony do poruszania się wyłącznie po powierzchni sfery (kuli) o stałym promieniu $R$, narzucone jest jedno równanie więzów holonomicznych: $x^2 + y^2 + z^2 = R^2$. Wobec tego liczba stopni swobody wynosi $f = 3 - 1 = 2$.

\textbf{Współrzędne i prędkości uogólnione:}
Zamiast używać kartezjańskich współrzędnych $x, y, z$ (które są ze sobą powiązane więzami), w formalizmie Lagrange'a wprowadza się \textbf{położenia uogólnione} $q_i$. Są to dowolne niezależne zmienne liczbowe (kąty, odległości), których liczba dokładnie równa się liczbie stopni swobody układu. 
Dla ruchu po sferze najwygodniej przyjąć klasyczne kąty ze współrzędnych sferycznych:
- $q_1 = \theta$ (kąt polarny, szerokość mierzona od bieguna, $0 \le \theta \le \pi$)
- $q_2 = \phi$ (kąt azymutalny, długość mierzona w płaszczyźnie równika, $0 \le \phi < 2\pi$)
Pochodne czasowe tych wielkości, czyli $\dot{\theta}$ oraz $\dot{\phi}$, nazywamy \textbf{prędkościami uogólnionymi}.

\textbf{Konstrukcja Lagranżjanu:}
Ponieważ ruch jest "swobodny" na powierzchni kuli (nie działa żadna zewnętrzna siła potencjalna np. od grawitacji, z wyjątkiem sił reakcji więzów, które ujęcie analityczne automatycznie eliminuje), energia potencjalna $V = 0$.
Funkcja Lagrange'a sprowadza się wyłącznie do energii kinetycznej. Wartość wektora prędkości w układzie sferycznym dla stałego promienia $R$ ($v_r = 0$) wyraża się jako $v^2 = v_\theta^2 + v_\phi^2 = (R\dot{\theta})^2 + (R\sin\theta\dot{\phi})^2$. Zatem Lagranżjan wynosi:
\[
L = T = \frac{1}{2}m R^2 \left(\dot{\theta}^2 + \sin^2\theta \dot{\phi}^2\right)
\]

\textbf{Równania Lagrange'a:}
Mamy dwa stopnie swobody, więc zapisujemy dwa równania ruchu.
1. Dla zmiennej azymutalnej $\phi$:
\[
\frac{\partial L}{\partial \phi} = 0 \quad \text{(ponieważ $\phi$ nie występuje w Lagranżjanie, tzw. zmienna cykliczna)}
\]
\[
\frac{\partial L}{\partial \dot{\phi}} = m R^2 \sin^2\theta \dot{\phi}
\]
Równanie Lagrange'a daje od razu $\frac{d}{dt}(m R^2 \sin^2\theta \dot{\phi}) = 0$, co prowadzi do całki ruchu $m R^2 \sin^2\theta \dot{\phi} = const$. Jest to zasada zachowania rzutu momentu pędu na oś Z ($L_z$).

2. Dla zmiennej polarnej $\theta$:
\[
\frac{\partial L}{\partial \theta} = m R^2 \sin\theta \cos\theta \dot{\phi}^2
\]
\[
\frac{\partial L}{\partial \dot{\theta}} = m R^2 \dot{\theta} \implies \frac{d}{dt}\left(\frac{\partial L}{\partial \dot{\theta}}\right) = m R^2 \ddot{\theta}
\]
Ostateczne równanie dla $\theta$ po skróceniu przez $mR^2$ to:
\[
\ddot{\theta} - \sin\theta \cos\theta \dot{\phi}^2 = 0
\]
Uzyskane równania precyzyjnie opisują tor punktu, który (jak można wykazać całkując te równania) jest z punktu widzenia geometrii odcinkiem wielkiego koła na sferze – tak zwanej linii geodezyjnej.

\section{Funkcja Hamiltona i równania Hamiltona na przykładzie oscylatora harmonicznego. Rozwiązania równań.}

Podczas gdy formalizm Lagrange'a używa jako zmiennych położeń uogólnionych $q_i$ i prędkości uogólnionych $\dot{q}_i$, mechanika Hamiltona opiera się na tzw. przestrzeni fazowej. Zmiennymi niezależnymi stają się w niej położenia uogólnione $q_i$ oraz \textbf{pędy uogólnione} $p_i$, zdefiniowane jako $p_i = \frac{\partial L}{\partial \dot{q}_i}$. Transformacja (zwana transformacją Legendre'a) prowadzi od Lagranżjanu do nowej kluczowej funkcji.

\textbf{Funkcja Hamiltona (Hamiltonian):}
Z definicji, funkcja Hamiltona opisana w zmiennych $(q, p, t)$ wyraża się wzorem:
\[
H(q_i, p_i, t) = \sum_i p_i \dot{q}_i - L(q_i, \dot{q}_i, t)
\]
W układach stacjonarnych (gdzie $L$ nie zależy jawnie od czasu, a więzy są skleronomiczne), Hamiltonian fizycznie odpowiada \textbf{całkowitej energii mechanicznej układu}: $H = T + V$. Należy jednak pamiętać, by koniecznie wyrazić prędkości $\dot{q}_i$ poprzez pędy $p_i$.

\textbf{Kanononiczne Równania Hamiltona:}
Ewolucja czasowa układu zadana jest przez zbiór $2s$ (gdzie $s$ to liczba stopni swobody) różniczkowych równań pierwszego rzędu, wykazujących piękną asymetrię znaku:
\[
\dot{q}_i = \frac{\partial H}{\partial p_i}, \quad \dot{p}_i = -\frac{\partial H}{\partial q_i}
\]

\textbf{Przykład: Oscylator harmoniczny:}
Rozważmy jednowymiarowy punkt na sprężynie: $T = \frac{1}{2}m\dot{x}^2$, $V = \frac{1}{2}kx^2$. 
Współrzędna uogólniona to $x$. Pęd uogólniony (tu tożsamy ze zwykłym pędem klasycznym) wynosi $p = m\dot{x}$, z czego wynika odwrotna relacja $\dot{x} = \frac{p}{m}$.
Przeliczmy energię kinetyczną na zmienną pędową: $T = \frac{p^2}{2m}$. Hamiltonian, będący sumą energii, przyjmuje postać:
\[
H(x, p) = \frac{p^2}{2m} + \frac{1}{2}kx^2
\]

Zapiszmy teraz kanoniczne równania ruchu:
1. Równanie na prędkość: 
\[
\dot{x} = \frac{\partial H}{\partial p} = \frac{p}{m}
\]
2. Równanie na tempo zmiany pędu (siłę):
\[
\dot{p} = -\frac{\partial H}{\partial x} = -kx
\]

\textbf{Rozwiązanie równań Hamiltona:}
Mamy układ dwóch sprzężonych równań różniczkowych pierwszego rzędu. Najłatwiej rozwiązać go, obliczając pochodną po czasie pierwszego równania: $\ddot{x} = \frac{\dot{p}}{m}$. Następnie z drugiego równania wstawiamy w to miejsce $\dot{p} = -kx$. Otrzymujemy:
\[
\ddot{x} = \frac{-kx}{m} \implies \ddot{x} + \omega_0^2 x = 0
\]
gdzie $\omega_0 = \sqrt{k/m}$. Widzimy więc pełną zgodność z modelem klasycznym. 
Zaletą ujęcia Hamiltona jest jednak geometryczny wgląd w \textbf{przestrzeń fazową} $(x, p)$. Rozwiązania tych równań:
\[
x(t) = A \cos(\omega_0 t + \phi_0), \quad p(t) = -m\omega_0 A \sin(\omega_0 t + \phi_0)
\]
Parametryczne wykreślenie ewolucji oscylatora w dwuwymiarowej przestrzeni $(x, p)$ z pominięciem parametru czasu ($t$) daje równanie elipsy:
\[
\frac{x^2}{2E/k} + \frac{p^2}{2mE} = 1
\]
Trajektoria fazowa oscylatora nietłumionego jest zawsze gładką, zamkniętą elipsą (lub okręgiem), co geometrycznie świadczy o okresowości i zachowaniu całkowitej energii układu (powierzchnia objęta elipsą jest całką działania $J = \oint p \, dx$).

\section{Położenia i pędy uogólnione na przykładzie wahadła matematycznego – równania Hamiltona; przybliżenie małych drgań.}

Wahadło matematyczne to model wyidealizowany: punkt materialny o masie $m$ zawieszony na nieważkiej, nierozciągliwej nici o stałej długości $l$, wychylony o kąt $\theta$ z pionu pod wpływem pola grawitacyjnego Ziemi ($g$). Ruch odbywa się na płaszczyźnie, po wycinku łuku okręgu.

\textbf{Współrzędne i pędy uogólnione:}
Z powodu więzów (stała odległość od punktu zawieszenia), układ ma zaledwie jeden stopień swobody. Najbardziej naturalnym położeniem uogólnionym jest kąt wychylenia $q = \theta$.
Pochodna kąta w czasie to prędkość kątowa $\dot{\theta}$. Prędkość liniowa masy to $v = l\dot{\theta}$.
Energia kinetyczna: $T = \frac{1}{2}mv^2 = \frac{1}{2}ml^2\dot{\theta}^2$.
Energia potencjalna (za zero przyjmując punkt zawieszenia): $V = -mgl\cos\theta$.
(Można przyjąć zero w najniższym punkcie, wtedy $V = mgl(1-\cos\theta)$; wybór wpływa tylko na stałą wartość Hamiltonianu, nie na same równania ruchu).
Lagranżjan układu to: $L = \frac{1}{2}ml^2\dot{\theta}^2 + mgl\cos\theta$.

Pęd uogólniony (zwany w tym przypadku momentem pędu koniugowanym z kątem $\theta$) definiujemy jako:
\[
p_\theta = \frac{\partial L}{\partial \dot{\theta}} = ml^2\dot{\theta}
\]
Z tej zależności wyznaczamy prędkość uogólnioną: $\dot{\theta} = \frac{p_\theta}{ml^2}$.

\textbf{Konstrukcja funkcji Hamiltona:}
Ponieważ układ jest stacjonarny i zachowawczy, Hamiltonian rygorystycznie równa się energii całkowitej ($H = T + V$), lecz wyrażonej ściśle jako funkcja położeń i pędów uogólnionych:
\[
T = \frac{1}{2}ml^2 \left(\frac{p_\theta}{ml^2}\right)^2 = \frac{p_\theta^2}{2ml^2}
\]
Zatem funkcja Hamiltona przyjmuje postać:
\[
H(\theta, p_\theta) = \frac{p_\theta^2}{2ml^2} - mgl\cos\theta
\]

\textbf{Równania kanoniczne Hamiltona:}
Korzystając z równań Hamiltona wyznaczamy ewolucję czasową przestrzeni fazowej ($\theta, p_\theta$):
1. Równanie prędkości (w tym przypadku prędkości kątowej):
\[
\dot{\theta} = \frac{\partial H}{\partial p_\theta} = \frac{p_\theta}{ml^2}
\]
2. Równanie dynamiki (zmiany pędu uogólnionego, tutaj oznaczające moment siły):
\[
\dot{p}_\theta = -\frac{\partial H}{\partial \theta} = -mgl\sin\theta
\]
Otrzymany układ równań różniczkowych jest precyzyjny (nieliniowy z powodu funkcji sinus), jednak na ogół trudny do analitycznego rozwiązania bez użycia całek eliptycznych.

\textbf{Przybliżenie małych drgań:}
Jeżeli założymy, że wahadło zostanie wychylone na bardzo mały kąt od położenia równowagi ($\theta \ll 1$ radian), możemy skorzystać z rozwinięcia funkcji sinus w szereg Taylora i ograniczyć się do pierwszego wyrazu, dokonując linearyzacji:
\[
\sin\theta \approx \theta
\]
Podstawiając to przybliżenie do równania dynamiki: $\dot{p}_\theta = -mgl\theta$. Z kolei z pierwszego równania, biorąc pochodną, mamy $\ddot{\theta} = \frac{\dot{p}_\theta}{ml^2}$. Połączmy oba równania:
\[
\ddot{\theta} = \frac{-mgl\theta}{ml^2} \implies \ddot{\theta} + \frac{g}{l}\theta = 0
\]
W ten sposób problem wahadła redukuje się rygorystycznie do jednowymiarowego, liniowego oscylatora harmonicznego o częstości kołowej $\omega_0 = \sqrt{g/l}$ oraz okresie drgań niezależnym od masy ciała ($T = 2\pi\sqrt{l/g}$).

\section{Pojecie ładunku próbnego i siła Lorentza – ruch ładunku w stałym polu elektrycznym i magnetycznym.}

\textbf{Pojęcie ładunku próbnego:}
W elektrodynamice punktowy ładunek próbny ($q_0$) to abstrakcyjny konstrukt myślowy stosowany do opisu, detekcji i pomiaru pół elektromagnetycznych, bez wpływania na stan tych pół. Definiuje się go jako obiekt o zaniedbywalnie małej wielkości przestrzennej oraz o tak nikłej wartości ładunku ($q_0 \to 0$), aby wprowadzony w dany obszar przestrzeni nie zakłócił pierwotnego rozkładu ładunków i prądów wywołujących badane pole. To na jego podstawie definiujemy precyzyjnie natężenie pola elektrycznego $\vec{E} = \lim_{q_0 \to 0} (\vec{F} / q_0)$.

\textbf{Siła Lorentza:}
W klasycznej elektrodynamice całkowita siła z jaką oddziałuje pole elektromagnetyczne na dowolny ładunek $q$ poruszający się z prędkością $\vec{v}$ to suma składowej elektrostatycznej i magnetycznej. Nosi ona miano siły Lorentza:
\[
\vec{F}_L = q\vec{E} + q(\vec{v} \times \vec{B})
\]
(Często samą część magnetyczną $q(\vec{v} \times \vec{B})$ nazywa się potocznie siłą Lorentza, jednak pełne, oficjalne wyrażenie zawiera oba człony). Z definicji iloczynu wektorowego siła magnetyczna działa zawsze prostopadle zarówno do wektora prędkości ładunku, jak i do linii pola $\vec{B}$, stąd nie wykonuje ona żadnej pracy (nie zmienia energii kinetycznej, ani modułu prędkości, a jedynie jej kierunek).

\textbf{Ruch ładunku w stałym, jednorodnym polu elektrycznym ($\vec{E} = const$, $\vec{B} = 0$):}
Jeżeli cząstka spoczywa lub wpada z prędkością do pola $\vec{E}$, doznaje stałej, niesłabnącej siły $\vec{F} = q\vec{E}$, powodującej rygorystyczne stałe przyspieszenie $\vec{a} = \frac{q}{m}\vec{E}$. Jest to ścisła analogia problemu rzutu poziomego/ukośnego ciała w polu grawitacyjnym:
- Jeżeli ładunek wbija się z początkową prędkością równoległą do linii pola, jego trajektorią pozostanie ruch prostoliniowy (jednostajnie przyspieszony lub opóźniony).
- Jeżeli ładunek wpada w pole prostopadle lub ukośnie, torem jego ruchu staje się gładka \textbf{parabola} (wykorzystywane m.in. w lampach oscyloskopowych do sterowania plamką elektronów, gdzie w locie pomiędzy dwiema okładkami odchylającymi tor zagina się parabolicznie).

\textbf{Ruch ładunku w stałym, jednorodnym polu magnetycznym ($\vec{E} = 0$, $\vec{B} = const$):}
Zjawiska zależą kategorycznie od kąta między wektorem prędkości $\vec{v}$ a wektorem pola $\vec{B}$.
1. **Ruch równoległy ($\vec{v} \parallel \vec{B}$):** Kąt wynosi zero lub $180^\circ$, sinus w iloczynie wektorowym wynosi $0$. Siła magnetyczna wynosi absolutnie zero, a ładunek porusza się swobodnie po linii prostej.
2. **Ruch prostopadły ($\vec{v} \perp \vec{B}$):** Siła wynosi maksymalnie $F = qvB$ i jest zawsze prostopadła do wektora prędkości, odgrywając fizyczną rolę idealnej siły dośrodkowej. Tor ruchu ładunku ulega rygorystycznemu zakrzywieniu do kołowego zjawiska. Cząstka krąży po zamkniętym \textbf{okręgu} z tzw. promieniem cyklotronowym (zwanym często promieniem Larmora):
\[
\frac{mv^2}{r} = qvB \implies r = \frac{mv}{qB}
\]
Jej okres obiegu (oraz kołowa częstość cyklotronowa $\omega_c = \frac{qB}{m}$) jest kompletnie niezależny od prędkości i promienia koła, co fundamentalnie pozwoliło E.O. Lawrence'owi zbudować cyklotron przyspieszający rezonansowo protony.
3. **Ruch z dowolnym kątem wlotu:** Ruch rozkłada się niezależnie na składową równoległą do linii pola ($v_\parallel$, stałą, bo siła w tej osi nie działa) oraz prostopadłą ($v_\perp$, wykonującą obrót jednostajny). Superpozycja ruchu liniowego stałego i kołowego powoduje potężne "owijanie się" cząstki wokoło linii pola w postaci **linii śrubowej (spirali geometrycznej)**. Tak objawia się zjawisko "butelki magnetycznej", która łapie kosmiczne ładunki formując pasy Van Allena nad ziemskim biegunem.

\section{Prawo Gaussa dla pola elektrycznego i magnetycznego (postać całkowa i różniczkowa). Prawo Coulomba jako konsekwencja.}

\textbf{Prawo Gaussa dla pola elektrycznego:}
Jest to pierwsze, fundamentalne równanie Maxwella, wiążące pole elektryczne ze źródłami tego pola, czyli z ładunkami. Prawo to orzeka, że strumień natężenia pola elektrycznego ($\vec{E}$) przenikający przez dowolną zamykającą się, abstrakcyjną powierzchnię (zwaną powierzchnią Gaussa, $S$) jest wprost proporcjonalny do całkowitego ładunku zamkniętego wewnątrz tej objętości ($Q_{wew}$). Zależność ta nie ma żadnego związku z rozkładem tych ładunków, liczy się jedynie sumaryczna algebraiczna wartość wewnątrz powierzchni.
\begin{itemize}
    \item \textbf{Postać całkowa:}
    \[
    \oint_S \vec{E} \cdot d\vec{S} = \frac{Q_{wew}}{\varepsilon_0} = \frac{1}{\varepsilon_0} \int_V \rho \, dV
    \]
    gdzie $\varepsilon_0$ to przenikalność próżni, a $\rho$ gęstość objętościowa ładunku w objętości $V$ zamkniętej powierzchnią $S$.
    \item \textbf{Postać różniczkowa:}
    Korzystając z twierdzenia Gaussa-Ostrogradskiego z analizy wektorowej, zamieniając całkę powierzchniową na objętościową z dywergencji wektora $\vec{E}$, otrzymujemy równanie lokalne, zachodzące w każdym punkcie przestrzeni:
    \[
    \nabla \cdot \vec{E} = \frac{\rho}{\varepsilon_0}
    \]
    Wynika z tego jasno, że ładunki elektryczne są skalarnymi źródłami pola elektrycznego (linie pola maja swoje początki na ładunkach dodatnich i końce na ujemnych).
\end{itemize}

\textbf{Prawo Gaussa dla pola magnetycznego:}
Zarazem drugie równanie Maxwella. Prawa przyrody nie ujawniły jak dotąd istnienia pojedynczych biegunów magnetycznych (tzw. monopoli magnetycznych). Linie indukcji pola magnetycznego $\vec{B}$ są liniami zawsze bezwzględnie zamkniętymi (nie mają początków ani końców, tworzą wiry wokół przewodników z prądem). Stąd strumień pola magnetycznego przez dowolną zamkniętą powierzchnię sumuje się zawsze dokładnie do zera (tyle samo linii "wchodzi", ile "wychodzi" na zewnątrz).
\begin{itemize}
    \item \textbf{Postać całkowa:}
    \[
    \oint_S \vec{B} \cdot d\vec{S} = 0
    \]
    \item \textbf{Postać różniczkowa:}
    Zastosowanie tego samego twierdzenia wektorowego, co uprzednio, obnaża fundamentalną cechę wektorowego pola $\vec{B}$, jaką jest bezźródłowość przestrzenna:
    \[
    \nabla \cdot \vec{B} = 0
    \]
\end{itemize}

\textbf{Prawo Coulomba jako konsekwencja Prawa Gaussa:}
Aby zademonstrować potęgę Prawa Gaussa, można analitycznie wyprowadzić z niego siłę oddziaływania punktowego ładunku ujętą wcześniej empirycznie prawem Coulomba. Rozpatrzmy pojedynczy, izotropowy ładunek punktowy $Q$ spoczywający w przestrzeni. 
Na zamykającą go abstrakcyjną powierzchnię wybieramy kulę (ze względu na sferyczną symetrię rzutu pola z ładunku centralnego) o promieniu r:
\[
\oint_S \vec{E} \cdot d\vec{S} = \frac{Q}{\varepsilon_0}
\]
Dzięki symetrii rzut wektora $\vec{E}$ zawsze ma ten sam kierunek co lokalny prostopadły wektor normalny $d\vec{S}$ łupiny kuli i na stałym rzędzie promienia moduł $|\vec{E}|$ pozostaje absolutnie niezmienny (stała matematyczna). Zatem przed całką dostajemy trywialne mnożenie: $E \oint dS = E \cdot 4\pi r^2$.
Stąd równanie pociąga:
\[
E \cdot 4\pi r^2 = \frac{Q}{\varepsilon_0} \implies E = \frac{1}{4\pi\varepsilon_0} \frac{Q}{r^2}
\]
Z definicji siły pola elektrostatycznego dla nowego małego ładunku próbnego $q$ wrzuconego w ową sferę, zachodzi $F = qE$. Podstawiając wyrażenie uzyskujemy finalne rygorystyczne równanie \textbf{Prawa Coulomba}: $F = \frac{1}{4\pi\varepsilon_0} \frac{qQ}{r^2}$.

\section{Prawo Faradaya (postać całkowa i różniczkowa); wzbudzanie siły elektromotorycznej w obwodach.}

Trzecie równanie Maxwella stanowi matematyczną formalizację genialnych eksperymentów Michaela Faradaya (1831), które wykazały zjawisko indukcji elektromagnetycznej, ugruntowując ścisłe sploty elektryczności z magnetyzmem u progu czasoprzestrzeni. Zasadniczą treścią prawa Faradaya jest twierdzenie, że \textbf{zmienne w czasie pole magnetyczne staje się automatycznie źródłem pola elektrycznego}, i to wysoce szczególnego: tzw. wirowego pola elektrycznego.

\textbf{Prawo Faradaya w postaci całkowej:}
Jeżeli przez dowolną zamkniętą krzywą matematyczną ułamkową pętli $C$ (tworzącą brzeg otwartej powierzchni $S$) przenika zmienny w czasie strumień wektora indukcji magnetycznej $\Phi_B$, to po obwodzie tejże krzywej wyidukowana (napędzona cyrkulacją wektora E) zostanie rygorystyczna Siła Elektromotoryczna (SEM, $\mathcal{E}$):
\[
\mathcal{E} = \oint_C \vec{E} \cdot d\vec{l} = - \frac{d\Phi_B}{dt} = - \frac{d}{dt} \int_S \vec{B} \cdot d\vec{S}
\]
Obecność wiązanego znaku minusa to hołd dla wspomnianej w poprzednich modułach reguły Lenza (praw zachowania u energii – układ broni się przed nowym wdzieraniem zmiany wywołującej). Pod całką rygorystyczną pojawia się wirowe natężenie powstałego wokół strumienia Zmiennego pola elektrycznego. Trzeba odnotować, że klasyczne elektrostatyczne pole (utworzone u Coulomba przez spoczywające swobodne ładunki punktowe) ma cyrkulację równą absolutnemu zeru ($\oint \vec{E} \cdot d\vec{l} = 0$). Zatem indukowane rygorystycznie wirowe pole $E$ wyłapuje fundamentalnie inną, bezpotencjałową naturę fizyczną.

\textbf{Prawo Faradaya w postaci różniczkowej:}
Przemiana twierdzenia całkowego Faradaya przy zastosowaniu twardej wektorowej maszynerii - mianowicie Twierdzenia Stokes'a, łączącego cyrkulację liniową pętli (rotor) z całką obszarową jej wnętrza - rzutuje na postać praw w rygorze dla pojedynczych punktów (lokalnych zawirowań).
Przechodzimy na rygor:
\[
\int_S (\nabla \times \vec{E}) \cdot d\vec{S} = \int_S \left( -\frac{\partial \vec{B}}{\partial t} \right) \cdot d\vec{S}
\]
Ponieważ owa tożsamość zmusza do równości zjawiskowych przez każde dowolnie pocięte mniejsze oczko objętości przestrzeni ($S$), same wektory funkcyjne ukryte głęboko we wnętrzu całki winny być trywialnie sobie równe. Wyprowadza to najpotężniejszą formę trzeciego Równania Maxwella różniczkowo:
\[
\nabla \times \vec{E} = -\frac{\partial \vec{B}}{\partial t}
\]
Wektorowa wielkość $\nabla \times \vec{E}$ to tzw. Rotacja lub Rotor wektora. Zatem powielenie pętli magnetyzmu z pochodną ułamka czasu owocuje zrodzeniem prostopadłego zawirowanego (rotacyjnego u osi pola czasowego) wirowego ładunku elektrycznego!

\textbf{Wzbudzanie siły elektromotorycznej (SEM) w fizycznych obwodach prądu:}
Za uwarunkowanie wyzwolenia siły w obwodzie fizycznego miedzianego drutu cewki obrotowej, zmiana strumienia $\Phi_B = \int \vec{B} \cdot d\vec{S} \approx B S \cos\alpha$ może wyniknąć fundamentalnie z 3 różnych przyczyn i mechanizmów powołanych przez inżynierię elektrodynamiczną:
\begin{itemize}
    \item \textbf{SEM transformacji:} Obwód rygorystycznie miedziany pętli tkwi sztywnie i obiektywnie nieruszony i bez żadnej fizycznej geometrii zniekształceń, podczas gdy mroczne same rygorystyczne wartości wektora na prąd magnetyczny $\vec{B}$ pęcznieją (gasną i uderzają prądem z sieci jak w przypadku zbliżenia magnesu na rurze) - jest to esencja serca Transformatora przesyłowego (drgania bez mechanicznych potęg na rzecz przeskoku prądu między cewkami).
    \item \textbf{SEM rotacji i kątowa (Generatora potęg u prądnic):} Niezwykle fundamentalne źródło dla całej u cywilizacji planety na pozysk potęgi rygorystycznej z energii. Potężny układ cewek z drutu umieszcza się bez zmiany rygoru u powierzchni w potężnym stalowym kleszczu magnesów u stałego $B$, z racji czego woda i turbina napędza z mrokiem łopat miedzianą pętlę na zmianę wyłącznie na rygor potęgi wektora kąta w czas - $\cos(\omega t)$ obracając wstrzykuje wygenerowane z rotacji pole wirowe prądu.
    \item \textbf{SEM ze zmiany pętli toru z obszarów powierzchni pola ($S$):} Geometria zgniatana lub rozciągana fizycznie prowadnicą jak i szynami obwodu u wyrw ze strumienia (użyte zjawiskowe pod indukcję magnetyczną przy miernikach, przy powolnym rygorystycznym tarciu mechanicznym na szynie hamulca oporu z przepustnic pola).
\end{itemize}

\section{Prawo Ampere’a – Maxwella (postać całkowa i różniczkowa); proces ładowania kondensatora płaskiego.}

Czwarte równanie Maxwella jest wielką kulminacją i syntezą całej elektrodynamiki. Prawa w swoim zarania badane z mroku na pętlach drutów z stałym stałym rygorystycznym u prądu przez Andre-Marie Ampere'a głosiły (użyte już u pętli w cz. 1 zestawienia) że do wyzwolenia wirowego pętlowego pola $\vec{B}$ potrzebny jest przepływ woluminu ze sznura zwykłego ładunku prądu przez oczko krzywej (prawo to rzutowało $\nabla \times \vec{B} = \mu_0 \vec{j}$).

Jednak James Clerk Maxwell w wybitnym umysłowym rewolucyjnym mroku rozpoznał rygorystycznie błąd w owych równaniach, co miało stać się gładką przyczyną by uwolnić z więzi światło jako Falę z czasoprzestrzeni (Prawo Ampere'a rygorystycznie o czystym rzucie nie spinało się bezsprzecznie pod prawami rygorów obwodów w miejscach takich jak bezprądowa wielka na wskroś próżnia kondensatora po drodze na naładowaniach).

\textbf{Ułamkowy paradoks prądu i postać całkowa Prawa Ampere'a - Maxwella:}
Mając w mroku proces wysoce silnego ładowania przez potężny drut u gęstego rygoru zwykłego prądu na blachy ciał ujętego u Kondensatora - do górnej tarczy wcieka po blachach potęga prądu przewodzenia o wartości rygorystycznej ($I$). Kiedy stworzymy wielką powierzchnię Gaussa (jak bańkę pętlową u brzegu nad drutem) w mroku owy prąd przenika ją rodząc zawirowania $\vec{B}$. Kiedy jednak rozciągniemy "wybrzuszenie brzucha z mydlanej u pętli powierzchni" aby nie przecinała ona rygorystycznego w strunie na mrok drutu przewodnika, a oplotła mrocznie potężnie we wnętrzu przestrzeń rozstępu pod tarczami obu kondensatorów u spacji pustej! Owy prąd we drutach nie rzutuje przebicia potęgi, zatem matematyka i prawa w rygorystycznym szoku nakazuje ułamkowi ułamać prawe zjawisko w zera ($I=0$), co jest potężną bzdurą niezachowania gładkich praw zjawisk natury i praw ładunków.
Maxwell wybitnie rygorystycznie u mroku dodał ratując pojęcie, mroczny, ułamkowy na ładunkach, tak zwany \textbf{Prąd przesunięcia (Id)} płynący potężnie niewidzialnie przez próżnię na pustce przestrzeni - jako zmiana i uderzenie narastającego u strzałek samego wektora Pola Elektrycznego!
Zatem rozbudowana i pełna finalna na równaniu, Poprawiona u Maxwella postać całkowa ukazała formę:
\[
\oint_C \vec{B} \cdot d\vec{l} = \mu_0 I_{wew} + \mu_0 \varepsilon_0 \frac{d}{dt} \int_S \vec{E} \cdot d\vec{S}
\]
Co dowodziło absolutnemu fenomenowi - rygorystycznie i bez prądu - \textbf{Zmienne na mroku pole Zjawiska Elektrycznego w próżni wytwarza absolutnie pełnoprawne pole Magnetyczne!}

\textbf{Postać Różniczkowa u równania Maxwella:}
Korzystając z definicji rotacji, zamiany wektora potęgi objętości prądów $I = \int \vec{j} \cdot dS$, przy tw. Stokes'a ściągając z powrotem formy, ułamek powraca i na gładkim zjawiskowym świetle ląduje słynny 4 i ostatni filar mechaniki Maxwella:
\[
\nabla \times \vec{B} = \mu_0 \vec{j} + \mu_0 \varepsilon_0 \frac{\partial \vec{E}}{\partial t}
\]
Gdzie rygorystyczny mroczny pierwszy wyraz z mrocznej gęstwy i łupku ($\vec{j}$) to twardy prąd przewodzenia z drutów (prawo klasyczne d. Ampere'a), zaś ten drugi gładki i wybitny (z rzędu pętli $\frac{\partial E}{\partial t}$) to legendarny ułamkowy dodatek - wybitna gęstość Prądu Przesunięcia u Maxwella, rygorystycznie odpowiadającego z symetryczności u łupku ze Zjawiskiem Faradaya dla fal obłoku na próżni.

\textbf{Prąd Przesunięcia a proces ładowania u kondensatora płaskiego:}
Mamy klasyczny prosty płaski kondensator z mroku powierzchni tarczy ($A$) i rygorystycznie wolnym u odstępie ($d$). Powołany do pętli do jego okładek w drutach wektory zwykłego prądu na rygor $I = dq/dt$ dodają u tarczy ładunku. Według elektrostatyki między u tarczy kondensatora u gładkiego zjawiska obłoku wektor Pola E rygorystycznie potęguje się jak ułamkiem $E = \frac{\sigma}{\varepsilon_0} = \frac{q}{A\varepsilon_0}$. Z racji pętli rygoru do mroku gładka siła wektora potężnieje o czasu różnice gęstwy.
Gęstość potężnego Maxwellowskiego i ułamkowego "Prądu przesunięcia" wewnątrz owego mrocznego rzutu i pustki to $\vec{j}_d = \varepsilon_0 \frac{d\vec{E}}{dt}$. Policzmy na rygorze jego wybitną wartość u potęg kondensatora po wymnożeniu na pole mroczne $A$:
\[
I_d = j_d \cdot A = \left(\varepsilon_0 \frac{d}{dt} \frac{q}{A\varepsilon_0}\right) \cdot A = \frac{dq}{dt} = I
\]
To absolutnie rygorystycznie fundamentalny w mrocznych wynikach cud u dowodu: Znika gładki drut i w próżni i przestrzeni we wnętrzu mroku kondensatora (nawet wyzbytego powietrza w szkle) potężnie płynie rygorystyczny ułamkowy nie z ładunku "Prąd u Przesunięcia (Id)", który na gładkim wykreślnym rachunku \textbf{co do absolutnego joty gładkiego zjawiska z wielkości Amperów równy JEST ułamkiem gładko płynącemu na zewnątrz drutów twardemu prawdziwemu rygorystycznemu Prądowi z przewodzenia (I)!} W ten rygor strumień prądu w uniwersum zawsze i wybitnie pozostaje pod pętlą całkowicie zamkniętym łupkiem ciągłości!

\section{Zasada zachowania ładunku jako konsekwencja równań Maxwella.}

\textbf{Zasada zachowania ładunku:}
Jest jednym z najbardziej ugruntowanych zjawiskowo fundamentalnych praw natury o absolutnym w mroku zrębów rygorze. Głosi ona na ułamek pętli matematycznej i słownie, iż całkowita mroczna suma wszystkich i absolutnie dodatnich, i obłocznie ujemnych ładunków u gęstwy we wnętrzu fizycznie Odizolowanego u rzutu klatki na ciał układu fizycznego bez cienia pętli do otoczenia, absolutnie pozostaje Stała na osi przebiegu u czasu rygorystycznego wektora mroku. Jeśli zaś wybitnie z otwartego gniazda uładunek znika w czasie ($\frac{dq}{dt}$ u ubytku mroku ujemnie ucieka ułamkiem mrocznych z gąszczy tarczy na powłoki do potęgi wybitnie), musi gładkim prawem gęstwy wektora potężnego i zrębem stać się wyłaniającym się rygorystycznie fizycznym prądem przewodzonym, wypływającym przez ową mroczną w ułamku zrębów ścianę powierzchni ($I = \oint \vec{j} \cdot d\vec{S}$). Formalizm owej ciągłości to klasyczne potężnie \textbf{Równanie ciągłości ładunku} w skali wybitnie ułamkowo różniczkowej gniazda u mroku:
\[
\nabla \cdot \vec{j} + \frac{\partial \rho}{\partial t} = 0
\]
gdzie u mroku wybitnie $\vec{j}$ stanowi wektor lokalnej małej gęstości z pętli prądu a gładkie zrębem ciał $\rho$ u mroku ułamku to rygorystyczna zebrana w punkt przestrzeni wolumetryczna mroczna z gęstwy ciał ładunkowa potęga potężnego u gęstości (kulomb na sześcian z zjawiska metra potęg)!

\textbf{Wyprowadzenie mroczne i rygorystyczne z filarów ciał u Maxwella:}
Wybitna geniuszu mroku w ułamku potężnego splecionego u Maxwella matematycznego geniuszu ciał rzutu tkwiła również wybitnie na tym potęg ułamku, iż potężnie dołączony od gładkich jego wybitnych u mroku dłoni dodatkowy "Prąd z Przesunięć" upewnił i z cementował ciał fundamentalnie zasady! Owo potężnie Równanie o ułamkach Ciągłości dla zrębów ładunku nie musiało gładkim zjawiskiem być postulatem niezależnym doczepionym sztucznie do kartki – wynikało czystym zderzeniem praw wprost i twardo do mechanicznego rygorystycznego zrębów u różniczek \textbf{wewnątrz Prawa zrębów ciał u rygoru Ampere'a-Maxwella i zjawiska u Gaussa}.

Dla twardego na mroku zjawisk u dowodu, weźmy na warsztat ciał to owo zmodyfikowane prawo:
\[
\nabla \times \vec{B} = \mu_0 \vec{j} + \mu_0 \varepsilon_0 \frac{\partial \vec{E}}{\partial t}
\]
Podzielmy dla czystego ułamkowego rygoru całe równanie przez nieinwazyjny na zręb mroku człon u potęg o ułamek z $\mu_0$:
\[
\frac{1}{\mu_0} \nabla \times \vec{B} = \vec{j} + \varepsilon_0 \frac{\partial \vec{E}}{\partial t}
\]
Zaaplikujmy u mrocznej gładkiej, matematycznie i zjawiskowo potężnej operacji tak zwanej wektorowej dywergencji wektora na mroku z ułamkiem u po obydwu z ciał gładkich u łupkach na stronach u osi tego równania. Aplikujemy (nakładamy i obłożymy dywergencją $\nabla \cdot$ lewą u potęg rygorystyczną oraz ułamkową i mroczną prawą burtę). Z podstawowych zrębów analiz wektorów wyłania się fundamentalna zasada o zręb u mroku: \textbf{Dywergencja z każdej rotacji jakiegokolwiek w mroku rygorystycznego wektora daje absolutnie ułamkowe rygorystyczne Zero na twardym betonie praw i pętli ($\nabla \cdot (\nabla \times \vec{B}) \equiv 0$)}. Lewa mroczna i gładka ze skrzyń na strona równania po wymuszeniu u zerowania znika. Pozostaje z ciał bezwzględna powłoka prawej z pętli strony:
\[
0 = \nabla \cdot \vec{j} + \nabla \cdot \left( \varepsilon_0 \frac{\partial \vec{E}}{\partial t} \right)
\]
Matematycznie pod mrocznym pędem pochodna ułamku czasu po wektorze ($\partial/\partial t$) ze zmiennych wymiarowych jest do gładkiej zamiany i wyrzutu niezależnie komutatywna w pętli dla wektora u mrocznym operatora dywergencji na gładkiej osi rygorystycznego ciał przestrzeni ($\nabla \cdot$). Możemy zatem zamienić bezpiecznie gładkie kolejności ciał mroku różniczek:
\[
0 = \nabla \cdot \vec{j} + \frac{\partial}{\partial t} \left( \varepsilon_0 (\nabla \cdot \vec{E}) \right)
\]
Spójrzmy na wielkość zamkniętą w rygorystycznym ze zrębów okrągłym u mroku z ciał nawiasie. Tam ciał od ułamka gęstwy to gładkie nic o prawach innego w pętli wybitnych cudów jak ułamek Prawo Gaussa o ułamkach na rygor z rzędu dla zjawisk elektryczności ($\nabla \cdot \vec{E} = \frac{\rho}{\varepsilon_0}$)!, Po uprzątnięciu rygorystycznie i gładkim pod mrok zjawisk ciał u skróceniu pętli $\varepsilon_0$:
\[
0 = \nabla \cdot \vec{j} + \frac{\partial}{\partial t} \left( \varepsilon_0 \frac{\rho}{\varepsilon_0} \right) \implies \nabla \cdot \vec{j} + \frac{\partial \rho}{\partial t} = 0
\]
Dowód z wstrząsającą czystością na mroku jest absolutnie bezwzględnie z ułamka zjawisk kompletny! Potwierdza, że układy równań elektrodynamiki są u mroku ciał na wielkim rygorze na fundamentach gładkiej czasoprzestrzeni ciał spójne!

\section{Płaskie fale elektromagnetyczne jako rozwiązania równań Maxwella w próżni.}

Fakt i matematyczna tytaniczna twarda konkluzja w rygorze dla równań Maxwella stanowi u wybitnego ułamka z ciał, że pola elektryczne ($\vec{E}$) ze spacji czasu i magnetyczne ($\vec{B}$) w mrocznej gładkiej czasoprzestrzeni ze zrębem o wolnej gęstej pustej i gładkiej Próżni fizycznej (przy całkowitym na rygor absolutnym zjawisku braku twardych ładunków, czyli gęstość ułamkowa ładunku $\rho=0$ z rygoru potęg, z brakiem w ułamku z prądów $j=0$) mogą nie tylko obiegać potęgi ale całkowicie swobodnie i wolno \textbf{odrywać się od źródeł} i niezależnie po oddaleniu propagować po ułamkowych ciał przestrzeniach wszechświata na mroku w postaci drgającej do zjawiska Fali.

Wielki punktowy model rzutu od mrocznej próżni nakazuje obarczyć równania Maxwella wyzerowaniem gładkim starych ułamków ciał materiału (zero ładunków, zero rygoru potęg u wektorów $J$):
1. $\nabla \cdot \vec{E} = 0$ (Pole $E$ mrocznym zerem nie wybija na gładkie rygorystyczne na wolno pędy).
2. $\nabla \cdot \vec{B} = 0$ (Brak ułamu w zjawisk mroku ładunków u biegunów).
3. $\nabla \times \vec{E} = -\frac{\partial \vec{B}}{\partial t}$ (Prawo Faraday'a gęstwy i zrębów o czasie mroku).
4. $\nabla \times \vec{B} = \mu_0 \varepsilon_0 \frac{\partial \vec{E}}{\partial t}$ (Prąd Przesunięcia o tarczy wybitnie sam pcha na mrok pole na oś czasu $E$).

\textbf{Równanie falowe próżni u rygorów:}
Aby w mroku matematyki za rygor otrzymać ułamkowe z ciał drgające z czasem wybitnie gładkie i od osi Równanie ułamek Falowe, na Równanie ciał Maxwella obłoku nr. 3 nakłada się wybitnie ponownie matematyczny wektor po stronie rygoru obydwu z osi - z mroku ponownie operator rygorystyczny rotacji u gładkim ciał wektora. Z tożsamości z analiz matematyki $\nabla \times (\nabla \times \vec{E}) = \nabla(\nabla \cdot \vec{E}) - \Delta \vec{E}$ i korzystając w gęstwy wybitnie z pierwszego rygoru mroku o zerowej u pętli gęstwy ułamu ładunków ($\nabla \cdot \vec{E} = 0$ w pustce u próżni) po na zręby połączeniach lewa i potężna ze w mroku strona i podstawieniu po rygorystyce Prawej pętli i z ciał na gęstwie za z rygorem ciał $\nabla \times B$ równania do pętli mroku nr 4 wybitnie uzyskujemy bezapelacyjny ułamkowy i rygorystycznie podziwu cud fizyki rzutu na matematyce ułamka fali w czasie!:
\[
\Delta \vec{E} = \mu_0 \varepsilon_0 \frac{\partial^2 \vec{E}}{\partial t^2} \quad \text{oraz z potęg } \quad \Delta \vec{B} = \mu_0 \varepsilon_0 \frac{\partial^2 \vec{B}}{\partial t^2}
\]

\textbf{Rozwiązania (Płaska Fala o pętli potęg rygoru elektromagnetycznego):}
Oba rygorystycznie wydarte na pętle ciał równania mają gładki i wybitny kształt idealnego gładkiego w pętlę Równania rzutu na Falę $\Delta \psi = \frac{1}{v^2} \frac{\partial^2 \psi}{\partial t^2}$ gdzie po wyrzuceniu dla pętli stałej wyłoniła się u pętli słynna potęgowa rygorystycznie i ułamkowo ze zjawiska osi czasu fundamentalna Prędkość u mroku propagacji mrocznych gładkich fal w ciele na wolnej gęstej próżni. W świetle potęg matematyki Maxwella stała stała się światłem!:
\[
v = c = \frac{1}{\sqrt{\mu_0 \varepsilon_0}} \approx 3 \cdot 10^8 \, \text{m/s}
\]
Najprostszym w ułamku, a rygorystycznie bezwzględnie dopasowanym na zręb rozwiązaniem dla obu wielkich potęg mroku w polach pętli są drgające ze zjawiskiem, w mroku równoległe (co u osi rygoru do frontu płaskie) Płaskie Fale Harmoniczne twardej o drganiu sinusoidy przemieszczających na tor ułamkowej fali ($z$) od rygoru (reprezentacja z liczb w exp u urojonej zespolonej postaci) w pętli do mrocznych obłokach rygoru fal czasoprzestrzennych z wektorem fali fali gęstej ($k$) i wybitnie czasowej i mrocznej o pędzie pulsu kołowego ($\omega$):
\[
\vec{E}(z,t) = \vec{E}_0 e^{i(kz - \omega t)} \quad \text{oraz} \quad \vec{B}(z,t) = \vec{B}_0 e^{i(kz - \omega t)}
\]

Z równań tych, u mrocznego i potęg gęstych splecień w Maxwella wychodzi i za ułamek wyrzutu płaskiego 3 gładkich i wybitnie rygorystycznych w przyrodzie u ciał cech praw fal w rzucie Fotonów z pętli u fali e-m:
\begin{itemize}
    \item Pętle z rygorów o Fali to w 100\% w zjawisk poprzeczne (dywergencja gładko bez ciał ułamu pętli dowodzi z wektora fali k i prawych $\vec{k} \cdot \vec{E} = 0$, więc pole mroku jest ciał pędem na oś fal u ciał k!).
    \item Na gładkim obłoku ciał Obydwa wektory wybitnych u potęg wektorów na Pola $\vec{E}$ rygorystycznie na $\vec{B}$ są potężnie w każdej sekundzie fali bezwzględnie Prostopadłe w rzucie z ciał w obłoku i przestrzeni na zjawisko fali Do Samego rygoru w mroku Siebie Nawzajem z rygorystycznie zgodną \textbf{Tą samą idealną ze zjawisk u pętli ułamu mrocznego Fazą o czasu skoku mroku z rygorów (górki osiągają 0 i pik rygorystycznie pętli z ciał bezwzględnie razem z ułamu u drgań fal u mroku!) }.
    \item Rygorystyczny ułamek ze ciał mroku a zarazem gładki i wybitny Moduł o osi pola rygoru fali Elektrycznej we zderzeniu zawsze gładkim u rygoru równa się na pętli wybitnemu Modułowi z pola wektora u mrocznym mroku na pętlę rygorystycznym u fali ciał i przestrzennej pętli osi z o magnetycznym w ułamku przemnożonym gładkim na pęd Prędkości $E_0 = c B_0$.
\end{itemize}

\section{Przykłady eksperymentów, które prowadziły do powstania teorii kwantowej; porównanie z fizyką klasyczną.}

Przełom XIX i XX wieku przyniósł szereg wyników eksperymentalnych, które w sposób drastyczny i niepodważalny łamały przewidywania klasycznej mechaniki Newtona oraz elektrodynamiki Maxwella. Konieczność ich wyjaśnienia zrodziła nową gałąź nauki – mechanikę kwantową.

\textbf{1. Promieniowanie ciała doskonale czarnego (1900 r.):}
\begin{itemize}
    \item \textbf{Oczekiwania klasyczne:} Zgodnie z termodynamiką i elektrodynamiką klasyczną, wyprowadzono prawo Rayleigha-Jeansa opisujące rozkład gęstości energii promieniowania wnękowego. Przewidywało ono, że natężenie promieniowania rośnie nieograniczenie wraz ze spadkiem długości fali ($\rho \propto 1/\lambda^4$). Prowadziło to do tzw. \textit{katastrofy w nadfiolecie} – ciało w dowolnej temperaturze musiałoby emitować nieskończoną ilość energii, co jest fizycznie niemożliwe.
    \item \textbf{Rozwiązanie kwantowe:} Max Planck dopasował do danych eksperymentalnych nową krzywą, wprowadzając zdesperowany postulat: energia oscylatorów w ściankach wnęki nie może zmieniać się w sposób ciągły. Energia może być emitowana i absorbowana wyłącznie w ściśle określonych, dyskretnych paczkach (kwantach) o wartości $E = nh\nu$ (gdzie $h$ to stała Plancka, a $\nu$ to częstotliwość). Wynik ten idealnie pokrył się z widmem eksperymentalnym na wszystkich długościach fal.
\end{itemize}

\textbf{2. Zjawisko fotoelektryczne zewnętrzne (1905 r.):}
\begin{itemize}
    \item \textbf{Oczekiwania klasyczne:} Klasyczna fala elektromagnetyczna powinna z biegiem czasu "rozbujać" elektrony w metalu, dostarczając im energii w sposób ciągły. Oczekiwano, że emisja elektronów zajdzie po pewnym czasie naświetlania dowolnym światłem, a energia wybijanych elektronów będzie zależeć od natężenia światła (amplitudy fali).
    \item \textbf{Wyniki eksperymentalne i kwantowe rozwiązanie:} Doświadczenia Lenarda ukazały coś zupełnie odwrotnego: zjawisko zachodzi \textit{natychmiastowo}, ale wysoce wybiórczo. Dla fal zbyt długich (np. światła czerwonego) elektrony nigdy nie były wybijane, niezależnie od użytej mocy lampy. Z kolei słabe światło fioletowe natychmiast wybijało elektrony. Albert Einstein wyjaśnił to postulując, że światło samo w sobie to strumień cząstek – \textit{fotonów} o energii $E = h\nu$. Energia kinetyczna wybitego elektronu wynosi $E_k = h\nu - W$ (gdzie $W$ to praca wyjścia). Decyduje pojedynczy akt zderzenia fotonu z elektronem, a nie ciągłe grzanie falą.
\end{itemize}

\textbf{3. Efekt Comptona (1923 r.):}
Rozpraszanie promieniowania rentgenowskiego na swobodnych elektronach wykazało, że ugięte promieniowanie zmienia swoją długość fali na dłuższą. Na gruncie elektrodynamiki klasycznej elektron uderzony falą powinien oscylować i reemitować falę o \textit{dokładnie tej samej} częstotliwości. Arthur Compton udowodnił jednak, że fotony zachowują się jak cząstki klasyczne obdarzone pędem $p = h/\lambda$. Zjawisko to potraktowane po prostu jak zderzenie sprężyste dwóch bil kręglowych (fotonu i elektronu) – z praw zachowania energii i pędu – dało wzór idealnie oddający eksperymentalne przesunięcie $\Delta \lambda$, koronując korpuskularną naturę światła.

\textbf{4. Widma liniowe atomów i stabilność atomu:}
Zgodnie z klasycznym modelem Rutherforda, elektron na orbicie doznaje przyspieszenia dośrodkowego, zatem (zgodnie z prawami Maxwella) musi ciągle emitować promieniowanie, tracić energię i w ułamku sekundy spaść na jądro. Widmo emisji powinno być ciągłe. Eksperymenty pokazywały jednak bezwzględną stabilność atomów i ostre linie widmowe (np. serie Balmera). Wyjaśnił to kwantowy model Bohra, postulujący istnienie \textit{stanów stacjonarnych}, w których elektron nie promieniuje, a emisja następuje jedynie jako dyskretny przeskok (kwant) między dozwolonymi orbitami ($h\nu = E_m - E_n$).

\section{Zasada nieoznaczoności w fizyce kwantowej. Kiedy występuje dla danej pary wielkości fizycznych.}

\textbf{Zasada nieoznaczoności Heisenberga:}
Sformułowana przez Wernera Heisenberga w 1927 roku zasada mówi o fundamentalnym, wpisanym w samą naturę Wszechświata, ograniczeniu określoności wielkości fizycznych. Nie jest ona wynikiem niedoskonałości przyrządów pomiarowych, lecz falowej natury materii. Orzeka ona, że nie jest możliwe jednoczesne wyznaczenie w tym samym stanie kwantowym, ze stuprocentową dokładnością, wartości pewnych par wielkości fizycznych (zjawisko to określa relacja komutacji).

Najsłynniejsza relacja dotyczy pary \textbf{położenie ($x$) i pęd ($p_x$)} w tym samym kierunku:
\[
\Delta x \cdot \Delta p_x \ge \frac{\hbar}{2}
\]
Gdzie $\Delta x$ i $\Delta p_x$ to odchylenia standardowe (niepewności pomiarowe) tych wielkości dla danego stanu kwantowego, a $\hbar = h/(2\pi)$ to zredukowana stała Plancka. Oznacza to, że im dokładniej wyznaczymy pozycję cząstki (lokalizując ją przestrzennie, ściskając $\Delta x \to 0$), tym bardziej nieokreślony staje się jej pęd (cząstka dostaje ogromnego rozrzutu położeń w momentach późniejszych, bo $\Delta p_x \to \infty$). W obrazie falowym wynika to z faktu, że aby zbudować bardzo wąską paczkę falową przestrzennie, musimy zsumować ze sobą fale z ogromnego, ciągłego spektrum długości fal (czyli pędów, bo $p = \hbar k$).

Podobne relacje wiążą inne specyficzne pary sprzężonych wielkości fizycznych, na przykład:
\begin{itemize}
    \item \textbf{Kąt obrotu i moment pędu:} $\Delta \phi \cdot \Delta L_z \ge \frac{\hbar}{2}$ (ścisłe zdefiniowanie momentu pędu np. w orbitach kwantowych wokół osi z całkowicie rozmywa położenie kątowe elektronu).
    \item \textbf{Energia i czas:} $\Delta E \cdot \Delta t \ge \frac{\hbar}{2}$ (co ma ogromne implikacje dla tzw. czasu życia stanów wzbudzonych atomów – im krócej cząstka przebywa w jakimś stanie, z tym większym rozmyciem energetycznym (szerokością naturalną linii widmowej) ów stan się wiąże. Pozwala to też na pojawianie się urojonych "cząstek wirtualnych" z próżni na czas łamiący zasadę zachowania energii, byle w ramach $\Delta t$).
\end{itemize}

\textbf{Kiedy należy oczekiwać wystąpienia zasady nieoznaczoności dla danej pary?}
W formalizmie mechaniki kwantowej (operatory w przestrzeni Hilberta), każdą mierzalną wielkość fizyczną (obserwablę) reprezentuje operator hermitowski $\hat{A}$. Wynik jednoczesnego pomiaru zależy od tego, czy operatory te są przemienne. Miara tego jest \textbf{komutator} dwóch operatorów:
\[
[\hat{A}, \hat{B}] = \hat{A}\hat{B} - \hat{B}\hat{A}
\]
Zjawisko nieoznaczoności pomiędzy dwiema wielkościami występuje \textbf{zawsze i tylko wtedy, gdy operatory odpowiadające tym wielkościom fizycznym nie komutują (ich komutator jest różny od zera)}: $[\hat{A}, \hat{B}] \neq 0$. 
Matematycznie uogólniona zasada nieoznaczoności Robertsona brzmi:
\[
\Delta A \cdot \Delta B \ge \frac{1}{2} \left| \langle [\hat{A}, \hat{B}] \rangle \right|
\]
Na przykład: operatory położenia i pędu w tej samej osi dają komutator $[\hat{x}, \hat{p}_x] = i\hbar$, co bezpośrednio generuje nierówność Heisenberga. Z kolei operatory położenia w osi x i pędu w osi y komutują $[\hat{x}, \hat{p}_y] = 0$, dlatego można jednocześnie znać pozycję cząstki $x$ i jej pęd pionowy $p_y$ z dowolną, absolutną precyzją.

\section{Efekt tunelowy: stacjonarne rozwiązanie równania Schrödingera i przykłady zjawisk.}

Efekt tunelowy to na wskroś kwantowe i niefizyczne na gruncie Newtonowskim zjawisko przenikania cząstek przez zapory energetyczne bez posiadania wystarczającej energii do "przeskoczenia" nad zaporą.
W mechanice klasycznej uwięziona cząstka o energii $E$ rzucona na stacjonarną barierę (mur) o potencjale $V_0$ gdzie $E < V_0$ odbije się rygorystycznie i bezwzględnie w 100\% przypadków. Cząstka nie ma możliwości przebywania wewnątrz bariery.

\textbf{Rozwiązanie równania Schrödingera dla prostokątnej bariery potencjału:}
Rozważmy barierę przestrzenną $V(x) = V_0$ dla z przedziału $0 \le x \le a$ oraz $V(x) = 0$ poza tym obszarem. Niech z lewej strony (od $x=-\infty$) pada wiązka cząstek o ustalonej energii $E < V_0$. Używamy stacjonarnego (niezależnego od czasu) równania Schrödingera w jednym wymiarze:
\[
-\frac{\hbar^2}{2m} \frac{d^2\psi}{dx^2} + V(x)\psi(x) = E\psi(x)
\]
Przestrzeń dzielimy na trzy obszary. Rozwiązujemy to równanie kawałkami:
\begin{itemize}
    \item \textbf{Obszar I ($x < 0$, przed barierą, $V=0$):} Równanie ma postać $\psi'' + k^2\psi = 0$, gdzie $k = \sqrt{2mE}/\hbar$. Rozwiązaniem jest superpozycja płaskiej fali padającej i odbitej: $\psi_I(x) = A e^{ikx} + B e^{-ikx}$.
    \item \textbf{Obszar II ($0 \le x \le a$, wewnątrz bariery, $V=V_0$):} Ponieważ $E < V_0$, równanie przybiera formę $\psi'' - \kappa^2\psi = 0$, gdzie $\kappa = \sqrt{2m(V_0 - E)}/\hbar$. Rozwiązaniem ułamkowym nie jest już funkcja oscylacyjna (co byłoby niefizyczne u klasyka), ale kombinacja eksponentów tłumionych i rosnących: $\psi_{II}(x) = C e^{-\kappa x} + D e^{\kappa x}$. Wewnątrz bariery funkcja falowa nie znika natychmiast, lecz szybko \textit{wykładniczo zanika}.
    \item \textbf{Obszar III ($x > a$, za barierą, $V=0$):} Rozwiązaniem jest jedynie rozchodząca się fala (bo nie ma odbicia z plus nieskończoności): $\psi_{III}(x) = F e^{ikx}$.
\end{itemize}

Ostateczne stałe $A, B, C, D, F$ uzyskuje się ze \textbf{zszycia funkcji falowej}, czyli z narzucenia matematycznych warunków ciągłości zarówno samej funkcji $\psi(x)$, jak i jej pierwszej pochodnej $\psi'(x)$ we współrzędnych kontaktowych $x=0$ oraz $x=a$. Zespolony Współczynnik Transmisji zdefiniowany jest jako $T = |F/A|^2$. Z obliczeń układu tych 4 równań analitycznych wynika, że jeśli bariera jest długa lub gruba (duże $\kappa a$), czynnik ten można zredukować do formy:
\[
T \approx T_0 e^{-2\kappa a} = T_0 \exp\left( -2\frac{\sqrt{2m(V_0 - E)}}{\hbar} a \right)
\]
Zauważmy absolutny cud teorii falowej: współczynnik ten, choć ułamkowo mały i stłamszony przez wykładnicze uderzenie szerokości i grubości energetycznej bariery, \textbf{nie jest zerem}! Istnieje stacjonarne prawdopodobieństwo, że cząstka "przeciśnie się" do obszaru III niczym duch.

\textbf{Przykłady zjawisk fizycznych objaśnianych efektem tunelowym:}
\begin{itemize}
    \item \textbf{Rozpad Alfa ($\alpha$) jąder atomowych (teoria Gamowa):} Nukleony spięte są we wnętrzu potężną studnią potencjału (oddziaływania silne), odgradzaną bardzo wysoką kulombowską barierą. Choć cząstka alfa nie ma kinetycznej energii, by przeskoczyć barierę, tuneluje przez nią powolnym prawdopodobieństwem (co tłumaczy czasy półtrwania rzędu miliardów lat u Izotopów, jak i rzędu mikrosekund przy niewielkich wariacjach bariery – zgodność z regułą Geigera-Nuttalla).
    \item \textbf{Fuzja termojądrowa w Słońcu:} Temperatury w rdzeniu (15 mln Kelwinów) nie dają z energii kinetycznej plazmy dość mocy na pokonanie potężnego odpychania elektromagnetycznego zderzających się jąder wodoru (klasycznie Słońce nie powinno świecić). Jednak dzięki zjawisku tunelowania reakcja fuzji zyskuje kwantowe wsparcie.
    \item \textbf{Skaningowy mikroskop tunelowy (STM):} Urządzenie przykładające super-ostrą igłę nanometryczną odległość od tarczy. Napięcie elektryczne tuneluje elektrony pomiędzy igłą i próbką. Skala prądu ($e^{-2\kappa a}$) jest tak gigantycznie czuła na najdrobniejszą zmianę dystansu, że oprogramowanie wyrysowuje tarcze dosłownie ułamkowego kształtu pojedynczych ułożonych atomów!
\end{itemize}

\section{Równanie Schrödingera dla cząstki kwantowej: omów zachowanie się swobodnej cząstki kwantowej.}

Erwin Schrödinger w 1926 roku (zainspirowany przez falową hipotezę de Broglie'a) wysunął potężne twierdzenie ułamkowe: ewolucja każdego nie-relatywistycznego kwantowego układu mechanicznego może być z całkowitym zdeterminowaniem opisana na w oparciu o wektor stanu zamknięty w postaci funkcji falowej $\Psi(\vec{r}, t)$, ewoluujący według odpowiedniego równania operatorowego, opierającego się bezpośrednio na konserwatywnych założeniach Hamiltonianu.

\textbf{Równanie Schrödingera:}
Zależne od czasu (ewolucyjne) równanie Schrödingera wywodzi się przez zamianę w klasycznym wzorze na energię ($E = \frac{p^2}{2m} + V$) odpowiednich klasycznych wielkości przez kwantowe odpowiedniki w postaci operatorów ($E \to i\hbar \frac{\partial}{\partial t}$, $\vec{p} \to -i\hbar \nabla$). Prowadzi to do podstawowego równania falowego świata mikroskopowego:
\[
i\hbar \frac{\partial \Psi}{\partial t} = \hat{H} \Psi \implies i\hbar \frac{\partial \Psi(\vec{r}, t)}{\partial t} = \left( -\frac{\hbar^2}{2m}\nabla^2 + V(\vec{r}, t) \right) \Psi(\vec{r}, t)
\]
Kwadrat modułu funkcji falowej ($|\Psi(\vec{r}, t)|^2$) to gęstość prawdopodobieństwa znalezienia cząstki w danym punkcie przestrzeni i w danej chwili czasu (interpretacja Borna).

\textbf{Swobodna cząstka kwantowa (Ruch swobodny):}
Cząstka swobodna to taka, na którą nie działają żadne siły fizyczne (zewnętrzne pole potencjału jest zerowe w całej przestrzeni, zatem wstawiamy $V(\vec{r}, t) \equiv 0$). Zakładając przypadek jednowymiarowy (dla osi $x$) u układzie, równanie do rozplatania przyjmuje bezpotencjałową formę:
\[
i\hbar \frac{\partial \Psi}{\partial t} = -\frac{\hbar^2}{2m} \frac{\partial^2 \Psi}{\partial x^2}
\]
Ponieważ potencjał nie zależy od czasu (jak zresztą każde $V=0$), używamy metody separacji zmiennych do wydobycia stanów stacjonarnych $\Psi(x,t) = \psi(x)e^{-iEt/\hbar}$. Część przestrzenna dla układów bezczasowych musi spełniać równanie oscylacyjne: $\frac{d^2\psi}{dx^2} + k^2\psi = 0$, z wielkością liczby falowej $k = \frac{\sqrt{2mE}}{\hbar}$.

Prowadzi to do ogólnego rozwiązania w postaci gładkiej \textbf{zespolonej płaskiej fali harmonicznej}:
\[
\Psi_k(x,t) = A e^{i(kx - \omega t)}
\]
gdzie u ułamka mrocznych fal $\omega = E/\hbar$ jest częstotliwością kołową. Cząstka swobodna z określonym stanem fali płaskiej dysponuje idealnie i nieskończenie określonym, ostrym rygorystycznym pędem wynoszącym $p = \hbar k$. Ale ta zaleta u Heisenberga stanowi potworną fizyczną udrękę – gęstość prawdopodobieństwa odnalezienia takiej fali dla tego stanu $|\Psi_k|^2 = |A|^2$ okazuje się po rygorze na obłoku absolutnie \textbf{stała (jednostajna) we wszechświecie na całej przestrzeni ($-\infty, \infty$)}! Cząstka z idealnie znanym pędem znajduje się wszędzie ułamkowo ze stałym losem. Takiej funkcji falowej nie da się poprawnie znormalizować w sensie ułamu rachunków całkowych ($\int |A|^2 dx \to \infty$), więc nie ma ona bytu fizycznego sensu jako oddzielna "punktowa cząstka".

Aby uświadczyć na obłoku fizyczną kwantową i mroczną cząstkę swobodną (zlokalizowaną w przestrzeni do skrzyżowania z tolerancją $\Delta x$), nie operujemy nieskończoną idealną falą płaską $\Psi_k$, a sumujemy z nich jako baz transformaty Fouriera całe \textbf{paczki falowe}:
\[
\Psi(x,t) = \frac{1}{\sqrt{2\pi}} \int_{-\infty}^{\infty} g(k) e^{i(kx - \omega(k) t)} dk
\]
Dla ułamkowej paczki falowej swobodnej cząstki kwantowej nieodzowną ułamkową kwantową konsekwencją dyspersji równania $E = p^2/2m \implies \omega(k) \propto k^2$ na obłoku pętli w próżni jest tzw. ułamkowe mroczne \textbf{rozpływanie się paczki falowej z upływem czasu}. Grupa i centrum masy pędzą sobie ze średnim oczekiwanym w próżni pędem bez przyspieszeń (jako środek lokalizacji cząstki wg praw Ehrenfesta dla dynamiki swobodnej). Jednak wariancja rozmycia przestrzennego ($\Delta x$) bezwzględnie rygorystycznie puchnie i się degraduje poszerzając – elektron wrzucony w próżnię do przestrzeni pętli o początkowo ciasno zdefiniowanej u Gaussa niepewności w punkcie wyjściowym po długim czasie zjawiska sam z siebie staje się niezlokalizowanym mrocznym, grubym ułamkowym szumem o gigantycznym rozmyciu pętli! 

\section{Funkcje falowe układów cząstek kwantowych i równanie Schrödingera: czym charakteryzują się bozony i fermiony.}

W klasycznej fizyce układ dwóch cząstek możemy zawsze traktować jako w pełni rozróżnialny (możemy każdej przyporządkować "numerek" lub farbę, prześledzić precyzyjnie ich tory, by zawsze wiedzieć "która jest którą").
W rygorach mikroskopowych mechaniki kwantowej zasada nieoznaczoności zamazuje trajektorie położeń, sprawiając potężnie, że identyczne cząstki elementarne tego samego typu (np. dwa elektrony ułamu w atomie helu, bądź dwa fotony) uważa się za absolutnie i fundamentalnie \textbf{nierozróżnialne}. Właściwość ta narzuca wybitne ułamkowe na obłoku kwantowe ograniczenia na budowę wielocząstkowej funkcji falowej opisującej stan całego układu, niosące dla struktury Wszechświata konsekwencje nie do podważenia.

\textbf{Funkcja falowa a zasada antysymetryzacji Pauli'ego (Twierdzenie o związku spinu ze statystyką):}
Dla układu wybitnie na mroku wielu identycznych z rygorów ciał kwantowych, Hamiltonian operujący całym gniazdem atomowym, rygorystycznie i wymownie jest i absolutnie na obłoku symetryczny na zamianę obydwu (jakichkolwiek z potęg) współrzędnych cząstek $(\vec{r}_i, s_i) \leftrightarrow (\vec{r}_j, s_j)$ (miejsc w argumencie zespolonej wielocząstkowej i gęstej o mroku funkcji na pętli). Gęstość ułamków prawdopodobieństwa $|\Psi(1, 2, \dots)|^2$ bezwzględnie musi być niewrażliwa i identyczna po zamianie ciał w układzie numerków u ciał mroku ($1 \leftrightarrow 2$). Sama u ułamku pętli zespolona funkcja u falowej paczki ($\Psi$) posiada na rygor zderzenia jednak wybitnie 2 poprawne i wykluczające się możliwości ewolucyjne: po permutacji cząstek z uderzenia wybitnie albo u pętli zwraca obłokowi gładko ten sam wyjściowy znak (funkcja ułamkowa symetryczna), albo rzuca ułamek pętli z przeciwnym wektorem znaku i spacji ciał w ciał fazie (funkcja anty-symetryczna $\Psi(1,2) = -\Psi(2,1)$).

Wolfgang Pauli objaśnił ten szokujący na twardym betonie paradoks włączając Spin do równań i udowodnił (co jest filarem Modelu Standardowego), że \textbf{cały znany kosmos cząstek dzieli się wybitnie i rygorystycznie na dwa klany ustroju kwantowego, powiązane ich wrodzonym momentem pędu krętu (spinem)}:

\textbf{1. FERMIONY:}
Cząstki dysponujące tak zwanym \textbf{połówkowym spinem ułamkowym} (np. $s = 1/2, 3/2, \dots$ wyrażone na jednostkach $\hbar$). 
\begin{itemize}
    \item Funkcja wielocząstkowa układu fermionów musi bezwzględnie w mrocznych ciałach opisu pozostawać ze znakiem \textbf{antysymetrycznym} przy zamianie ze sobą miejsc u ciał obojętnie rygorystycznych 2 identycznych cząstek! ($\Psi(1,2) = -\Psi(2,1)$)
    \item Prowadzi to na ułamek osi rygoru bezwzględnie do sformułowania \textbf{Zakazu Pauliego}. Jeśli dla pętli na eksperymencie próbowano by wprowadzić, wpakować mroczne 2 te same fermiony idealnie bez ułamu pętli z do dokładnie z osi tej samej joty wybitnie do gładkiego Stanu kwantowego (ten sam pęd z gęstwy i spin ciał u orbity np. do gładkiej funkcji), z tożsamości po narzucie ciał dla antysymetrii w rygor wyszło by po odjęciu bezwzględnie ze wzoru $\Psi = 0$. Zjawisko u gniazd mroku na pętlę to oznajmia: \textbf{Dwa fermiony nie mogą w mroku atomu bezwzględnie zajmować tego samego powłokowego stanu na obłoku u orbitali!}
    \item Podlegają i statystycznie spoczywają pod gładką w mroku powłoką rzutowania o ułamkach na potęgi zjawiskowych krzywych rygorystycznie ze \textbf{Statystyki Fermiego-Diraca}. Konieczność zabudowania mrocznego "piętra i schodków powłokowych po 2 gładkie u dołu spinu do mroku elektrony i szczebel na górę u potęgi wybitnie kolejnego orbitala w atomach do góry na Mendelejewie" zamiast runięcia do najniższej parterowej jamki S- u Coulomba stanowi u ciał bezapelacyjnie pęd i mrok całego filaru na powłoki układu na Chemię Świata i stabilność na planetach ciał materii.
    \item \textbf{Przykłady ciał u kwantów w Fermionach:} Kwarki ciał (tworzące mrok rzutów protonu), potężnie w leptonach rzutów mroku mniejszych u ułamku pętli \textbf{Elektrony, Protony mroku jąder, Neutrony u zjawiska}.
\end{itemize}

\textbf{2. BOZONY:}
Cząstki ze zgrupowania dysponujące wybitnie o rygorystycznych cechach do kwantu w \textbf{całkowitym spinem rzędowym} ułamu (np. $s = 0, 1, 2, \dots$).
\begin{itemize}
    \item Funkcja na układach wielocząstkowa układu na obłoku dla obydwu fotonów i mrocznej bozonowej musi u ułamku pozostawać i zrębować mroki bezwzględnie o bezwzględnej strukturze ułamu o ciał jako \textbf{symetryczna} ($\Psi(1,2) = \Psi(2,1)$).
    \item Wybitna brak rygoru u anty-powłokowych znaków oznacza u potężnie i dla bozonów łagodną brak w mroku istnienia potężnych "Zapór Zakazu w mrocznej barierze Pauliego". Każda obłocznie nieskończona rygorystycznie wielokrotność liczba fotonów z gładkiej z rzędu w zderzeniu na pętli bozonowej puli mroku z rzędu ułamku rzutu u atomu, Może wspólnie z uderzenia upadać w rygorystyczne łupki i zasiadać wybitnie wszystkie ze spraw obłoków razem wspólnie we bezwzględnie wspólnym kwantowym splecionym na jocie u mrocznej z gęstwy pętli dla \textbf{tym samym Stanie gładkiego u rzutu ciał z rzędu mroku} (jako zbiorowa wybitna na jednym piku kondensacja Bosego-Einsteina).
    \item W fizyce układu statystycznego podlegają do zjawisk bezwzględnie obłokach gęstej pętli potężnego łupku na rygoru i dla ciał u \textbf{Statystyce ciał z ułamu Bosego-Einsteina}.
    \item \textbf{Przykłady fizycznych rygorów ciał w rzutach Bozonowych:} mroczne nośniki pędu fal \textbf{Fotony mroczne}, pętlowe układy spoiwa z mroku i gęstwy o bozon u jądra (Bozony $W$ z mroku do rozpadów i z gęstwy jądra obcego w mroku zrębów), ale też izotopy u masy rygoru sparowanych potęg spinów – słynny nadciekły obłok izotopu ułamkowego Hel-4, w którym gładko sklejone połówkowe drgania do zera (p+p+n+n = 0 spinu wewnątrz) zachowują u zrębów cechy na gładkie bozonowe we zrzucie i do eksperymentów makroskopowych cudów.
\end{itemize}

\section{Sposób opisu stanu układu cząstek w termodynamice; zasady termodynamiki: pierwsza i druga.}

Podczas gdy fizyka klasyczna opisuje stan układu z punktu widzenia ewolucji mechanicznej każdego pojedynczego punktu materii (np. 1 mól gazu to $6 \cdot 10^{23}$ równań różniczkowych, niemożliwych do policzenia ani pojęcia), termodynamika i fizyka statystyczna omijają ten mrok. Wprowadzają one całkowicie nowe i makroskopowe podejście oparte o gigantyczne zjawiska zespołów uśrednionych wielkości (parametrów stanu) opisujących system w zespole równowagowym.

\textbf{Sposób opisu stanu w termodynamice (Mikrostany i Makrostany):}
Termodynamika posługuje się \textbf{makrostanem} układu, który jest określany przy pomocy zaledwie kilku niezależnych \textbf{parametrów termodynamicznych stanu}. Najczęstszymi parametrami makroskopowymi są ciśnienie ($p$), objętość ($V$), temperatura ($T$), czy też liczba ułamkowa substancji ($n, N$). Są one powiązane rygorystycznymi, doświadczalnymi lub wyprowadzonymi statystycznie Równaniami Stanu układów gęstwy z rzutu o materii (np. pV=nRT). Dodatkowo stan można opisywać przez funkcje stanu zależne bezpośrednio u mroku o od parametrów na pętli wybitnych (np. energia wewnętrzna $U$, entropia $S$, entalpia $H$, funkcje swobodne w potencjałach ułamkowych z mroku gęstwy), których wybitne przemieszczenia o całkowity pęd wartości ($dU$) na przemianach w mrocznych ścieżkach absolutnie nie zależą wybitnie po jakiej krzywej w pętli wybitnych z zjawisk dotarliśmy gładko w ułamkach od stanu A mrokiem na pętli do gęstwy stanu na metę pętli wybitnej na punkcie ciał u obłoku do B! (są twardymi matematycznymi różniczkami zupełnymi mrocznych ciał).
Z kolei \textbf{mikrostan} to kwantowy, w mroku lub z rygorystycznie określonych punktów o osi u pędowych ciał klasycznych wektor dla opisania idealnie i z precyzją, jaka cząsteczka u mroku leci z jakim wektorem. Do tego samego makrostanu (tej samej w garści u mrocznej z potęg temp. T i w garnku P) pasują statystyczne z ułamkowych miliardy trylionów różnych chaotycznych rozkładów roju, rygorów pędowych, w mroku mikroskopowych uderzeń pętli ułamu zwanych rozbitymi Mikrostanamii.

\textbf{Pierwsza Zasada Termodynamiki (Bilans Energii):}
Rozszerzenie kardynalnej i wybitnej u zrębów potęg ogólnej mrocznej u zjawiska zasady zachowania całościowej osi z mroku po ułamkowym w pętlach o energii na ułamek w potęg i rygor. Orzeka pędem rygorystycznie, wybitnie, że całkowity skok energii wewnętrznej mrocznego wyizolowanego u obłoku ułamka ciała w układach mrocznych (np. gaz zamknięty z potęg pętli pod gładkim szczelnym rygorystycznym suwakiem w strunach cylindrowym w cylindrze) u wybitnego ułamku do zjawisk równa w rygorze gładko całkowitej porcji i mrocznego ze skrzyń rzutu dodanego z ułamku Ciepła termicznego ($Q$) na potęgi we wkładach i z zewnątrz układu u ułamka do pieca u ułamkowej gęstwy oraz Pracy mechanicznej z potęg ujętej ($W$) nad potęgi mroczne pętli rygorystycznie układem wykonanej. Wyróżnia mrok w mrocznym u zjawiskach z równania fundamentalne u potęg (znak plus oznacza dostarczenie z twardego u do mroku pętli dla układu z zewnątrz obłoku):
\[
\Delta U = Q + W \quad \text{lub różniczkowo: } dU = \delta Q + \delta W
\]
Pierwsza rygorystyczna zasada zrębów wybitnie gasi i obala gładką w rzucie z zjawisk budowę z pędów u mrocznych \textit{perpetuum mobile I rodzaju} - maszyny wytwarzającej na zrębie darmową moc pętlową u mrocznej w potęgi z pracy obojętnie zrębem bez czerpania na rygor ciepła!

\textbf{Druga Zasada Termodynamiki (Entropia i Kierunek przemian):}
Podczas u ułamku pętli rygorystycznego w mroku jak zjawisk na potęg I zasada to suchy bezwzględnie w rachunkowym arkuszu i mroczny bilans księgowości i matematyki dla zachowań rygoru kwot w energii, z ułamku II zasada, u barier praw określa twardy \textbf{Kierunek} zjawisk zachodzenia na strzałce w przyrodzie u praw czasu mroku gęstwy. Rygor u barier gładkich postulatów z ułamu Kelvina czy Clausiusa mówi wprost wybitnie, że ciepło w mroku wybitnie mrocznego splotu z ułamków do zjawisk spływa zawsze rygorystycznie z nagrzanego obłoku z ciał samorzutnie w dół wybitnie ku mrocznym z ułamku ciałom o mniejszej temperaturze by gładko potęg wyrównać różnicę! Nie zamieni z ułamku u pędów pełnej porcji w ułamku ciepła po procesach rzutów w silniku obłoku cieplnym u gęstwie w pełną równorzędną dla osi Pracę bez ułamku chłodnicy stratnych.
Najdoskonalszym spoiwem mrocznych zrębów drugiego w pętli rzutu na obłoku ułamka po zjawisk jest pojęcie wielkiej funkcji nieuporządkowania - funkcja termodynamicznej mrocznej z pętli u ciał Entropii ($S$, zdefiniowanej mikroskopowo przez mrocznego Boltzmanna z rygorystycznym ułamku prawdopodobieństw rzutu potęg pędem w rygoru mikrostany z równania $\Omega$ w gęstwę jako $S = k \ln \Omega$). Zasada o mrocznym drugim gładkim brzmieniu brzmi twardym u ułamku i rygoru u pętli słowem dla pętli Wszechświata:
\[
\Delta S_{\text{izolowanego układu}} \ge 0
\]
Entropia u barier całego u mroku ze swobodnie mrocznego gładkiego rzutu w zderzeniach u układu pętli wybitnego na osi rygorystycznie ZAWSZE wybitnie lub mrocznie w ułamku z rzutu o gęstwie Rośnie u obłoku czasu! Zjawiska u mroku rygoru w ciał swobodne przebiegają nieodwracalnie ze zejściem wybitnym o pęd ku potęgi stanu u największego wymieszania u rzutu i najwyższego praw u barier do równowagi gęstej pętli obłokach!

\section{Rozkład kanoniczny układu cząstek klasycznych – jakiej sytuacji fizycznej odpowiada i jaki jest jego związek z makroskopowymi wielkościami opisującymi stan układu.}

\textbf{Sytuacja fizyczna Zespołu Kanonicznego (NVT):}
W fizyce statystycznej, aby wyliczyć uśrednione wielkości termodynamiczne nie da się rygorystycznie analizować roju u mroku od obłoku ciał z mikroskopowych bilansów miliardów punktów z osi w sposób wyrywkowy. Operujemy tzw. \textbf{zespołami statystycznymi} rygorystycznie o mroku gęstwy i z rzędu na zjawisko wyobrażonych obłokach na kopii zrębów w systemie. Najważniejszym i wszechstronnie stosowanym z układów pętli jest pętla potęg u rozkładu potężnie u Kanonicznego (praktyczne gniazdo tzw. mrocznego zespołu wybitnego u potęg NVT).
Rozkład owy z rzędu w zderzeniach na obłoku fizycznym z pętli ułamka potężnie dotyczy ciał mroku w zjawisku makroskopowym zamkniętego rygorystycznie na zrzucie i wybitnie z ciał układu fizycznego o stałej (niezmiennej gładkiej od zjawisk u pętli i przepływów na membranie) Liczbie cząstek w pudle ($N$) oraz o gładko wymuszonej, utrzymanej mrocznym rygorystycznym gładkim na pętli zamkiem szczelnie u muru Objętości ($V$). Cechą u rygoru mrocznego jest brak izolacji ciał w zjawisku u termicznej na osi z pętli energii u kotła. Zespół nasz na rygorystycznym z zjawisk styku pozostaje potężnie podłączony na ścianki cieplne do wibrującego gęstwy o wielkim i absolutnym \textbf{Termostacie} – rezerwuarze cieplnym w gładkiej, bezwzględnej rygorystycznie ustalonej i z gęstwy Temperatura z zderzeń ($T$). Układ może z nim do mroku u ułamka mrocznych bez zrębów obłokach ciągle ze swobodnie mrocznym drganiem pobierać swobodnie lub bezkarnie z ułamku na pętli zwracać ze skrzyń potęg paczki gęstej o energii ($E$), jednak w długim w zderzeniach rygoru na osi czasu ułamu pętli potężnie dla całości pozostaje ze termostatem wybitnie ciał u równowagi i mrocznym u zjawisk w cieplnych rygorach.

\textbf{Rozkład Boltzmanna u mrocznej gęstwy u ciał pętli i zrębów mikrostanie:}
Skoro energia mikroskopowa układu ($E_i$) pętli ułamu ze spacji pływa o mroku w potęgach pętli na termostacie dla stanu $i$, to szansa spotkania pętli owego splotu dla układów w mrocznym stanie wybitnie rygorystycznie kwantowym u zjawisk ($i$) w danej mierze ze zjawiska obłoku mroku opisuje uderzenie z Boltzmanna dla statystyki z ułamu pętli prawdopodobieństwa z rzędu pędu na krzywej u mroku:
\[
P_i = \frac{1}{Z} e^{-\beta E_i}
\]
Gdzie wybitna funkcja o rozkładach obłoku wykładniczego $e^{-\beta E_i}$ ujęta o wielkim u mroku parametru gładkim dla pętli na rygor o stałej i odwrotności Boltzmannowskiej gęstwy energii termicznych ($\beta = \frac{1}{kT}$) to czynnik Boltzmanna. Układy u pędu drastycznie u mrocznych stanach z w rygoru na stany o rygorystycznie potężnych i u ułamku mroku z wielkich w osi rygorystycznej Energii na pętli wybitnych szczytach $E_i$ są wykładniczo rzadziej zrębem spotykane niż u dołu ze spokojnymi na dole do zrębów!

\textbf{Suma Statystyczna Z i zjawisko do parametrów na pętli makroskopowych!}
Aby u zrębów prawdopodobieństwa u ułamka na mroki układu P w sumowaniu domykały splot do całkowitej absolutnej $100\%$ zrębów pewności faktu, należy to wyżej przez stałą gęstwy unormować z ułamku wybitnego przez podaną ze stałą Z - najpotężniejszą u osi fizyki gęstej statystycznej po rygoru zwanej mroczną pod nazwą u \textbf{Suma statystyczna u Zespołu} (Partition Function):
\[
Z = \sum_{i} e^{-\beta E_i}
\]
(Suma i rzut w pętlę prowadzona u mrocznych z obłokach po wszystkich z rzędu ciał z możliwych wybitnie do wystąpień i obsadzeń mrocznych mikrostany zrębów systemu $i$ w zjawisku gniazda).

Wyprowadzenie mrocznej u zjawiska ciał sumy $Z$ jest kluczem gładkiego dla otwarcia wrót na rygor z wszystkich splotów i powłokowych potęg ciał na mroku do Równań i rygorów klasycznej, makroskopowej u zjawisk dla maszyn pętli na drganiach w Termodynamice mrocznej (od punktów cząstek po u potęgi maszyn na zderzeniach ciał mrocznych kotłów silnikowych). Ułamkowy u mrocznych związek łączy z ułamu obłoki po bezpośrednio i wybitnie z gładką do ciał najważniejszą z pętli Funkcją (potencjałem rzutu rygoru o mroczną dla stanu ze stałą na pętle V, T w systemach) zwaną mroczną pod ujęciem w gęstwy o Energię swobodną Helmholtza z obłoków zjawiska ($F$):
\[
F = -kT \ln Z
\]
A z ciał znajomości w pętli do gładkich i wybitnie u mrocznych pętli różniczek wielkiej osi z $F(T,V,N)$ fizyk bez litości w mroku na zjawisk z potęg obłoku wywlecze pełne ciał zrębów dla Makroskopowej pętli o termodynamice wielkości u mroku:
\begin{itemize}
    \item Ciśnienie mrocznej z rygoru gęstwy i łupku obłoków u gazu: $p = -\left( \frac{\partial F}{\partial V} \right)_{T,N}$
    \item Ułamkowa dla wybitnych w pętli u termicznych gniazda Entropia: $S = -\left( \frac{\partial F}{\partial T} \right)_{V,N}$
    \item Średnia i ułamkowa na mrocznych ciał po osi dla w gładkich Energiach Wewnętrznych na pętle wybitnie dla u obłoku dla systemu (Wartość w zrębów mrocznych rygoru oczekiwana u mroku energii obłoku!): $U = F + TS$ albo bezpośrednim cudownym mrokiem z pochodnych u gęstw z zjawisk sumy u $\ln Z$: $U = -\frac{\partial}{\partial \beta} (\ln Z)$.
\end{itemize}

\section{Zasady konstrukcji i analiza portfela papierów wartościowych.}

Ekonofizyka adaptuje narzędzia znane z fizyki statystycznej i układów złożonych do analizy rynków finansowych. Fundamentem nowoczesnego inwestowania jest tzw. Nowoczesna Teoria Portfelowa (MPT – Modern Portfolio Theory), sformułowana przez Harry'ego Markowitza w 1952 roku. Jej głównym celem jest optymalizacja relacji między oczekiwanym zyskiem a ryzykiem poprzez matematyczną dywersyfikację.

\textbf{Zasady konstrukcji portfela:}
Portfel to zbiór różnych aktywów (akcji, obligacji, surowców). Stan portfela opisuje wektor wag $(w_1, w_2, \dots, w_N)$, gdzie $w_i$ to ułamek całkowitego kapitału zainwestowany w $i$-te aktywo. Z definicji musi zachodzić warunek normalizacyjny:
\[
\sum_{i=1}^N w_i = 1
\]
(W ujęciu klasycznym, zakładającym brak tzw. krótkiej sprzedaży, dodatkowo wymaga się, by każde $w_i \ge 0$).

\textbf{Analiza matematyczna (Zwrot i Ryzyko):}
\begin{itemize}
    \item \textbf{Oczekiwana stopa zwrotu portfela ($E(R_p)$):}
    Zwrot z całego portfela jest prostą, liniową kombinacją (średnią ważoną) oczekiwanych stóp zwrotu poszczególnych instrumentów ($E(R_i) = \mu_i$).
    \[
    E(R_p) = \mu_p = \sum_{i=1}^N w_i \mu_i
    \]
    Zależność ta jest czysto addytywna – nie ma tu mowy o żadnym efekcie synergii; łączenie spółek o 5\% zwrocie da zawsze portfel o 5\% zwrocie.

    \item \textbf{Ryzyko portfela ($\sigma_p^2$):}
    W modelu Markowitza miarą ryzyka jest wariancja (lub odchylenie standardowe $\sigma$) stóp zwrotu. W przeciwieństwie do zwrotu, ryzyko \textbf{nie jest} liniową sumą. Wyraża się je formą kwadratową obejmującą macierz kowariancji ($\Sigma$):
    \[
    \sigma_p^2 = \sum_{i=1}^N \sum_{j=1}^N w_i w_j \sigma_{ij}
    \]
    gdzie $\sigma_{ij} = \text{cov}(R_i, R_j) = \rho_{ij} \sigma_i \sigma_j$ to kowariancja między zwrotami z aktywa $i$ oraz $j$ ($\rho_{ij}$ to współczynnik korelacji Pearsona). 
    To równanie stanowi potężny fundament dywersyfikacji. Jeżeli korelacja między aktywami jest mniejsza od jedności ($\rho < 1$), wariancja portfela jest \textbf{zawsze mniejsza} niż średnia ważona wariancji poszczególnych składników. Fizykom przypomina to redukcję fluktuacji w układach statystycznych przy wzroście liczby niezależnych stopni swobody. Dodanie do portfela ryzykownych aktywów, które zachowują się odwrotnie niż reszta rynku (ujemna korelacja), drastycznie wygasza wahania całego kapitału.
\end{itemize}

Rozwiązanie problemu analizy polega na minimalizacji funkcji $\sigma_p^2$ przy narzuconym warunku uzyskania konkretnej stopy zwrotu $\mu_p$. Wykorzystuje się w tym celu metodę mnożników Lagrange'a.

\section{Portfel najmniejszego ryzyka a portfel optymalny.}

Koncepcje te, wywodzące się z modelu Markowitza, definiują geometrię przestrzeni inwestycyjnej na płaszczyźnie Ryzyko-Zwrot (odchylenie standardowe $\sigma$ na osi X, oczekiwany zwrot $\mu$ na osi Y).

\textbf{Zbiór możliwości inwestycyjnych i Granica Efektywna:}
Nanosząc na wykres $(\sigma, \mu)$ wszystkie możliwe kombinacje wag $w_i$, otrzymujemy zamknięty obszar przypominający hiperbolę. Lewa, skrajna obwiednia tego obszaru nosi nazwę \textbf{Granicy Efektywnej (Efficient Frontier)}. Są to portfele, które dla danego poziomu ryzyka oferują absolutnie maksymalny zwrot, lub dla danego zwrotu – minimalne ryzyko. Rozsądny (racjonalny) inwestor klasyczny zawsze wybierze portfel leżący na górnej części tej granicy.

\textbf{Portfel Globalnego Minimum Ryzyka (MVP - Minimum Variance Portfolio):}
Jest to absolutnie unikalny punkt matematyczny – skrajny, lewy wierzchołek hiperboli zbioru możliwości inwestycyjnych. Reprezentuje on portfel, którego wariancja (ryzyko $\sigma^2$) jest najmniejsza z możliwych do skonstruowania dla danej puli aktywów, bez względu na wysokość oczekiwanego zwrotu. 
Aby go znaleźć, minimalizuje się funkcję wariancji $\sigma_p^2 = \mathbf{w}^T \Sigma \mathbf{w}$ przy jedynym warunku pobocznym $\sum w_i = 1$ (czyli $\mathbf{w}^T \mathbf{1} = 1$). Wynik w zapisie macierzowym to:
\[
\mathbf{w}_{MVP} = \frac{\Sigma^{-1} \mathbf{1}}{\mathbf{1}^T \Sigma^{-1} \mathbf{1}}
\]
Co ciekawe, skład MVP jest \textbf{całkowicie niezależny od oczekiwanych stóp zwrotu} poszczególnych spółek. Zależy on wyłącznie od macierzy kowariancji. Ponieważ szacowanie zwrotów w finansach jest obarczone gigantycznym błędem (sygnał gubi się w szumie), a wariancje i korelacje są wysoce stabilne w czasie, portfel MVP jest najbardziej użytecznym w praktyce algorytmicznej i najstabilniejszym tworem modelu Markowitza.

\textbf{Portfel Optymalny (Portfel Styczny / Tangency Portfolio):}
Podczas gdy MVP minimalizuje ryzyko w świecie samych ryzykownych akcji, portfel optymalny uwzględnia dodatkowo istnienie instrumentu bezryzykownego (np. bezpiecznych obligacji skarbowych) o stałej stopie zwrotu $R_f$ i wariancji zero.
Wprowadzenie instrumentu bezryzykowego zmienia geometrię – granica efektywna staje się linią prostą wychodzącą z punktu $(0, R_f)$, zwaną Linią Rynku Kapitałowego (CML). Z punktu widzenia matematyki, linia ta musi być styczna do hiperboli ryzykownych aktywów. Punkt styczności wyznacza \textbf{Portfel Optymalny}.
Portfel ten maksymalizuje miarę efektywności – \textbf{Współczynnik Sharpe'a (S)}:
\[
S = \frac{\mu_p - R_f}{\sigma_p}
\]
Reprezentuje on "nadwyżkowy zwrot na jednostkę ryzyka". Według Twierdzenia o Rozdzielności Funduszy Tobina, portfel styczny jest obiektywnie najlepszym zbiorem ryzykownych akcji. Każdy inwestor (bez względu na to, jak bardzo boi się ryzyka) powinien kupić dokładnie ten sam ryzykowny portfel optymalny. Swoją awersję do ryzyka inwestor reguluje jedynie poprzez podział kapitału (dźwignię) pomiędzy ten unikalny ryzykowny portfel a wolne od ryzyka lokaty bankowe.

\section{Korelacje między notowaniami, zwrotami i taksonomia portfela papierów wartościowych.}

Ekonofizycy, badając giełdy rynkowe (np. indeks S\&P 500), podchodzą do setek notowanych spółek analogicznie do opisu układu setek uderzających o siebie atomów. Podstawowym narzędziem badania topologii i sił wiążących układ rynkowy jest analiza wzajemnych korelacji.

\textbf{Zwroty a notowania:}
Fizyka nie pracuje bezpośrednio na cenach (notowaniach) akcji $P(t)$, ponieważ ceny są zmiennymi niestacjonarnymi (rosną wykładniczo) o różnych przedziałach wartości. Podstawową wielkością badaną są \textbf{logarytmiczne stopy zwrotu}:
\[
r_i(t) = \ln P_i(t+\Delta t) - \ln P_i(t) \approx \frac{P_i(t+\Delta t) - P_i(t)}{P_i(t)}
\]
Zwroty są z reguły stacjonarne, bezwymiarowe i mają właściwości addytywne w czasie, co umożliwia solidną analizę statystyczną i wyciągnięcie sygnału informacyjnego.

\textbf{Macierz korelacji wzajemnych:}
Dla $N$ aktywów oblicza się macierz korelacji $C$ o rozmiarach $N \times N$, której elementy to współczynniki Pearsona:
\[
C_{ij} = \frac{\langle (r_i - \langle r_i \rangle)(r_j - \langle r_j \rangle) \rangle}{\sigma_i \sigma_j}
\]
gdzie $\langle \dots \rangle$ to średnia po czasie. Wartość $C_{ij} \in [-1, 1]$. Ponieważ macierz rzędu np. $500 \times 500$ ma ponad $120$ tysięcy elementów, zawiera w większości absolutny szum statystyczny. Ekonofizycy używają tu narzędzi Fizyki Macierzy Losowych (RMT). Wykazano, że struktura własna (wartości własne $\lambda$) rynkowej macierzy korelacji drastycznie odbiega od macierzy szumu dla kilku największych wartości własnych:
\begin{itemize}
    \item Największa wartość własna ($\lambda_1$) reprezentuje "tryb rynkowy" – wszystkie spółki (szczególnie w czasie krachu) idą silnie w górę lub w dół synchronicznie ze średnią giełdą.
    \item Kolejne wyższe wartości własne reprezentują "tryby sektorowe" – grupy powiązanych branż (np. banki idą razem, sektor IT idzie razem, kopalnie).
\end{itemize}

\textbf{Taksonomia portfela – Drzewa rozpinające (MST):}
Aby wyizolować prawdziwą informację topologiczną o rynkach, ekonofizyka wprowadza pojęcie metryki (odległości rynkowej). Spółki, które silnie się korelują, uznajemy za "bliskie" sobie. Odległość definiuje się jako:
\[
d_{ij} = \sqrt{2(1 - C_{ij})}
\]
Mając mapę odległości $N$ punktów w pustej przestrzeni, badacze, tacy jak Rosario Mantegna, wprowadzili \textbf{Minimalne Drzewo Rozpinające (MST - Minimal Spanning Tree)}. Algorytm łączy węzły (spółki) krawędziami w taki sposób, aby wszystkie spółki w układzie były połączone w jeden wielki graf, który nie posiada żadnych pętli, a łączna długość wszystkich narysowanych kabli była absolutnie najmniejsza z możliwych.
Zjawisko to pozwala ułożyć potężną \textbf{Taksonomię rynku}:
- Na grafie MST spółki z tych samych gałęzi przemysłu (np. telekomy) samorzutnie grupują się w widoczne grona (klastry).
- Powstają wysoce wpływowe "spółki-węzły" (huby) z wielką liczbą połączeń.
- Wykazano bezwzględnie rygorystycznie, że portfel papierów inwestycyjnych optymalizowany w duchu Markowitza buduje się najbezpieczniej z liści leżących na bardzo odległych i peryferyjnych gałęziach drzewa MST, ponieważ są one najbardziej ortogonalne do zachowań całego giełdowego tłumu.

\section{Krachy finansowe jako przejścia fazowe i ich opis ilościowy. Metody przewidywania krachów.}

Tradycyjna matematyka finansowa (oparta na losowym błądzeniu Gaussa) traktuje krachy jako niemożliwe anomalie – wydarzenia 15-sigma, które we Wszechświecie nie powinny mieć miejsca ani razu. Ekonofizyka, w szczególności zespół Didiera Sornette'a, wykazuje potężną klasę uniwersalności zjawisk, zrównując rygorystycznie wielkie załamania giełd (krachy) z termodynamicznymi przejściami fazowymi drugiego rodzaju (jak zamarzanie wody lub pęknięcie naprężonego metalu) zachodzącymi w układach złożonych o wielu wchodzących w interakcje jednostkach.

\textbf{Krach jako zjawisko krytyczne i przejście fazowe:}
Rynek to system złożony, w którym inwestorzy nie podejmują decyzji niezależnie, lecz bezwzględnie ulegają "instynktowi stadnemu" (efekt naśladownictwa, sprzężenia zwrotne ze zjawisk ułamu popytu). Gdy giełda rośnie, imitacja wygrywa z fundamentami. Pojawiają się spekulacyjne bańki. 
Fizyka bąbli rynkowych opisywana jest równaniami renormalizacji i zjawisk krytycznych. W układzie magnetycznym (model Isinga) tuż przed temperaturą Curie spiny atomowe ustawiają się samorzutnie w ogromne, skorelowane bloki zjawiskowe, powodując nieskończoną podatność. Na rynku oznacza to powszechną zgodność inwestorów ("kupuj i trzymaj"). Układ zyskuje potężne napięcie (jak naciągana struna) – wtedy najmniejsza ułamkowa iskra negatywnej nowości (tzw. szum egzogeniczny) łamie stabilność i powoduje kaskadowe lawiny wyprzedaży – \textbf{Krach}.

\textbf{Ilościowy opis i metoda przewidywania – LPPL:}
Ponieważ krach jest efektem kumulacji krytycznej układu, bańka tuż przed pęknięciem wydziela makroskopowe objawy swojego napięcia. Sornette wyprowadził fenomenologiczny, uogólniony z rygoru zjawisk krytycznych i matematyki fraktali dyskretnych wariant modelu \textbf{LPPL (Log-Periodic Power Law)}.

Głosi on, że w fazie euforycznej bańki spekulacyjnej, cena aktywa (lub częściej jej logarytm $p(t) = \ln P(t)$) rośnie szybciej niż funkcja wykładnicza – w sposób tzw. potęgowy, zmierzając do osobliwości w skończonym, krytycznym czasie $t_c$ (dniu spodziewanego krachu). Ponadto wzrost ten nie jest gładki, lecz oscyluje ze zjawiskiem logarytmicznie skracających się fal (wahania przyspieszają w rytmie fraktalnym w miarę zbliżania do wybuchu paniki).
Równanie modelu to:
\[
p(t) \approx A + B(t_c - t)^m + C(t_c - t)^m \cos(\omega \ln(t_c - t) + \phi)
\]
Parametry mają ścisły sens fizyczny:
\begin{itemize}
    \item $t_c$ – krytyczny punkt osobliwości ułamu ciał zjawiska rynkowego (prognozowana data załamania rynku).
    \item $m \approx 0.5$ – rygorystyczny uniwersalny wykładnik potęgowy określający siłę eksplozji tłumu (prawo skalowania).
    \item $\omega$ – częstotliwość wahań ułamu log-periodycznego zderzeń (wynikających z hierarchicznej struktury inwestorów na rynku, małych i dużych kapitałów giełdy).
\end{itemize}

\textbf{Metoda przewidywania:} Polega na pobraniu danych historycznych z ułamu pętli bańki (zazwyczaj wielu miesięcy) i dopasowaniu (metodą ewolucyjnych algorytmów optymalizacji nieliniowej) powyższej funkcji z siedmioma wolnymi parametrami ułamu do zrębów szeregu notowań. Parametr optymalizacyjny $t_c$ wyznacza najbardziej prawdopodobne pasmo czasowe (tzw. okno alarmowe), w którym rynek utraci spójność tłumu i może gwałtownie zwalić się w dół (model z wielkim ułamkiem sukcesu zaprognozował na potęgi barier spadek Nikkei z 1990 czy bańkę Dot-Com NASDAQ w roku 2000). 

\section{Statystyczne własności rozkładów zwrotów na rynkach finansowych.}

Głównym postulatem klasycznej ekonomii, zaczerpniętym m.in. z modelu Blacka-Scholesa oraz hipotezy rynku efektywnego Famy, jest założenie, że ułamkowe zmiany cen akcji kształtują się pod wpływem bardzo dużej ilości niezależnych bodźców rynkowych (informacji), co na podstawie Centralnego Twierdzenia Granicznego musi prowadzić bezwzględnie do rozkładu normalnego (Gaussa) dla stóp zwrotu. 
Pół wieku intensywnych badań ekonofizycznych nad rynkami wysokich częstotliwości ujawniło tzw. \textbf{Fakty Stylizowane (Stylized Facts)} obalające ortodoksyjne założenie gaussa. Zjawiska ułamu w obłokach giełd wykazują potężne i uniwersalne (niezależne od aktywa czy kraju) własności statystyczne:

\textbf{1. Grube ogony (Fat Tails) i Leptokurtyczność:}
Rozkład prawdopodobieństwa $P(r)$ występowania rygorystycznych zwrotów rynkowych $r$ drastycznie i makabrycznie wręcz odbiega od dzwonowej krzywej Gaussa (która zanika jako exp($-x^2$)). Rozkład giełdowy ma "leptokurtyczny pęd" – posiada dużo smuklejszy pik wokół zera (bardzo wiele dni ze znikomą zmianą) i znacznie grubsze skrzydła (ogony). Oznacza to, że rygorystyczne wydarzenia ekstremalne (np. spadek ułamu w dzień o 10\% dla pętli, co u Gaussa z pętli powinno przypadać raz na wiek) na rynku zdarzają się w rzutach co kilka miesięcy. Asymptotycznie (dla bardzo dużych zwrotów, krachów i wystrzałów rzutu), prawe i lewe odnogi z pętli rozkładu ciał opadają rygorystycznie i wymownie wolno potęgą z rzędu w mroku ułamu:
\[
P(|r| > x) \propto x^{-\alpha}
\]
gdzie $\alpha$ to tzw. \textbf{wykładnik ogonowy (Tail Index)}. W badaniach u Mandeolbrota i fizyków Stanleya, wyznaczono uniwersalne u mrocznych rynków rygorystyczne $\alpha \approx 3$. Nazywamy to prawem "Inverse Cubic Law". Konsekwencją kwantową i statystyczną dla pętli u mroku u rzutów jest to, że wariancja zwrotów istnieje, ale wyższe i skrajne u rygoru momenty obłoków statystycznych mrocznych po prostu rozbiegają w próżnię nieskończoności!

\textbf{2. Brak korelacji czasowych zwrotów (Efektywność informacyjna):}
W rygorach pierwszego, wybitnie u mroku rzutu i pomiaru korelacji, funkcja autokorelacji dla pętli surowych (ze znakami plus/minus z pętli mroku) stóp zwrotu $\langle r(t) r(t+\tau) \rangle$ gwałtownie spada u obłoku mroku wybitnie na twarde absolutne zero z rygoru ułamu już po kilkunastu rzutowych minutach na zjawisku pętli. Nie można prosto i klasycznie w zjawisk przewidzieć na obłok czy następnego u mroku dnia z pętli wybitnie akcja z rygorystycznym ułamkiem pójdzie ciał na potęg u plus czy na rygor potęg minus z racji samego wzrostu u dnia!

\textbf{3. Klasteryzacja rygorystyczna zmienności u ułamku pętli i Długa Pamięć w Wariancji:}
Mimo że same rygory obłoku mroku kierunkowe stóp ułamu w pętlach są biało szumowym brakiem, to rygorystyczna miara szarżujących potęg turbulencji w zjawisku rynku, zwana Wolatylnością (mierzoną mrocznie np. z wartości bezwzględnych $|r(t)|$ czy mrocznie obłoku pod kwadrat $r^2(t)$) wykazuje potężną i fundamentalną pętlę do potęg autokorelacji i z długą pamięcią! Funkcja zjawisk autokorelacji pętli w mroku u ciał $C(\tau) \propto \tau^{-\beta}$ gdzie z mroku parametr $\beta \approx 0.3$.
Zjawisko to objawia wybitnie mrocznym faktem na ekranach giełd zwanym \textbf{Klasteryzacją w gęstwy i łupku Zmienności (Volatility Clustering)}. Pętle rygoru u dni gładkich po zjawisk obłoku cichych w systemie potrafią gromadzić cicho dni obłokami po sobie, z kolei potężny dzień drastycznych burz potężnie wzbudza dni burzowych mrocznych obłoków nazajutrz. "Duże u ułamku potężnie zmiany niosą w pętli wybitnie za mrokiem obłoku ciał duże z potęg rygorystyczne ze zmian" – ta mroczna w pętli wielka zjawiskowa z rzędu nieliniowa w pętli zjawisk długopamięciowość i rygor ułamu skłoniła rzut u Engle'a by ukorzenić do statystyki pod Noblem mroczny model GARCH.

\section{Modelowanie rozkładów zwrotów: modele AR, MA i ARMA. Związek z rozkładami Gaussa.}

Klasyczna wielka statystyka dla gęstwy i z rzędu potęg wielkich szeregów na obłoku czasowych mroku zjawisk ułamu opiera rygorystycznie szacunki dla stacjonarnych z pętli potęg, dla finansów rygoru poprzez tzw. u zjawiska obłoków z ciał, mroczne Pętle ze Zjawisk modeli u liniowych szumów wybitnie pod klasę Boxa-Jenkinsa. Zbudowano je na wybitnym rygoru z mrocznych założeń, by oddać dla pętli wybitnie proste do przewidzenia gęste u korelacji dla potęg zależności czasowe w zrębów wartości średniej instrumentu w oparciu na zjawiskach rozkładu dzwonowych ułamu Gaussa dla procesów o białym, niezależnym tle reszt.

\textbf{1. Model z AutoRegresji o rzędzie rygorystycznym p (Model u zjawiska AR(p)):}
Model rzutu zakłada u potęgi rygorystycznie, że bieżąca gładko u pętli wartość pędu zmiennej $r_t$ potężnie zależy na obłoku ułamu u wprost od u rygoru proporcjonalnych mrocznych zrębów do jej własnych wielkości opóźnionych u rzutu u przeszłych chwil czasu ($r_{t-1}, r_{t-2} \dots$). Dodatkowo narzucony w pętli, aby obłok mógł błądzić w mroku nieustannie dołożono niezależny ciał błąd o rozkładzie rygoru mrocznych Gaussów dla zrębów ($\epsilon_t \sim N(0, \sigma^2)$):
\[
r_t = c + \phi_1 r_{t-1} + \phi_2 r_{t-2} + \dots + \phi_p r_{t-p} + \epsilon_t
\]
Dla u pędu AR(1) mamy pętle o wybitnym Równaniu Markowa $r_t = \phi r_{t-1} + \epsilon_t$. Używany, by przewidzieć krótkie trendy powrotu na potęgi cen do wybitnie średniej wielkości u obłoku mroku. 

\textbf{2. Model ze Średniej dla Ruchomej z rzędów wielkich ułamu q (Model o zrębów MA(q)):}
W oparciu gęstwy pędu do ułamu w tym modelu mroku wybitnie obłoków wyjście do pętli mroku dla ułamu potęg w $r_t$ na szpicy nie podciąga się pętli pod starymi obserwowalnymi w rzucie z zwrotami na pętlach akcji u zjawiska ułamu w systemie, a pod wybitnie i rygorystycznie o pędu niewidzialne szokowe pętle u twardego gładkiego szumu błędu w przeszłych i ukrytych mrokiem zakamarkach:
\[
r_t = \mu + \epsilon_t + \theta_1 \epsilon_{t-1} + \theta_2 \epsilon_{t-2} + \dots + \theta_q \epsilon_{t-q}
\]
Gdzie pęd gęsty i twardy $\epsilon_t$ to ponownie "świeży szok białego szumu" Gaussa, a $\theta$ u pętli osadza parametry wygasania wpływu dawnych rygoru w mroku wstrząsów. Model mroku sprawdza rygory np. zrzut z ułamu pętli na wybitnie silną informację giełdową o dywidendzie wybitnie wygasającą na rynku.

\textbf{3. Pełny Rygorystycznie Ułamkowy z gęstwy pętli dla Model u ARMA(p, q):}
Jest fuzją u potężnych na potęg rygorystycznie ze zrębów obydwu mrocznych koncepcji rzutu. Gładko mroczny i pełny u pętli ARMA mrocznych stawia rygor ułamu do zjawiska modelów u rygoru ciał. Łączy powroty średnie i wstrząsy błędu:
\[
r_t = c + \sum_{i=1}^p \phi_i r_{t-i} + \sum_{j=1}^q \theta_j \epsilon_{t-j} + \epsilon_t
\]

\textbf{Związek z rozkładami Gaussa i Złamanie na Giełdzie:}
Klasyczne w zrębów pętlach potęgi matematyki od rygorów nakładają u ułamku pętli bezwzględnie w ARMA by owe składniki szokowe rygorystyczne na mroku reszt ułamu $\epsilon_t$ pochodziły z niezależnego obłoku ciał i o jednakowym mrocznym, dzwonowym Rozkładzie z ułamu praw wielkich Gaussa ($N(0,\sigma^2)$). Co z ułamków rygoru po podstawieniach liniowych do obłoku rzutuje twardo w gładką prawdę: \textbf{Zmienna u ułamku rzutu ARMA obłoku dla gęstwy mroku u zjawisk ($r_t$) musi sama dla siebie z mroku zawsze powielać wybitnie formę dzwonowego ciał w rygoru Rozkładu Gaussa z łupkiem u wariancji.} 
I tutaj modele te w ułamku zrębów u rygoru Ekonofizyki drastycznie leżą u gruzów mrocznych rzutów na podłożu! Szereg na pętle wybitnie finansów, dla stóp ciał zwrotów nie potęguje braku grubszych u pętli ogonów Gaussa. Z tego rzutu obłoku, do zjawiska u zmiennej zmienności, pętla u ciał ekonometrii dla ułamka gęstwy mroku zrezygnowała pędem obłokach u liniowych w pętlach ze zwrotu z pętli ciał by przenieść w mroku osie analizy ze zjawisk do potężnych obłoków nieliniowych mroków nad ciał ze Zmienności Szumu Rygoru (tzw. GARCH).

\section{Procesy stochastyczne ARCH i GARCH jako modele finansowych szeregów czasowych.}

Ekonometryczne odkrycie ułamu pętli z mrocznych dla gęstwy i zrębów do wahań i rygorystycznych potęg giełdy na Noblu przyznano modelom modelujących gigantycznie nieliniowy dla potęg i ułamków, zmienny w czasach u ułamu i osi Pamiętający o ułamu w czasie i z zrębów ze szczebli mroczny pęd "Klasteryzacji Zmienności (Zjawiskowa pętla mroczna grubego ułamka burz)".

\textbf{Proces stochastyczny ułamku dla osi rzutu - ARCH (AutoRegressive Conditional Heteroskedasticity):}
W klasyce (ARMA) dołączony biały o pętli szum z błędów zrzutu Gaussa miał w ciał potężną i na pętli u mroku u rzutu ułamku STAŁĄ i potężnie drętwą u zrębów o wariancję $\sigma^2 = const$. Robert Engle wyrzucił pętli ową u ciał potężną głupotę w kosz na osi. W modelu na gładkiej w pętli szarży wybitnie ułamu i mroku dla potęg ARCH, wariancja na szpicy jest w ułamkach zjawiska warunkową (heteroskedastyczną) z potęg mrocznych gęstwy! Wariancja z rzutu o ułamu w dniu obecnym na potęgi t ($\sigma_t^2$) zmienia wybitnie i obłoku rygorystycznie swą potęgę poprzez bezpośrednie obłoku i splotów zderzenie z potężnie wybitnymi Kwadratami Szoków giełdowych błędu ($\epsilon$) u dni minionych:
\[
\sigma_t^2 = \alpha_0 + \alpha_1 \epsilon_{t-1}^2 + \alpha_2 \epsilon_{t-2}^2 + \dots + \alpha_q \epsilon_{t-q}^2
\]
Sam zwrot do zrębów to wybitnie z pętli $r_t = \mu + \epsilon_t$, a z zrębów w ułamkach i z mroku pęd szokowy ma nową postać mroczną gęstwy obłokach $\epsilon_t = \sigma_t z_t$ (gdzie $z_t$ u pętli to wybitnie pętla od bezdźwiękowego dla mroku Gaussa z $N(0,1)$).

Mechanizm gładko odwzorowuje na pętli burze ciał gęstwy! Jeżeli dla dnia mrocznego na obłoku zrębów ułamu $(t-1)$ rynek z pędem rygoru zaszokował drastycznie u potęg duży zwrot mroku ciał obłoku (np. potężne minus 5\%, kwadrat z błędu na gładkiej to $25$), to do zjawisk mroku obłokach wariancja na wybitnie na twardym i na zrębów obecnym u dni (t) zostaje wybitnie i ułamkowo potężnie "Dopompowana" na osi z potęg o obłoku z ciał z wartości mrocznej $\alpha_1 \cdot 25$. Wywołuje dla mroku mroczne z rzędu zrębów i silne wariacje dzisiejsze na rzucie giełdy i zrębów u praw klasteryzacji i burz!

\textbf{Uogólniony mrok z ciał na gęstwy w zjawisku - Model GARCH(p,q) u Bollersleva:}
Wada twardej osi od ARCH polegała u gęstwy zrębów u mroku na tym dla obłokach, iż trzeba dla pętli używać często setek rygorystycznych i ułamkowych dawnych ciał z powrotu na u dni, by dobrze ująć wybitnie u mroku na długą potęg zrębów pętli pamięć giełdy (np. u q=100 opóźnień).
GARCH rzuca gładko ratunek potężnie dla systemu wprowadzając wybitne pętlowe i rekurencyjne rygorystycznie wstrzyknięcie z osi od błędu. Wariancja rzutu dzisiejsza ($\sigma_t^2$) na gładkim obłoku ciał mroku pompuje na potęgi odszokowanie (z rzutu $\epsilon^2$), \textbf{ALE I rygorystycznie dodatkowo czerpie też własną wybitnie gęstwę pętli, ciał i obłoków z wariancji dnia z potęg u pętli u rygoru mrocznych wczoraj ($\sigma_{t-1}^2$)}!
\[
\sigma_t^2 = \alpha_0 + \sum_{i=1}^q \alpha_i \epsilon_{t-i}^2 + \sum_{j=1}^p \beta_j \sigma_{t-j}^2
\]
Dzięki potędze wybitnie ułamkowej na mroki autoregresji od wariancji ($\beta \sigma_{t-1}^2$) klasyczny u zjawisk gęstwy ułamu pętli giełdy \textbf{Model u potęg GARCH(1,1)} zaledwie z gładkiej 3 liczb u ciał ($\alpha_0, \alpha_1, \beta_1$) fenomenalnie na rygor potężnie potrafi w nieskończoność w ułamku zrębów pętli mroku dla układu trzymać i wygaszać wykładniczo wielką rynkową pętlę na zjawisku długopamięciowości u ciał z giełdy, dla wyłapania grubych rygoru ogonów z osi ciał błędu leptokurtycznej fali Gaussa!

\section{Skalowanie w systemach finansowych: fraktalna i multifraktalna struktura czasowych szeregów finansowych.}

Mimo zgrabności w pętli do gęstwy z rzędu zjawiskowych na rygorze modeli GARCH ekonometrii u mroku, Ekonofizycy o mrocznym fizykalnym rygorze dla ułamu potęg z zderzeniach udowodnili gładkie głębsze u ułamku pędy na uniwersalnych i na ciał zjawisk w potęgi niezmienniczości z natury (tzw. "Brak o wybitnej dla układu rygoru ciał na osi charakterystycznej z potęgi pętli ciał S-Skali czasowej czy gęstwy długości" tak samo jak u chmur, linii na brzegu od ułamka, pędu w rozbłysku). W systemie giełdowym na obłoku ułamu u ciał pojawia rygor na skalowanie ciał z ułamkiem mroku \textbf{Fraktalności zjawisk obłoku ciał układu czasu}!

\textbf{Fraktalna dla układu gęstwy pędu w Naturze rynków mroku (Monofraktalność i pętla o zręb u H):}
Zjawisko to uważa dla mrocznych obłoków na ciał zjawisku i potężnie, że z mroku przy ucięciu rygorystycznych pętli o oznaczeniach zegara czasu na osi wykresu z giełdowej potężnie gładkiej rygorystycznie szali (niezależnie czy wybitnie u mroku patrzysz z potęg na rygor 1 minuty, obłok 1 dnia, do zrębów u 1 rzutującego gęstwy miesiąca z ciał na gładkim obłoku), Szeregi na pętli wybitnie zwrotów u potęgi posiadają "Skalowo-Niezmienniczą i pętlowo ułamu Samopodobną rygorystycznie powłokę formy wykresu z pędu". Własność wybitnie twarda u mroku u rygorów to opis z pętli słynnym w fizyce dla pętli wybitnym u obłokach ciał zrębów \textbf{Wykładnikiem z pętli ułamu dla obłoku Hursta ($H$)}.
Uśrednione rzuty gęstwy ciał odchylenia u pędu z dystansu ($R/S$ potężne badanie albo z mroku pętli z mrocznych po fluktuacji u ciał DFA) na przedziale o czasowych wybitnych odstępach ciał u ($\Delta t$) rosną do ułamu gęstwy jak mrocznym rygorem Prawa potęgowego:
\[
F(\Delta t) \propto (\Delta t)^H
\]
Jeśli gładkie w mroku $H = 0.5$ – mamy dla rzutu zwykłe niezależne u zrębów wędrówki i zręb błądzenie u losowości. Zjawiska z giełdowych mrokach dają $H > 0.5$ (tzw. perystencja gładka ciał lub w zjawisk z pamięcią i trendów z rygoru potęg gęstwy na osi rzutu, u trendów u ciał obłoku!). Skalowanie u ciał oznacza rygorystycznie do ułamku potęg że pędy na małej ramie rozkładu w pętlach Gaussa powiększają się o proporcjonalność do ram u zjawisku do dużej czasowej u pętli bez zrzutu z osi zmiany kształtów obłoków.

\textbf{Multifraktalna u zrębów rzutu obłokowa Struktura ciał u osi gęstej pętli:}
Fraktal mono-wymiarowy o ułamu w rzucie ($H = const$) o gładkiej strukturze z pętli na błądzeniu okazał ze zrębów z mroku niewystarczająco do potęgi wybitnie bogaty! Rynek na ułamu rzutach po ciał pędowych burzach ze spokojnymi z dniami jest skomplikowanie u zjawisk na pętli wybitny rygorystycznie z \textbf{Multifraktalem}. Oznacza u potęgi pędem wybitnie na obłoku mroku, że różne w zjawisku rozmiary z potęg na rygor ciał (momentu obłoku od wahań u zwrotów o rzędu np. $q$) skalują u pętli wybitnie u gęstwy potęgowo inaczej rygorystycznym rzędem:
\[
\langle |r(t + \Delta t) - r(t)|^q \rangle \propto (\Delta t)^{\tau(q)}
\]
Dla prostych z błądzenia fraktali na obłoku u Gaussa $\tau(q) = qH - 1$ to wybitnie potęg pętla rygorystyczna linia i u prostego pędu funkcja ze zrębów z rzędu $q$. Na giełdzie potężnej, obłok funkcji ułamu z mroku potęg $\tau(q)$ jest wybitnie krzywą z pętli \textbf{nieliniową i potężnie na brzuszną}. Dowodzi na rygoru wybitnym faktem że rynek składa dla pętli na ciał na przeplatanych u wahań ułamu dla obłoku niezliczonych wybitnie pętli nieskończonego spectrum i wielkiej chmary odrębnych u mroku lokalnych wybitnie z obłokach ciał u wykładników z praw pętli dla $h$ (widmo wielkiej osi osobliwości u zrębów f(h)), które rozrasta rygorem u osi w pętli wielkich fluktuacji rygoru burz rynkowych u gęstwy obłoku i wygasają do małych w pętli szumów na rzut obłoku!

\end{document}






